\documentclass[11pt]{article}
\usepackage[T1]{fontenc}
\usepackage{lmodern}
\usepackage[margin=1in]{geometry}
\usepackage{amsmath,amssymb,amsthm,mathtools,booktabs,enumitem,array,xcolor}
\usepackage[hidelinks]{hyperref}
\newtheorem{theorem}{Theorem}[section]
\newtheorem{proposition}[theorem]{Proposition}
\newtheorem{lemma}[theorem]{Lemma}
\newtheorem{corollary}[theorem]{Corollary}
\newtheorem*{namedthm*}{\thisthmname}
\newcommand{\thisthmname}{}
\newenvironment{named}[1]{\renewcommand{\thisthmname}{#1}\begin{namedthm*}}{\end{namedthm*}}
\theoremstyle{definition}
\newtheorem{definition}[theorem]{Definition}
\theoremstyle{remark}
\newtheorem{remark}[theorem]{Remark}
\newcommand{\fn}[1]{\texttt{\detokenize{#1}}}
\newcommand{\tC}{\textnormal{\textsc{[C]}}}
\newcommand{\tA}{\textnormal{\textsc{[A]}}}
\newcommand{\tD}{\textnormal{\textsc{[D]}}}
\newcommand{\tF}{\textnormal{\textsc{[F]}}}
\newcommand{\OPEN}[1]{\textcolor{red!75!black}{\textbf{\textsf{[open: #1]}}}}
\newcommand{\AUTH}[1]{\textcolor{blue!55!black}{\textsf{\small[authors: #1]}}}
\newcommand{\thin}[1]{\colorbox{orange!30}{#1}}
\newcommand{\ip}[2]{\langle #1,#2\rangle}
\newcommand{\abs}[1]{\lvert #1\rvert}
\newcommand{\norm}[1]{\lVert #1\rVert}
\newcommand{\kt}{\tilde\kappa}
\newcommand{\Gt}{\tilde\Gamma}
\newcommand{\lt}{\tilde\ell}
\newcommand{\Yb}{\bar Y}
\newcommand{\eps}{\varepsilon}
\newcommand{\cL}{\mathcal L}
\newcommand{\cN}{\mathcal N}
\newcommand{\cQ}{\mathcal Q}
\newcommand{\cR}{\mathcal R}
\newcommand{\cB}{\mathfrak B}
\newcommand{\Lam}{\Lambda}
\newcolumntype{L}[1]{>{\raggedright\arraybackslash}p{#1}}
\setlist{itemsep=2pt,topsep=3pt}
\setlength{\emergencystretch}{3em}
\renewcommand{\arraystretch}{1.12}
\title{Step 4 of the hair-clip proof: from the embedded clip $\Gamma_0$ to the step-3 trapping set near $S_1$\\[2pt]
\large The initial curve, the validated Fourier phase P1, the shadowing phase P2 and the hand-off}
\author{Micah Warren \and Jacob Ogden}
\date{Version 2, November 12, 2026}
\begin{document}
\sloppy
\maketitle
\begin{center}\fbox{\parbox{0.95\linewidth}{\small\textbf{Status} (version 2, Nov~12, 2026; it supersedes draft v1 of Nov~11, kept as \fn{step4_proof_v1.tex}). Every step-4 computation is done in \fn{arb} with outward rounding and no sampling. The final tail is pinned by one manifest, \fn{step4/p2/out/FINAL_manifest.json}, which records every stage with its command, environment and the SHA-256 of its inputs and outputs; a second pass reproduces it byte for byte. The gaps A1--A3 and the findings B1--B9 of the independent review (\fn{step4/STEP4_REVIEW.md}) are closed; Table~\ref{tab:review} gives each resolution. The numbers are those of \fn{step4/STEP4_RESULT.md} (Nov~12), rounded outward.

\textbf{Hand-off.} Theorem~4 delivers the hand-off conditions (H1)--(H6) of step~3 v3.1 \emph{as amended by v3.2}, for every good time in base slab~180 and every $S_1$ in the certified set; Lemma~O H1 for the far family is now closed (cleanup O1). The amendment v3.2 has been approved by the authors and is built in \fn{writeup/step3_proof_v32/}, \fn{writeup/lemmaO_v32/} and \fn{writeup/step3_amend_v32/}. It proves Lemma~O H6 in step~3. \emph{Step-3 dependency.} Lemma~R is proved, with arb-certified constants (\fn{writeup/lemmaO_v32/lemmaR.pdf}); it closes the outer ramp. An independent referee has confirmed a gap in step~3 itself, present already in v3.1 (\fn{writeup/step3_gap_v32/gap_v32.pdf}, \fn{REFEREE_check.md}): the far-field window-energy inequalities of Lemma~S(B) drop the window-edge (total-derivative) terms, and step~3's trapping argument uses the resulting far-field bound $\abs f\le\frac12$ (W4). This is a missing estimate with large numerical room (derived $\abs f\le0.093$ against $\frac12$), not a failed inequality. The main theorem is therefore stated \emph{conditional on Hypothesis~(F)} (\fn{writeup/main/main.pdf}). Step~4's own results do not depend on (F); its v3.2 H6 closure margins use the strip input, which is part of (F). A fix, a uniformly local Lemma~S(B), is being evaluated (go/no-go tests). Under v3.1 alone, Lemma~O H6 is not met (\S\ref{ss:H6v31}).

\textbf{Open} (\S\ref{s:open}): Hypothesis~(F) of step~3 (the Lemma~S(B) gap; the main theorem is conditional on it); the true curve's $k_{\max}$ at $t_1$ (only the base value is certified); kit~B: pending. \emph{Closed in this version:} Lemma~O H1 for the far family, Lemma~R, the script items O6a--c.

\textbf{Thinnest margins} (Table~\ref{tab:margins}): $|\alpha|$ $\times1.27$; cone slope datum $\times1.40$; cone $\sup|h_0''|$ datum $\times1.59$; initial error $Y(\tau_a)$ against $\Yb$ $\times1.6$; far $|e|$ box $\times1.69$.}}\end{center}

\begin{abstract}
\noindent Step 4 of the computer-assisted proof that an embedded symmetric ``hair-clip'' curve develops, under curve diffusion flow (CDF), a finite-time singularity modelled on the non-embedded one-loop shrinker $S_1$ is the \emph{certified passage} from an explicit initial curve $\Gamma_0$ into the trapping set of step~3. (i)~$\Gamma_0$ is an exact sine series for the tangent angle, closed by a Krawczyk test and certified embedded (\S\ref{s:G0}). (ii)~P1 validates the flow in Fourier variables up to $k_{\max}\approx29.35$, through 96 slabs with Lohner blocks and Duhamel tail bounds; a Galerkin-limit lemma makes this an existence and uniqueness theorem (\S\ref{s:P1}). (iii)~After a certified switch (\S\ref{s:switch}), P2 shadows a polished piecewise-polynomial base solution in rescaled variables up to $k\approx2.37\cdot10^4$. A comparison system for modulated tip, intermediate and far-field energies controls the error, with every nonlinear term inside the certified closure (\S\ref{s:base}, \S\ref{s:P2}). (iv)~The hand-off computes the step-3 coordinates uniformly on a whole base slab and checks the hypotheses of step~3 as amended by v3.2, including the Lemma~O data per cylinder (\S\ref{s:T9}--\ref{s:H6}). The curve self-intersects on the way, as it must, since $S_1$ is not embedded; step~3 does not need embeddedness. \S\ref{s:main} combines steps 1--4, \S\ref{s:open} lists what is open, and Appendix~\ref{app:how} says what is computed and how.

\smallskip\noindent\textbf{Status tags.} \tC\ certified by interval/ball arithmetic, or a proof whose only non-closed-form step is so certified; \tA\ analytic argument, standard or from the literature; \tD\ derived on paper in a design note, not machine-checked; \tF\ float only; \OPEN{\dots}: open.
\end{abstract}

\tableofcontents

\section{Summary and Theorem 4}\label{s:sum}

\subsection{Conventions}\label{ss:conv}
CDF moves a closed plane curve with normal velocity $-\kappa_{ss}$. Everything is symmetric about the $x$-axis. In Fourier variables (P1) the tangent angle is $\theta(\alpha)=-\pi/2+\alpha+\sum_{n\ge1}u_n\sin n\alpha$, $\alpha\in[0,2\pi]$, where $\alpha$ is proportional to arclength. The sine class is invariant under the flow, and invariance of the class is the symmetry. Near the singularity (P2, step~3) we write $\gamma=c\,e_1+\lambda\,\Gamma$, $d\tau=\lambda^{-4}dt$, $a:=\tfrac14+\lambda_\tau/\lambda$, $b:=-c_\tau/\lambda$. Then
\[
\ip{\Gamma_\tau}{\nu}=-\kappa_{ss}+(\tfrac14-a)\ip\Gamma\nu+b\ip{e_1}\nu,
\]
which is the frame of \fn{step3_proof.tex} \S1.1. We write $x$ for arclength from the tip and $y:=x-X_0$, $X_0:=2.2490776515$. The step-3 objects are taken from step~3 without change: $S_1$, $\hat\psi_3$, $\hat h_2$ (\fn{agent_int/v3/c4fun.py}), the cutoff $\chi=\chi_{32}$, the weight $\rho$, the constraint functions $f_1,f_2,f_3$, the functional $\Lam=E+\mu E_I$ ($\mu=0.08$), the shells, the windows and Lemma~O. ``Step~3'' means \fn{step3_proof.tex} v3.1 as amended by v3.2 (\S\ref{s:H6}) unless v3.1 is named. The P2 base is stored in slabs of $\Delta\tau=0.15$: slab $m$ covers $\tau\in[0.15m,0.15(m+1)]$, parametrised by $t_c\in[-1,1]$. ``Time since $\Gamma_0$'' is the CDF time; the clip time of the float runs is larger by $0.001$, because $\Gamma_0$ was taken at $t=0.001$ of the original clip run. A number written as a bound is a certified bound, rounded outward.

\subsection{The chain}\label{ss:chain}
\begin{center}\small
\begin{tabular}{@{}L{0.16\linewidth}L{0.60\linewidth}L{0.17\linewidth}@{}}\toprule
link & result (file) & status, \S\\\midrule
$\Gamma_0$ (T2) & 3762 sine modes; $b_2$ by Krawczyk (half-width $8.14\cdot10^{-24}$); embedded (\fn{t2/gamma0_final.json}) & \tC, \S\ref{s:G0}\\
P1 (T3), slabs 0--96 & $t_{96}=4.597850539485830831\pm7.8\cdot10^{-19}$; $k_{\max}\in[29.352874545325,\,29.352874545328]$ (\fn{t3/p1/status.json}) & \tC, Lemma~\ref{lem:gal}, \S\ref{s:P1}\\
base v2 (T1) & 181 slabs, $\tau\in[0,27.15]$, $k$ from $26.25$ to $2.37\cdot10^4$; slab 0 is P1 slab 95 exactly; (SE)-corrected, $|a_2|\le1.6\cdot10^{-14}$ & exact data, \S\ref{s:base}\\
residuals (T6) & T6's engine on the (SE)-corrected base itself, slabs 2--180 (\fn{p2/t6c}; review B1) & \tC, \S\ref{ss:res}\\
switch & P1 slab 96 to base slab 2 at $t_{c,a}=982986695892350/2^{50}$; $Y(\tau_a)\le6.21\cdot10^{-13}$ (\fn{switch_s2.json}) & \tC, $\times1.6$, \S\ref{s:switch}\\
T7 & 26 coercivity certificates, $\min c_I\ge0.05284$; nonlinear closure (Lemma~4c) on the 25 intervals of slabs 2--180, $\min\varrho_c\ge0.04711$; tip box (Lemma~4b) (\fn{t7/constants.json}) & \tC, \S\ref{ss:coer}\\
(Y) & level $\le3.22\cdot10^{-14}$; majorant $\Yb=10^{-12}\ge Y(\tau_a)$, asserted (\fn{ycomp_final.json}) & \tC, \S\ref{ss:Y}\\
far field (T8) & Lemmas F, B, J, N3$'$, slabs 2--180; $B_j\le(7.61\cdot10^{-7},1.76\cdot10^{-9},5.11\cdot10^{-8},7.82\cdot10^{-6})$; $D_2(\tau^*)\le9.40\cdot10^{-12}$ (\fn{amp_A1_w180.json}) & \tC, \S\ref{ss:far}\\
boxes & far boxes at $\Yb=10^{-12}$: $\times6.42$, $\times4.69$, $\times1.99$, $\times1.69$; tip $|e|,|e'|\le1.61\cdot10^{-8},1.64\cdot10^{-8}$ against $10^{-5}$ & \tC, \S\ref{ss:boxes}\\
hand-off (T9) & (H1)--(H4), (H6) for every $\tau^*$ in slab 180 (Table~\ref{tab:handoff}); $t_1\in[4.5978574775,4.5978586144]$ & \tC, \S\ref{s:T9}\\
Lemma~O data & H1--H5 per cylinder and the v3.2 thresholds (Table~\ref{tab:LO}); far H1 closed (arc by arc, \fn{t9_lemmaO_far.py}) & \tC, \S\ref{s:LO}\\
H6, v3.2 & not met under v3.1; v3.2 approved and built: H6 proved in step~3 (Lemma~R proved), conditional on Hypothesis~(F); step-4 closure $\times1.69$, same condition & \S\ref{s:H6}\\\bottomrule
\end{tabular}\end{center}

\subsection{Theorem 4}\label{ss:thm4}
\begin{named}{Theorem 4 (step 4)}
Let $S_1$ be any shrinker in the certified set of step~1, that is, the profile of any zero of the shooting function $M$ in the certified parallelepiped $Y$ of \fn{s1_tight/in_br2.txt} (\S\ref{ss:S1}). Let $\Gamma_0$ be the closed curve
\[
X(\alpha)=X_0'+\frac{L_0}{2\pi}\int_0^\alpha e^{i\theta(a)}\,da,\qquad\theta(\alpha)=-\frac\pi2+\alpha+\sum_{n=1}^{3762}b_n\sin n\alpha,\qquad\alpha\in[0,2\pi],
\]
with the data of \fn{t2/gamma0_final.json}: $b_n$ the stored doubles (\fn{b_hex}) for $n\neq2$; $b_2$ the unique zero of $C(b_2)=\int_0^\pi\cos(\theta(\alpha)+\frac\pi2)\,d\alpha$ in
\[
[-1.00245655773267901818047080,\ -1.00245655773267901818045451];
\]
$L_0=\texttt{0x1.9e9c2d62f0874p+6}=103.65251688569134$; $X_0'=(\texttt{-0x1.bff02d9838f3cp-1},0)$. Let $\gamma$ be the curve diffusion flow with $\gamma(0)=\Gamma_0$. Then:
\begin{enumerate}[label=(\alph*)]
\item $\Gamma_0$ is smooth, embedded and symmetric about the $x$-axis.
\item $\gamma$ exists, and is smooth and symmetric about the $x$-axis, on $[0,t_1]$. Here $t_1=t(\tau^*)$ for a good time $\tau^*$ in base slab~180 and, measured from $\Gamma_0$,
\[
t_1\in[4.5978574775,\ 4.5978586144].
\]
The base curve's maximal curvature at $t_1$ lies in $[2.279\cdot10^4,\ 2.376\cdot10^4]$.
\item There are $\lambda_1>0$ and $c_1\in\mathbb R$ such that $\Gamma_1:=(\gamma(t_1)-c_1e_1)/\lambda_1$ satisfies the hand-off conditions of Theorem~W of step~3, v3.1 as amended by v3.2, with the step-3 modulation $(x_c,\lambda,\alpha)$ at $\tau_1$:
\begin{enumerate}[label=(H\arabic*),leftmargin=2.6em]
\item the modulated decomposition exists, with the three constraints satisfied exactly; $|\alpha|\le1.1739\cdot10^{-4}+7.3\cdot10^{-13}$ and $\sqrt\Lam\le2.29\cdot10^{-8}$ (required: $1.5\cdot10^{-4}$ and $5\cdot10^{-8}$);
\item on $y\in[14,60]$, $(E_0,E_1,E_2)\le(2\cdot10^{-7},10^{-4},0.2)$: the base values are $(8.01\cdot10^{-11},1.34\cdot10^{-10},2.94\cdot10^{-10})$, and the error keeps them below the thresholds for every $\Yb\le3.63\cdot10^{-10}$;
\item on every window $[y,1.5y]$, $y\in[50,253]$, the deviation energies satisfy $E_0\le2\cdot10^{-7}$ and $E_1,E_2,E_3\le10^{-5}$ (base values $3.48\cdot10^{-8}$, $1.27\cdot10^{-7}$, $3.76\cdot10^{-7}$, $9.71\cdot10^{-7}$; the error keeps them below for every $\Yb\le1.575\cdot10^{-11}$); and $K_4\le5.60\cdot10^{-3}$ on $y\in[50,350]$ (required: $0.06$);
\item $\Gamma_1$ is a graph over the $S_1$ ray on $y\in[48,350]$, with $|\theta-\Theta|\le0.078+5.5\cdot10^{-4}$;
\item the Lemma~O data of v3.2, read arc by arc (Table~\ref{tab:LO}): for the cone family ($r_c\in[200,380]$, $X=3$, $M=0.12$) and the far family (centres $\ge376$, $X=2$, $L=113$, $M=0.15$), Lemma~O H2 (slope), H3 (jets), H4 ($\int(h_0^{(5)})^2\le M^2$ per cylinder) and H5 hold, together with the six data thresholds of v3.2 and the far thresholds in $\lambda_1$ units; Lemma~O H1 (graph) holds for both families, arc by arc (\S\ref{ss:LOH1});
\item $\gamma(t_1)$ is a smooth closed immersed curve.
\end{enumerate}
\end{enumerate}
Every bound is a certified upper bound, rounded outward, and holds for every $\tau^*$ in base slab~180 and every $S_1$ in the certified set.
\end{named}

\noindent\textbf{Under v3.1 alone.} Without the amendment, (H5) also asks Lemma~O H6 (the a-priori bounds close at every centre), and with the v3.1 far family this is not met (\S\ref{ss:H6v31}). The amendment v3.2 re-specifies the far family at $(M,K)=(0.15,0.3)$, adds six data thresholds to (H5), and proves H6 in step~3 from them and from the strip statement of Lemma~S(B) (\S\ref{ss:v32}). It has been approved by the authors and is built in \fn{writeup/step3_proof_v32/}, \fn{writeup/lemmaO_v32/} and \fn{writeup/step3_amend_v32/}; Lemma~R is proved, and the step-3 gap in Lemma~S(B) is covered by Hypothesis~(F) (\S\ref{ss:v32status}). Theorem~4(c) delivers the v3.2 data, which do not depend on (F). The v3.1 values are recorded in \S\ref{ss:LOv31}.

\subsection{What is proved, and what is open}\label{ss:status}
\begin{itemize}
\item \textbf{(a)} is certified (Proposition~\ref{prop:G0}).
\item \textbf{(b)} Up to $t_{96}$, existence, uniqueness and regularity are certified by the Galerkin-limit Lemma~\ref{lem:gal}. From $t_{96}$ to $t_1$ they follow from Theorem~P2 (\S\ref{s:P2}), which is now certified in full: the nonlinear terms are inside the certified closure (Lemma~4c) and the tip box is certified (Lemma~4b). Continuation uses the literature fact (P1) of step~3 \tA. The interval for $t_1$ and the base $k_{\max}$ are produced by \fn{t1clock.py} \tC; The script includes the scale factor $(\lambda/\tilde\lambda)^4$ and $x_{\rm NL}$ slab by slab (Remark~\ref{rem:t1}); the quoted interval stands. The true curve's $k_{\max}$ at $t_1$ is \emph{not} certified: it needs a pointwise tip bound on $e''$ \OPEN{\S\ref{s:open}}. No hypothesis of step~3 uses it.
\item \textbf{(c)} (H1)--(H4) are certified for every $\tau^*$ in slab 180 (\S\ref{s:T9}); (H6) follows from (b) \tA. In (H5), Lemma~O H2--H5, the six data thresholds and the far thresholds are certified for both families, and H1 for both families (for the far family arc by arc, by a pass/fail script, \S\ref{ss:LOH1}). Lemma~O H7 (time ratio) and H9 (regularity) are internal to step~3 and hold there. H8 is not required. Lemma~O H6 is no longer a hand-off check: step~3 v3.2 proves it (D5b) from these data and from the strip statement of Lemma~S(B) at levels 0--2, whose level-2 part uses Lemma~R, which is proved. The strip statement is part of Hypothesis~(F), so H6 in step~3, and step~4's certified closure margins for it (Table~\ref{tab:v32}), are conditional on (F) (\S\ref{ss:v32status}).
\end{itemize}
Consequently the main theorem (\S\ref{s:main}) holds once the far-family H1 item is closed and Lemma~R is proved.

\begin{table}[htbp]\centering\small
\caption{Hand-off items, certified uniformly for $\tau^*$ anywhere in base slab 180 ($t_c\in[-1,1]$, $\Delta\tau=0.15$, $k\in[2.27,2.38]\cdot10^4$). Base values in \fn{arb}; the error enters at $\Yb=10^{-12}$. A $\Yb$-limit is the largest $\Yb$ for which the item holds. Files are in \fn{step4/p2/out/} unless marked.}\label{tab:handoff}
\begin{tabular}{@{}L{0.13\linewidth}L{0.19\linewidth}L{0.47\linewidth}L{0.13\linewidth}@{}}\toprule
item & requirement & certified (file) & margin\\\midrule
(H1) $\alpha$ & $|\alpha|\le1.5\cdot10^{-4}$ & $|\tilde\alpha|\le1.1739\cdot10^{-4}$, plus $C_\alpha\Yb\le7.3\cdot10^{-13}$ (\fn{h1slab/summary.json}, \fn{t9_h1_clam_s180.json}) & \thin{$\times1.27$}\\
(H1) $\sqrt\Lam$ & $\le5\cdot10^{-8}$ & base $\le2.282\cdot10^{-8}$; with error $\le2.29\cdot10^{-8}$ ($C_\Lam\Yb\le2.42\cdot10^{-11}$); holds for $\Yb\le1.12\cdot10^{-9}$ & $\times2.18$\\
(H2) shells, $y\in[14,60]$ & $E_{0,1,2}\le(2\cdot10^{-7},10^{-4},0.2)$ & base $(8.01\cdot10^{-11},1.34\cdot10^{-10},2.94\cdot10^{-10})$; $\Yb$-limit $3.63\cdot10^{-10}$ ($E_2$; own bound by $\Lam_e$) (\fn{h1slab/win_summary.json}, \fn{t9_cone_kch.json}) & $\times363$ in $\Yb$\\
(H3) windows & $E_0\le2\cdot10^{-7}$, $E_{1,2,3}\le10^{-5}$, $y\in[50,253]$ & base $(3.48\cdot10^{-8},1.27\cdot10^{-7},3.76\cdot10^{-7},9.71\cdot10^{-7})$; $\Yb$-limit $1.575\cdot10^{-11}$ ($E_3$, binding) (\fn{t8se/constants/amp_A1_w180.json}) & $\times15.7$ in $\Yb$\\
(H3) $K_4$ & $\le0.06$ on $y\in[50,350]$ & base $3.1633\cdot10^{-4}$; total $\le5.60\cdot10^{-3}$ (\fn{t9_cone_kch.json}) & $\times10.7$\\
(H4) graph & over the $S_1$ ray, $y\in[48,350]$ & $|\theta-\Theta|\le0.078$ (base) $+\,5.5\cdot10^{-4}$ (error) & ---\\
(H5) & Lemma~O data & Table~\ref{tab:LO} & \S\ref{s:LO}\\
(H6) & smooth, closed, immersed & Theorem~4(b) & \tA\\\bottomrule
\end{tabular}\end{table}

\subsection{Margins}\label{ss:margins}
Table~\ref{tab:margins} lists the certified margins, thinnest first within each group. A margin is threshold over bound, rounded down. The highlighted entries are the thinnest of the proof. The step-3 side of v3.2 is listed for reference; it is not a step-4 deliverable.

\begin{table}[htbp]\centering\small
\caption{Certified margins. ``In $\Yb$'': the item holds for all $\Yb$ up to the stated limit, compared with the $\Yb=10^{-12}$ in use.}\label{tab:margins}
\begin{tabular}{@{}L{0.42\linewidth}L{0.21\linewidth}L{0.11\linewidth}L{0.10\linewidth}l@{}}\toprule
item & bound & threshold & margin & \S\\\midrule
\multicolumn{5}{@{}l}{\emph{P2 bootstrap}}\\
initial error $Y(\tau_a)$ against $\Yb$ & $6.21\cdot10^{-13}$ & $10^{-12}$ & \thin{$\times1.6$} & \ref{s:switch}\\
far box $|e|$: $\sqrt{F_0}$ ratio & $1.00001$ & $1.7$ & \thin{$\times1.69$} & \ref{ss:boxes}\\
far box $x^2|e'''|$ & $8.34\cdot10^{-3}$ at $\Yb=1.7\cdot10^{-12}$ & $10^{-2}$ & $\times1.99$ at $10^{-12}$ & \ref{ss:boxes}\\
far box $x|e''|$ & $1.08\cdot10^{-4}$ at $1.7\cdot10^{-12}$ & $3\cdot10^{-4}$ & $\times4.69$ at $10^{-12}$ & \ref{ss:boxes}\\
far box $|e'|$ & $2.38\cdot10^{-4}$ at $1.7\cdot10^{-12}$ & $1.5\cdot10^{-3}$ & $\times6.42$ at $10^{-12}$ & \ref{ss:boxes}\\
Lemma~J rate condition $\lambda<2\mu_1$ (a free parameter; see \S\ref{ss:far}) & $\lambda_{\max}=0.2520$ & $0.30$ & $\times1.19$ & \ref{ss:far}\\
\multicolumn{5}{@{}l}{\emph{Hand-off (H1)--(H4)}}\\
(H1) $|\alpha|$ & $1.1739\cdot10^{-4}+7.3\cdot10^{-13}$ & $1.5\cdot10^{-4}$ & \thin{$\times1.27$} & \ref{ss:alpha}\\
(H1) $\sqrt\Lam$ & $2.29\cdot10^{-8}$ & $5\cdot10^{-8}$ & $\times2.18$ & \ref{ss:alpha}\\
(H3) $K_4$, $y\in[50,350]$ & $5.60\cdot10^{-3}$ & $0.06$ & $\times10.7$ & \ref{ss:cone}\\
(H3) windows ($E_3$ binding) & $\Yb$-limit $1.575\cdot10^{-11}$ & $\Yb=10^{-12}$ & $\times15.7$ in $\Yb$ & \ref{ss:win}\\
(H2) shells ($E_2$ binding) & $\Yb$-limit $3.63\cdot10^{-10}$ & $\Yb=10^{-12}$ & $\times363$ in $\Yb$ & \ref{ss:win}\\
\multicolumn{5}{@{}l}{\emph{Lemma~O data, v3.2 (own cylinder units unless marked)}}\\
cone slope datum $s_0$ & $0.05724$ & $0.08$ & \thin{$\times1.40$} & \ref{ss:LOjets}\\
cone $\sup|h_0''|$ datum $m_2$ & $0.01883$ & $0.03$ & \thin{$\times1.59$} & \ref{ss:LOjets}\\
far slope datum $s_0$ & $0.04153$ & $0.08$ & $\times1.92$ & \ref{ss:LOjets}\\
far $|h_{xx}|$ ($\lambda_1$ units); far $m_2$ & $1.797\cdot10^{-4}$; $0.02031$ & $3.53\cdot10^{-4}$; $0.04$ & $\times1.96$ & \ref{ss:LOjets}\\
cone $\|h_0'''\|_{L^2}$ datum $n_3$ & $1.170\cdot10^{-3}$ & $0.004$ & $\times3.41$ & \ref{ss:LOjets}\\
far slope against $M=0.15$ & $0.0416$ & $0.15$ & $\times3.61$ & \ref{ss:LOjets}\\
cone Lemma~O H4 (worst, $r_c\in[370,380]$) & $0.2393\,M^2$ & $M^2$ & $\times4.18$ & \ref{ss:H4}\\
far Lemma~O H4 at $M=0.15$ (centre 376) & $0.1761\,M^2$ & $M^2$ & $\times5.68$ & \ref{ss:H4}\\
far $\|h_{xxx}\|_{L^2}$ ($\lambda_1$ units); far $n_3$ & $5.516\cdot10^{-7}$; $6.626\cdot10^{-4}$ & $3.32\cdot10^{-6}$; $0.004$ & $\times6.01$; $\times6.03$ & \ref{ss:LOjets}\\
cone $K_2$; cone $K_4$ & $1.89\cdot10^{-2}$; $1.23\cdot10^{-2}$ & $0.12$ & $\times6.3$; $\times9.7$ & \ref{ss:LOjets}\\
far $|h_{xxxx}|$ ($\lambda_1$ units) & $6.964\cdot10^{-9}$ & $1.03\cdot10^{-7}$ & $\times14.79$ & \ref{ss:LOjets}\\
cone $K_3$; far $|h_{xxx}|$ ($\lambda_1$ units) & $4.78\cdot10^{-4}$; $2.595\cdot10^{-8}$ & $0.12$; $1.17\cdot10^{-5}$ & $\times251$; $\times450.9$ & \ref{ss:LOjets}\\
\multicolumn{5}{@{}l}{\emph{Step-3 side of v3.2 (for reference; \fn{t9_h6_v32.json}); the closure rows are conditional on Hypothesis~(F)}}\\
closure: cone $K_2$ ($r_c=380$, $r_p=150$) & $0.10664$ & $0.18$ & $\times1.69$ & \ref{ss:v32}\\
closure: cone slope; cone $K_3$ & $0.13532$; $0.09348$ & $0.25$; $0.18$ & $\times1.85$; $\times1.93$ & \ref{ss:v32}\\
closure: far $K_2$; far $K_3$; far slope & $0.13125$; $0.12549$; $0.08753$ & $0.3$; $0.3$; $0.25$ & $\times2.29$; $\times2.39$; $\times2.86$ & \ref{ss:v32}\\
Lemma~O $t^*$: cone; far at $(0.15,0.3)$ & $3.683\cdot10^{-7}$; $9.511\cdot10^{-8}$ & $1.6001\cdot10^{-7}$; $6.134\cdot10^{-9}$ & $\times2.30$; $\times15.5$ & \ref{ss:v32}\\\bottomrule
\end{tabular}\end{table}

\subsection{The review and its resolution}\label{ss:review}
\begin{table}[htbp]\centering\small
\caption{Findings of the review (\fn{STEP4_REVIEW.md}, Nov~11) and their resolutions (\fn{STEP4_RESULT.md} \S3).}\label{tab:review}
\begin{tabular}{@{}L{0.05\linewidth}L{0.27\linewidth}L{0.53\linewidth}L{0.07\linewidth}@{}}\toprule
id & finding & resolution & \S\\\midrule
A1 & Lemma~O H4, $\int(h_0^{(5)})^2\le M^2$, not delivered; on the window $W'$ of length $1.46\cdot10^{-4}$ it failed by $\times5$ to $\times280$ & The good-time window is all of slab 180. Every hand-off item is certified uniformly in $\tau^*$ (for H1: a 16-point Krawczyk/IFT test and 32 derivative balls). $D_2(\tau^*)\le9.40\cdot10^{-12}$. H4 per cylinder: $\le0.2393\,M^2$ (cone), $\le0.1761\,M^2$ (far, $M=0.15$) & \ref{ss:window}, \ref{ss:H4}\\
A2 & nonlinear terms of (Y) only in the allowance $\mathrm{DN}=10^{-6}$ & T7 Lemma~4c: 55 explicit terms per interval inside the certified closure; $\min\varrho_c^{\rm NL}\ge0.04711$ (losses $\le8\cdot10^{-5}$); $x_{\rm NL}\le1.95\cdot10^{-15}$ in (M) & \ref{ss:lem4c}\\
A3 & tip box $2\cdot10^3\Yb/\sqrt{\mu_e}$ not certified; its constant is too small near $x=34.25$ & T7 Lemma~4b: $|e|\le1.6031\cdot10^4\Lam_e^{1/2}$, $|e'|\le1.6352\cdot10^4\Lam_e^{1/2}$ on $[0,34.25]$; used by \fn{merge.tip_lemma}, \fn{n3far}, \fn{clam.py}, \fn{n3boxes.py} and the Lemma~O scripts & \ref{ss:lem4b}\\
B1 & \fn{se_corr.py}'s majorants not derived; they set the far residuals & The residual of the (SE)-corrected base is computed directly by T6's engine (\fn{p2/t6c}). \fn{se_corr} is retired; its $\cR_0$ increment was invalid. $a_2$ is unique by Banach & \ref{ss:res}\\
B2 & box margins reported against the wrong boxes & \fn{n3boxes.py} reads the boxes from the N3$'$ output and asserts all four: $\times6.42$, $\times4.69$, $\times1.99$, $\times1.69$ & \ref{ss:boxes}\\
B3 & H1 by a Newton step with a typed remainder & Krawczyk test at 16 exact points; exact $u''''$ cell hull; no fall-back bound (any failure is an error); derivative balls & \ref{ss:alpha}\\
B4 & typed base values in \fn{camp.py} & \fn{camp.py} reads the JSON outputs (\fn{T9WIN}, \fn{t9_cone_kch.json}) & \ref{ss:win}\\
B5 & typed allowances in \fn{t9.py}, \fn{t9_far.py} & KCH, $c_0$ and $c_0^{\rm tip}$ computed exactly (\fn{p2/t9chart.py}, arb Taylor mode); exact normal-graph jets (\fn{curvjets.py}); tip-anchored $y_{\rm ray}$; shells by their own $\Lam_e$ bound & \ref{ss:win}\\
B6 & $|x|\le C_x\Yb$ omits $x_r$ and the nonlinear part of (M) & $x_r\le2.0\cdot10^{-16}$ added; $x_{\rm NL}\le1.95\cdot10^{-15}$ (Lemma~4c) is at most $8.0\cdot10^{-4}$ of the slack $C_x(\Yb-Y_{\rm cert})\ge2.44\cdot10^{-12}$ & \ref{ss:mod}\\
B7 & unscripted: $Y_0$ in (Y); placement of $F_j(\tau_0)$; $t_1$, $k_{\max}$; (H6) & \fn{ycomp.py} runs with $Y_0=6.21\cdot10^{-13}$ and asserts $\Yb\ge Y_0$; \fn{camp.py} starts the $F_j$ at the switch ($\mathrm{SWF}=0.9365$); \fn{t1clock.py} gives $t_1$ and the base $k_{\max}$; (H6) is stated & \ref{s:switch}, \ref{ss:t1}\\
B8 & $S_1$ is not proved unique & Theorem~4 is stated for any $S_1$ in the certified set; \fn{check_outA5.py}: all $1025+1025+501$ rows lie inside the step-3 hulls & \ref{ss:S1}\\
B9 & rounding spots & outward \fn{nextafter} in \fn{switch.py} and \fn{ycomp.py}; exact dyadic cells in \fn{h1.py}; T6 hulls widened by 2 ulps on read & \ref{app:arb}\\\bottomrule
\end{tabular}\end{table}

\section{\texorpdfstring{The initial curve $\Gamma_0$}{The initial curve Gamma0}}\label{s:G0}
\emph{Verified by} \fn{t2/gamma0.py final} (input \fn{t2/gamma0s_src.npz}); \emph{certificate} \fn{t2/gamma0_final.json}, log \fn{t2/final.log}; design in \fn{t3/T3_REPORT.md} \S1. (\fn{T2_REPORT.md} describes the unslaved predecessor; its closure width $4\cdot10^{-16}$ is superseded.)

\paragraph{Definition.} $\Gamma_0$ is the curve of Theorem~4. Its tangent angle is the exact sine series $\theta=-\pi/2+\varphi$, $\varphi=\alpha+\sum_{n\le3762}b_n\sin n\alpha$. The point $\alpha=0$ is an axis crossing with $\theta(0)=-\pi/2$, and the half-curve $\alpha\in[0,\pi]$ lies in $\{y<0\}$. The coefficients are exact doubles except $b_2$, which is the exact root of the closure condition $C(b_2)=\int_0^\pi\cos\varphi=0$ (then $\operatorname{Im}X(\pi)=0$). The curve was obtained from the float snapshot of the clip run at $t=0.001$ (\fn{p2_t0.001.npz}) by truncation and two slaving passes. These steps define $\Gamma_0$; they are not part of the proof, and the closeness of $\Gamma_0$ to the snapshot is not certified (and not needed).
\begin{itemize}
\item \emph{Truncation.} $M=3762$, chosen by $\sum_{n>M}n^4\abs{b_n}\le10^{-12}$ (tail $9.96\cdot10^{-13}$). The discarded even part has $\max\abs{a_n}=7.3\cdot10^{-11}$.
\item \emph{Slaving} (T3 option (c)). The float snapshot is off the slow manifold of the Galerkin flow ($10^{-10}$ at $n=186$--$350$). A polynomial slab cannot carry that layer: a slab $\Delta t=t/2$ from the raw curve has residual $2\cdot10^{-3}$. Pass~1 replaces modes $186..1000$ by the $s=0$ values of a free-start Lobatto collocation on $[0,1.8t]$ ($\abs{\Delta b_n}\le1.6\cdot10^{-10}$). Pass~2 replaces modes $\ge250$ by the $s=0$ values of slab~1's own base (largest change $5.5\cdot10^{-18}$, at $n=250$), so slab~1 starts with mismatch $9.4\cdot10^{-23}$. Modes $\le185$ and $>1000$ keep the snapshot values.
\end{itemize}

\begin{proposition}[$\Gamma_0$ \tC]\label{prop:G0}
(i) \emph{Closure.} Validated quadrature (Taylor models of degree $24$ per arc with rigorous remainder, 512 leaves) gives $C(m)\in[\pm4.07\cdot10^{-24}]$ at the centre $m$. On $X=m\pm8.137\cdot10^{-24}$ (the interval of Theorem~4) the Krawczyk operator satisfies $K(X)\subset\operatorname{int}X$, with $\abs{1-C'(X)/c'}\le3.01\cdot10^{-12}$ ($c'=-1.99704$). So $C$ has a unique zero $b_2^*\in X$, and $\Gamma_0$ is closed. Also $\operatorname{Re}X(\pi)=-0.8227219774300355680063\pm3.71\cdot10^{-23}$ and $\abs{\operatorname{Im}X(\pi)}\le3.29\cdot10^{-22}$.
(ii) \emph{Embeddedness.} With $b_2\in X$, the half-curve splits into 701 arcs, each of total turning $\le0.04997<\pi/2$, with adjacent pairs turning $\le0.0960<\pi$; so each arc and each adjacent pair is injective. The boxes of non-adjacent arcs are pairwise disjoint after 3 refinement rounds (minimal $L^\infty$ separation $8.89\cdot10^{-3}$). The middle arcs $49..655$ have certified $y<0$. On the left zone (arcs $0..48$) $\varphi\in(-1.5,1.5)$, so $y$ strictly decreases from $y(0)=0$. On the right zone ($656..700$) $\varphi\in(1.6,4.7)$, so $y(\alpha)=\frac{L_0}{2\pi}\int_\alpha^\pi\cos\varphi<0$, using $y(\pi)=0$, i.e.\ $C(b_2^*)=0$. Hence the half-curve is injective and meets the axis exactly at its two endpoints. These are distinct ($X(0)\approx-0.8749$, $X(\pi)\approx-0.8227$). With the mirror image, $\Gamma_0$ is embedded. The slot half-gap is $\ge0.0248219$ (float $0.0248223$).
\end{proposition}
\noindent The disjointness test rounds the \fn{arb} bounds outward to doubles with \fn{nextafter}. The right-zone sign argument uses the exact closure at $b_2^*$, not the box bound $\abs{\operatorname{Im}X(\pi)}\le3.3\cdot10^{-22}$. The review checked that $M_0=1000$ and $N_c=4000\ge3762$ in P1's first slab, so every $\Gamma_0$ mode is in the centre or in the box; that the $b_2$ interval enters the initial radii correctly; and that $L_0$ matches. Kit~A reproduces this certificate bit for bit on Windows~11 (\S\ref{app:kits}).

\section{P1: validated Fourier slabs}\label{s:P1}
\emph{Verified by} \fn{t3/run_p1.py} with the engine \fn{t3/slab.py}; \emph{certificates} \fn{t3/p1/ckpt/slab_NNNNN.json} and \fn{t3/p1/status.json}; design \fn{t3/T3_DESIGN.md} \S0--6, report \fn{t3/T3_REPORT.md}.

\subsection{Equation}
Let $\theta=-\pi/2+\alpha+\mathrm{th}$, $\mathrm{th}=\sum_{n\ge1}u_n\sin n\alpha$, and clock $\tau'$ with $d\tau'=(2\pi/L)^4dt$. Then, with $\lambda_n=n^4-n^2$,
\[
u_n'=-\lambda_nu_n+N_n(u),\qquad N=-\partial_\alpha\bigl(v^2+\tfrac13v^3\bigr)+(1+v)\,\partial_\alpha^{-1}\bigl(w^2-\ip{w^2}{}\bigr),\quad v=\mathrm{th}_\alpha,\ w=\mathrm{th}_{\alpha\alpha},
\]
where $\partial^{-1}$ is the zero-mean antiderivative. This is the CDF equation after one integration by parts, checked against \fn{hls.py} in floats. The slaved quantities are $(\ln L)'=-\frac12\sum n^4u_n^2$, $t'=(L/2\pi)^4$ and $X_0'=(L/2\pi)(\sum n^3u_n,0)$. They are propagated by exact Chebyshev algebra in \fn{arb}.

\subsection{One slab}
On a slab $s\in[0,h]$, with $h\,\theta_{\max}^4\approx2$:
\begin{enumerate}[label=(\arabic*),leftmargin=*]
\item \emph{Base.} $\bar u$ is a Chebyshev--Lobatto polynomial of degree $K=24$ ($K=16$ in the swing) in $s$, for $M$ modes. It is computed by a float collocation Newton step and then polished three times against the \fn{arb} collocation residual. The centre is carried as exact 128-bit midpoints. The residual $r=\bar u'+\Lambda\bar u-N(\bar u)$ (degree $3K$) is enclosed exactly.
\item \emph{Error split.} $e=u-\bar u=P(s)\rho+z(s)$ with $\abs{\rho_i}\le w_i$. Here $P$ solves the variational equation on the slow block $n\le M_a$, $\lambda_{M_a}h\le8$, with slaved high rows. Its defect $D=P'+\Lambda P-JP$ is enclosed in \fn{arb}.
\item \emph{Bounds on $z$.} $z'=-\Lambda z+Jz+Q(e)-r-D\rho$. On ten sub-intervals, componentwise Duhamel bounds with exact $e^{-\lambda_ns}$ are used for $n\le N_c=4M$, with a tail $T\,n^{-6}$ beyond. A strict \fn{arb} isolation check (the first-exit argument of Lemma~\ref{lem:gal}) closes the bounds with a 3\% margin. Intra-slab absorption moves the $z_a$ part into $\rho$.
\item \emph{Lohner step} at $s=h$: QR of $P_a(h)\,\mathrm{diag}(w)$ with column sorting; new $w'$, box and tail. Block and $M$ are resized by exact changes of representation.
\end{enumerate}

\subsection{The Galerkin-limit lemma}
Let $Y=(Z_1,\dots,Z_{N_c},T(N_c+1)^{-6},\dots)$ and $W=\{z:\abs{z_n}\le Y_n\ \forall n\}$. Let $\hat J,\hat Q$ be the componentwise majorants of \fn{T3_DESIGN.md} \S3, so that $\abs{(Jz)_n}\le\hat J(Y)_n$ and $\abs{Q_n(s,P\rho+z)}\le\hat Q(Y)_n$ for $z\in W$, $s\in I=[0,\Delta]$. Let $\abs{r_n}\le\hat r_n$, $\abs{(D\rho)_n}\le\hat d_n$, and $F_n:=\hat J(Y)_n+\hat Q(Y)_n+\hat r_n+\hat d_n$. Put $g_n(t)=e^{-\lambda_nt}a_n+\phi_n(t)F_n$ with $\phi_n(t)=(1-e^{-\lambda_nt})/\lambda_n$.
\begin{lemma}[first exit and Galerkin limit; \fn{T3_DESIGN.md} \S6]\label{lem:gal}
Assume the isolation condition \textup{(H)}: $Y_n>\max(a_n,g_n(\Delta))$ for all $n$. This is checked in \fn{arb} componentwise for $n\le N_c$, and through the strict tail inequality on $T$ for $n>N_c$. Let $u(0)$ be data in the sine class with $z(0)=u(0)-\bar u(0)-P(0)\rho$, $\abs{z_n(0)}\le a_n$. Then there is a unique solution $u\in C^0(I;H^5)\cap C^1(I;L^2)$ of $u_s=-(\partial^4+\partial^2)u+N(u)$ in the sine class, with $z(s)\in W$ for all $s\in I$ and $\abs{z_n(\Delta)}\le g_n(\Delta)$.
\end{lemma}
\noindent The proof is in Appendix~\ref{app:gal}. It constructs the solution, so chaining sub-intervals and slabs gives existence and uniqueness on all of P1; no external existence theorem is used before $t_{96}$. The review found the lemma complete as written.

\subsection{Results \tC}
P1 was run from $\Gamma_0$ through 96 validated slabs: graded slabs $\Delta t=t/2$ to $t\approx0.03$, then $h=2/\theta_{\max}^4$, the swing and the approach. The run used the M-step rule $\max(80,0.15M)$. It stopped at slab 95 ($k_{\max}\ge25$, certified $26.2537786005$) and was then extended by one slab to the switch state.
\begin{center}\small\begin{tabular}{@{}ll@{}}\toprule
quantity at slab 96 (\fn{t3/p1/status.json}) & value\\\midrule
time since $\Gamma_0$ & $t_{96}=4.597850539485830831\pm7.8\cdot10^{-19}$\\
$L$ & $48.29333418442305\pm3.7\cdot10^{-15}$\\
$X_0$ & $0.84121611132611\pm4.3\cdot10^{-15}$\\
$k_{\max}$ & $[29.352874545325992,\,29.352874545327012]$\\
$M$, block $M_a$, $\theta_{\max}$ & $2800$ ($3200$ next), $285\to319$, $225.61$\\
$\abs e_j=\sum n^j\abs{e_n}$, $j=0..4$ & $4.58\cdot10^{-15}$, $4.68\cdot10^{-13}$, $6.87\cdot10^{-11}$,\\ & $1.29\cdot10^{-8}$, $2.97\cdot10^{-6}$\\
physical $C^2$, $C^5$ (lead) & $1.97\cdot10^{-13}$, $8.50\cdot10^{-10}$\\
residual (truncation part), CPU & $2.52\cdot10^{-11}$, $6.24$ h (1 core)\\\bottomrule\end{tabular}\end{center}
The T5 design asked for $C^2\le3\cdot10^{-14}$ at the switch. It is not met, and it is not needed: the switch (\S\ref{s:switch}) computes the modulated initial error directly from the P1 enclosure. P1 is a computation in Fourier variables, so it is insensitive to the loss of embeddedness that happens in its first slabs (\S\ref{ss:embed}). Kit~A reproduces P1 on Windows~11: 96 of 96 slabs validated, final enclosures overlapping (\S\ref{app:kits}).

\section{The base}\label{s:base}
\emph{Sources:} \fn{t1/POLISH_V2.md}, \fn{t1/polished_v2/manifest.json}, \fn{STEP4_DESIGN.md} \S21--22. The base is exact data; nothing about it is claimed beyond what T6 encloses (\S\ref{ss:res}).

The P2 base is a half curve $\Gt(x,\tau)$, $x\in[0,\lt(\tau)]$, given by its curvature $\kt$. In $x$, $\kt$ is a $C^4$ piecewise polynomial with exact symmetry at both ends ($\kt'=\kt'''=0$ at $x=0,\lt$). The discretisation is weak (Legendre--Galerkin), on a geometric mesh whose outer knots move with $\lt$. In $\tau$ it is degree-8 Radau collocation per slab, continuous across slabs. Exact quadrature then gives $\Gt$ and the scalars $\tilde a,\tilde b$.

\paragraph{Seed.} Slab~0 is P1 slab~95 \emph{exactly}. The 2800-term series of slab~95 was evaluated in \fn{arb} at 192 bits, by blocked \fn{arb_mat} products: $\sum_qe^{iqB\alpha}\sum_rc_{qB+r}e^{ir\alpha}$ with $B=64$, because Horner on $\abs z=1$ inflates radii like $(\abs{\cos}+\abs{\sin})^n$. This gives $k=26.2537786005622$, $\lambda=K_0/\kappa_{\rm tip}$, $\lt=411.97336$ and a mesh of $N=120$ elements. The nodal $\kt^{(0..4)}$ are exact at the knots. The bubbles match the Legendre moments 0 and 1. A zero-mean bump of size $7.6\cdot10^{-23}$ removes the far-point offset $3.4\cdot10^{-21}$. The seed is stored as exact 128-bit dyadics. Sampled agreement, at 300 random points: $\abs{\kappa_h-\kappa_{P1}}\le4.7\cdot10^{-17}$ for $x\le40$ \tF. This check is not used: the switch encloses the mismatch rigorously.

\paragraph{Store.} The store \fn{polished_v2/} has 181 slabs, $\tau\in[0,27.15]$, $\Delta\tau=0.15$, $k$ from $26.2538$ to $23718.7$, and $N$ from $120$ to $171$ elements ($844\to1201$ DOF). The maximal polish residual is $9.7\cdot10^{-27}$. Fresh-process reloads of slabs 0, 1, 160, 170, 180 give bit-identical norms. A server restart interrupted the polish after slab 164; it was resumed losslessly, with Newton paths that differ only within $\abs E<10^{-26}$. At slab 180, $\lt\in[3.56,3.70]\cdot10^5$.

\paragraph{(SE) defect of the stored base.} At slab ends the tilt is $\le4.2\cdot10^{-28}$. Inside slabs, at the Radau stage rows, the tilt is $\le3.2\cdot10^{-18}$ and the far-point offset $h=\operatorname{Im}\Gt(\lt)$ is $\le1.05\cdot10^{-15}$ (rescaled). The cause is that $\int\kt$ and $h$ are nonlinear in the moving outer mesh, and $\mathrm{Leg}_0$ conserves $\int\kt$ only at the collocation nodes. A comoving correction makes (SE) exact; its residual is computed directly by T6's engine (\S\ref{ss:res}).

\begin{remark}[development notes]
(1) \emph{Mis-seeded base, caught and fixed.} The first base (\fn{t1/polished}) was seeded from the float snapshot \fn{p2_k20} of the original clip's HLS run, not from P1's validated trajectory. The assembly found this at the switch (Nov~10): the float mismatch was $Y(\tau_0)\approx2.8\cdot10^{-9}$, $\times173$ over the budget of the time, and P2 could not close on that base. The base was rebuilt from P1's exact centre (\fn{polished_v2}); the certified modulated mismatch is now $6.21\cdot10^{-13}$. All T6, T7 and T8 results on the old base are archived and superseded. The lesson: the base must continue the validated trajectory.
(2) \emph{Hand-off scale.} The float go/no-go (milestone 1) met the outer-window thresholds of step 3 only from $k\approx1.6$--$2.3\cdot10^4$. The deviation there is $-40(y/k)^3$, so the window energies scale like $k^{-5.9}$. The hand-off was therefore raised to $k\approx2.3\cdot10^4$, and step 3 moved to $\alpha_1=1.5\cdot10^{-4}$.
\end{remark}

\section{The switch}\label{s:switch}
\emph{Verified by} \fn{p2/switch.py} (stage ``switch'' of the final tail, \fn{T6DIR=p2/t6c}); \emph{certificate} \fn{p2/out/switch_s2.json}, log \fn{switch_final.log}; the manifest check \fn{switch_vs_typed} confirms the typed inputs of later stages.

\begin{lemma}[switch \tC]\label{lem:switch}
Let $u$ be the P1 solution at the end of slab 96 (\fn{t3/p1/ckpt/slab_00096.json}), at time $t_{96}$. Take base slab 2 at $t_{c,a}=982986695892350/2^{50}=0.873067570143915$ ($\tau_a=0.44048$; $k_{\rm tip}=29.3528745453264\pm4.8\cdot10^{-14}$). Then there are $(\lambda_0,c_0,\tau_0)$ in the a-priori ball of the implicit function theorem such that $\gamma(t_{96})$ has the modulated decomposition of Theorem~P2 over the base, with
\[
Y(\tau_a)\le6.21\cdot10^{-13},\qquad F_j(\tau_a)\le(3.75\cdot10^{-23},\ 2.45\cdot10^{-22},\ 4.72\cdot10^{-19},\ 3.19\cdot10^{-15})\quad(j=0..3).
\]
\end{lemma}
\noindent\emph{Proof.} The P1 ball is converted into the base chart: $\alpha=\pi-s\,x$, $\kappa=-s\theta_\alpha$, $\lambda_{P}=K_0/\kappa_{\rm tip}$, ball norms with the $n^{-6}$ tail, the length factor, and 1902 cells. The physical times agree: the base clock at the switch is $t_{96}$ to within $2.6\cdot10^{-20}$. An a-priori ball feeds the IFT for $(\lambda_0,c_0,\tau_0)$, and the energies are evaluated in \fn{arb} on the modulated error.\hfill$\square$
\begin{center}\small\begin{tabular}{@{}L{0.30\linewidth}L{0.42\linewidth}L{0.22\linewidth}@{}}\toprule
quantity & value & role\\\midrule
IFT: $\bar m$, contraction $q$, shifts & $5.104$, $4.18\cdot10^{-9}$, all $\le1.26\cdot10^{-12}$ (clock shift $\le1.24\cdot10^{-12}$) & start inside its a-priori ball\\
raw $Y$ (no modulation) & $2.19\cdot10^{-14}$ & ---\\
modulated $Y(\tau_a)$ & $6.2032\cdot10^{-13}$ ($A\le6.203\cdot10^{-13}$, $B\le1.512\cdot10^{-12}$) & initial value of (Y)\\
$F_j(\tau_a)$, $j=0..3$ & $3.75\cdot10^{-23}$, $2.45\cdot10^{-22}$, $4.72\cdot10^{-19}$, $3.19\cdot10^{-15}$ & rounded up into \fn{--F0} $4,3,5,3.4$\\\bottomrule\end{tabular}\end{center}
The modulation moves the frame onto the constraint manifold, which is why the modulated $Y$ is larger than the raw one. The margin against $\Yb=10^{-12}$ is $\times1.6$. Review B7: \fn{ycomp.py} now receives $Y_0=6.21\cdot10^{-13}$ and asserts $\Yb\ge Y_0$, and \fn{camp.py} starts the $F_j$ at the switch point, a fraction $\mathrm{SWF}=(1+t_{c,a})/2=0.9365$ through slab 2. Review B9: the weight arguments in \fn{switch.py} are rounded outward with \fn{nextafter}. Rerun on the direct residuals of \fn{p2/t6c}, the switch is unchanged: $Y(\tau_a)\le6.2032\cdot10^{-13}$ and the same $F_j$.

\section{P2: energy shadowing}\label{s:P2}
\emph{Sources:} \fn{t5/T5_DESIGN.md} (design; tags \tD), \fn{t7/T7_REPORT.md}, \fn{t8/LEMMA_F.md}, \fn{t8/T8_REPORT.md}, \fn{t6/T6_REPORT.md}, \fn{p2/*.py}, \fn{STEP4_RESULT.md}. The statement below consolidates T5 with the choices made later: the fixed late-base constraint functions of T7, levels $j\le3$ only, the boxes of Lemmas 4a, 4b and N3$'$, the nonlinear closure of Lemma 4c, and the (SE)-corrected base with its directly computed residual.

\subsection{Objects}\label{ss:objects}
\emph{Base.} Let $\hat\kappa$ be the (SE)-corrected base curvature (\S\ref{ss:res}), and $\Gt,\tilde a,\tilde b$ the corresponding curve and scalars. Put $c:=\frac14-\tilde a$. Derived quantities: the frame $\tilde T,\tilde\nu$; $\tilde p=\ip\Gt{\tilde\nu}$, $\tilde q=\ip\Gt{\tilde T}$, $\tilde n=\ip{e_1}{\tilde\nu}$, $\tilde t_1=\ip{e_1}{\tilde T}$; the normal velocity $\tilde V=\ip{\partial_\tau\Gt|_x}{\tilde\nu}$; the transport $\varpi=\int_0^x\kt\tilde V$; and $Q:=\int_0^x\kt'^2$.
\emph{Residual.} $r:=\tilde V+\kt_{xx}-c\tilde p-\tilde b\tilde n$, with $S_1=\int_0^xr$ and $S_2=\int_0^x(x-x')r\,dx'$. Exactly, $r+iA=e^{i\theta}\int_0^xe^{-i\theta}\rho_\theta$, with $A'=\kt r$; this is T6's identity.
\emph{Linearisation} (signs confirmed symbolically in T1 \fn{signs.py}):
\[
\cL e=-e''''+b_2e''+b_1e'+b_0e,\quad b_2=-\kt^2,\quad b_1=-5\kt\kt'-c\tilde q-\tilde b\tilde t_1+\varpi=-cy+\beta,\quad b_0=c+c_0,
\]
with $\beta:=Q-cX_0-6\kt\kt'+A$ and $c_0:=-3\kt'^2-4\kt\kt''$. The terms $-c\tilde q$ and $-\tilde b\tilde t_1$ cancel exactly in $b_1$.
\emph{Constraints and directions.} $\varphi_1=\chi p_L$ and $\varphi_2=\chi n_L$ are fixed functions of a late base (\fn{t7/phi_store.json}, exactly $C^4$); $\varphi_3=\chi\hat\psi_3$. The directions are $g=(\tilde p,\tilde n,\tilde V/n)$, $n=\norm{\chi\tilde V}_\rho$, each normalised in $L^2_\rho$.
\emph{Energies of the error $e$.} Let $u=\chi e$ and $v=\zeta_\infty e'$, where $\zeta_\infty$ is the step-3 ramp on $[10.25,16.25]$ and equals $1$ beyond. Set $A,W_1,B,W_3,D_4=\norm{u^{(j)}}_\rho$ and $V_j=\norm{v^{(j)}}_{W_\infty}$, with $W_\infty=(1+x^2)^{-2}$ on all of $[10.25,\lt]$. Then
\[
\Lam_e=A^2+\eps B^2+\mu_e(V_0^2+\eps_IV_2^2),\qquad Y=\sqrt{\Lam_e},\qquad\eps=10^{-5},\ \eps_I=3\cdot10^{-5},\ \mu_e=0.036 .
\]
Since $W_{\rm step3}\le W_\infty$ (checked by the review: the blend has $\phi'\ge4x/(1+x^2)$ for $x\ge6.4$), the step-3 functional satisfies $\Lam^{S3}[e]\le(0.08/\mu_e)\Lam_e$. The far energies $F_j$, $D_j$ (levels $j\le3$) are those of Lemma~F (\S\ref{ss:far}).
\emph{Boxes.} $\mathbb B$ (tip, $x\in[0,34.25]$): $\abs e,\abs{e'}\le\eta_1=10^{-5}$, with no condition on $e''$ (Lemma~4a); $B\le\eta_B=10^{-4}$; $\abs\vartheta\le10^{-3}$. $\mathbb B'$ (far; the boxes of N3$'$, certified at $\Yb_{\rm box}=1.7\cdot10^{-12}$): for $y\ge28$, $\abs e\le1.7\,E(x)$, $\abs{e'}\le b_0=1.5\cdot10^{-3}$, $x\abs{e''}\le b_1=3\cdot10^{-4}$, $x^2\abs{e'''}\le b_2=10^{-2}$; for $y<28$, $\abs e,\abs{e'}\le\eta_1$ and $x\abs{e''},x^2\abs{e'''}\le\eta_F=10^{-2}$. Here $E(x):=c_e\Yb+\sqrt{F_0}\,(y^3/84)^{1/2}$, with $c_e=1.6031\cdot10^4$ from Lemma~4b.

\subsection{The shadowing theorem}\label{ss:thmP2}
\begin{named}{Theorem P2 (shadowing; consolidated from T5 \S1.3)}
Assume:
\begin{itemize}[leftmargin=2em]
\item[(I)] at $\tau_a$ the true curve has a modulated decomposition over the base, with $Y(\tau_a)\le6.21\cdot10^{-13}$ and $F_j(\tau_a)\le F_j^0=(4\cdot10^{-23},3\cdot10^{-22},5\cdot10^{-19},3.4\cdot10^{-15})$ (Lemma~\ref{lem:switch});
\item[(B)] the base is exactly symmetric at both ends, with $\hat\kappa\in C^4$ and $\hat\kappa^{(5)}\in L^\infty$, enclosed jets, (SE) exact, and $n\ge2.32\cdot10^{-4}$;
\item[(R)] the residual norms of \fn{p2/t6c/residuals.json} (Table~\ref{tab:res}) hold per slab;
\item[(C)] T7's constants and certificates hold, including the nonlinear closure of Lemma~4c (\S\ref{ss:coer});
\item[(F)] T8's per-slab constants of Lemmas F, J and N3$'$ hold (\S\ref{ss:far});
\item[($\mathbb B$)] the bars stay strictly inside $\mathbb B\cup\mathbb B'$ (\S\ref{ss:boxes}).
\end{itemize}
Then the CDF solution through the P1 enclosure exists on $[t_{96},t_1]$, and there are $C^1$ functions $\lambda>0$, $c$ and $t$ of the label $\tau\in[\tau_a,27.15]$, with $t_\tau=\lambda^4(1-\vartheta)$ and $\abs\vartheta\le3.37\cdot10^{-8}$, such that
\[
\frac{\gamma(t(\tau))-c(\tau)e_1}{\lambda(\tau)}=\Gt(\cdot,\tau)+e(\cdot,\tau)\,\tilde\nu(\cdot,\tau),\qquad\ip{\chi e}{\varphi_j}_\rho=0\ (j=1,2,3),
\]
\[
\sqrt{\Lam_e(\tau)}\le\Yb=10^{-12},\qquad\sqrt{F_j(\tau)}\le A_j\Yb+B_j\ (j\le3),
\]
with $A_j,B_j$ as in Table~\ref{tab:far}; and there is a good time $\tau^*$ in slab 180 with $D_2(\tau^*)\le9.40\cdot10^{-12}$.
\end{named}
\noindent\emph{Status.} \tC: every constant is certified, and the nonlinear terms are inside the certified closure (review A2 closed by Lemma~4c, A3 by Lemma~4b). The analytic steps of the proof below are \tD\ in T5 and are now written out in \fn{LEMMA_F.md} and \fn{T7_REPORT.md} \S4a--4c; continuation uses the literature fact (P1) of step~3 \tA.

\subsection{Proof}\label{ss:proofP2}
\emph{Step 1: representation.} For the true solution put $\Gamma=(\gamma-ce_1)/\lambda$ as a function of the label $\tau$, with $t_\tau=\lambda^4(1-\vartheta)$. Then $\ip{\Gamma_\tau}\nu=-(1-\vartheta)\kappa_{ss}+(\frac14-a)\ip\Gamma\nu+b\ip{e_1}\nu$. Subtract the base identity and write $\Gamma$ as a normal graph over $\Gt(\tau)$. The $\vartheta$-direction becomes $\tilde V$, and
\[
e_\tau=\cL e-\sum_{i=1}^3x_ig_i-(1-\vartheta)r+\cN(e)+\cQ(x;e),\quad x_1=(a-\tilde a)-\vartheta c,\ x_2=-(b-\tilde b)-\vartheta\tilde b,\ x_3=n\vartheta .
\]
Here $\cN$ is the nonlinear remainder of the graph operator and $\cQ$ is bilinear in $x$ and the jets of $e$. Since $e''=\kappa-\kt+\dots$ with $\kappa$ smooth, $e\in W^{7,\infty}$ when $\hat\kappa\in C^4$ with $\hat\kappa^{(5)}\in L^\infty$, which is enough for every identity below. The inequalities hold a.e.\ in $\tau$ and are integrated.

\emph{Step 2: modulation (M).} Differentiate the constraints and integrate by parts onto the smooth $\varphi_j$, as in step 3's \fn{MOD_IBP}:
\[
\sum_ix_i(M_{ij}+E_{ij}(e))=\ip e{\ell_j}_\rho-(1-\vartheta)\ip{S_2}{\rho^{-1}(\rho\varphi_j)''}_\rho+\ip{\chi(\cN+\cQ)}{\varphi_j}_\rho+\text{strip},
\]
with $M_{ij}=\ip{\chi g_i}{\varphi_j}_\rho$, $\ell_j=\rho^{-1}\cL^\dagger(\rho\varphi_j)$, $\cL^\dagger\psi=-\psi''''+(b_2\psi)''-(b_1\psi)'+b_0\psi$. Hence $\abs x\le C_xA+x_r+x_{\rm NL}$, with $x_r=\bar m\abs{c_\varphi}\cR_2\le2.0\cdot10^{-16}$ and $x_{\rm NL}\le1.95\cdot10^{-15}$ (Lemma~4c). Existence and continuity of $(\lambda,c,t)$ follow from the implicit function theorem with matrix $M+E$, which is valid while $\bar m\norm{E}<\frac12$ (at the box $\bar m\norm{\hat E}\le1.04\cdot10^{-4}$).

\emph{Step 3: clock.} By (M), $\abs\vartheta=\abs{x_3}/n\le\bar x/n_{\min}\le3.37\cdot10^{-8}$ at $\Yb=10^{-12}$ (\fn{t1clock.json}), far inside $10^{-3}$. The third direction is an exact symmetry, the shift of $\tau$ along the base orbit, so $Y$ is not forced by the accumulated clock shift, and the $2\times2$ system of T5 \S2.3 is triangular.

\emph{Step 4: energy inequalities.} The four inequalities of T5 \S2.4 are, with $\Phi_1\ge\sup\abs{\rho'/\rho}$ and $\Phi_2\ge\sup\abs{\rho''/\rho}$ on $[0,34.25]$,
\begin{align*}
\tfrac12\tfrac d{d\tau}A^2&\le-c_IA^2-\theta_0B^2+(s_A+s_B)V_0^2+\cR_2\bigl(B+2\Phi_1W_1+(\Phi_2+\bar mc_\varphi)A\bigr)+N_0,\tag{E0}\\
\tfrac\eps2\tfrac d{d\tau}B^2&\le\eps\bigl[-\tfrac12D_4^2+G_{22}B^2+G_{12}W_1B+(G_{02}+G_x)AB\\&\qquad+\cR_s(D_4+2\Phi_1W_3+\Phi_2B)+\text{strip}_2\bigr]+N_2,\tag{E2}\\
\tfrac12\tfrac d{d\tau}V_0^2&\le-\mu_1V_0^2-\tfrac12V_2^2+a_fW_1^2+b_fB^2+\cR_1(V_2+2\Psi_1V_1+\Psi_2V_0)+\abs x\gamma_0V_0+N_{I0},\tag{I0}\\
\tfrac{\eps_I}2\tfrac d{d\tau}V_2^2&\le\eps_I\bigl[-\tfrac12V_4^2-\mu_2V_2^2\\&\qquad+(c_3D_4+c_2W_3+c_1B+c_0W_1+\cR^I_s+\abs x\gamma_2)(V_4+2\Psi_1V_3+2c_{1W}V_2)\bigr]+N_{I2}.\tag{I2}
\end{align*}
In (E0) the residual is paired weakly: $\ip{\chi r}u_\rho=\int S_2(\rho\chi u)''$, so only $\kt$ is needed. In (I0), $\ip{\zeta r_x}v_W=\int S_1(W\zeta v)''$. At the far end, $b_1(\lt)=\varpi(\lt)=-\lt_\tau$ there, so the transport boundary term cancels the moving-domain term. The nonlinear terms $N_0,N_2,N_{I0},N_{I2}$ and the nonlinear part of (M) are bounded by Lemma~4c by 55 explicit terms per interval, and the certified quadratic form in nonlinear mode (\fn{closure_cert.py --nl}) gives
\[
\mathfrak Q_c(z)-\mathrm{NL}(z)\ge\varrho_c\Lam_e+\tfrac{\theta_0}4B^2+\tfrac\eps8D_4^2+\tfrac{\mu_e}8V_2^2+\tfrac{\mu_e\eps_I}8V_4^2=:\varrho_c\Lam_e+\mathfrak D,\qquad z=(A,B,D_4,V_0,V_2,V_4),
\]
with $\min\varrho_c=0.04711$ on the 25 intervals of slabs 2--180.

\emph{Step 5: comparison.} By Lemma~\ref{lem:Y} (\S\ref{ss:Y}), $Y\le\Yb=10^{-12}$ on $[\tau_a,27.15]$, as long as $\mathbb B\cup\mathbb B'$ holds.

\emph{Step 6: far field.} Lemmas F, B, J and N3$'$ and Corollary~C (\S\ref{ss:far}) give $\sqrt{F_j}\le A_j\Yb+B_j$ and the good time $\tau^*$ with $D_2(\tau^*)\le9.40\cdot10^{-12}$.

\emph{Step 7: closing.} Run the continuity argument on the largest interval on which $\mathbb B\cup\mathbb B'$ holds. Steps 2--6 and the pointwise bounds of \S\ref{ss:boxes} give $\mathbb B\cup\mathbb B'$ with strict inequalities, so the interval is all of $[\tau_a,27.15]$. N3$'$ is certified at $\Yb_{\rm box}=1.7\cdot10^{-12}$, above the $\Yb=10^{-12}$ in use. The bounds control the curvature and its derivatives, so the solution continues: this is the continuation fact (P1) of step 3, from the literature \tA. No tail lemma and no barrier march are needed: the system is linear with small coefficients, and the horizon is finite ($\Delta\tau\le27.15$).\hfill$\square$

\subsection{The modulation constants}\label{ss:mod}
From T7 (C3), per slab (\fn{c2r3}): $\bar m\le5.40$, $\bar\ell\le9.57$, $C_x\le7.63$, $C_{x,\rm full}\le7.82$, and $c_\varphi=(1.895,1.534,0.318)$ for normalised $\varphi_j$, with Euclidean norm $\abs{c_\varphi}\le2.459$ over all certificates. On the direct residuals, $n\ge2.32\cdot10^{-4}$ (slab 180; \fn{constants_b1.json}). At $\tau_c$, from $k=27$ to $1.6\cdot10^4$: $\bar m$ $5.06\to5.39$, $\bar\ell$ $9.56\to9.41$, $n_V=\norm{\chi V_s}$ from $0.167$ down to $3.4\cdot10^{-4}$. The residual corrections to the directions ($\tilde V=V_s+r$ in $g_3$, $b_1\mathrel{+}=\int\kt r$) cost at most $9.7\cdot10^{-9}$ in $c_I$. The modulation norms after \fn{fargref} are $C_{\gamma0}\le0.2069$, $C_{\gamma2}\le0.0160$, $G_x\le14.06$. Four $t_c$-balls per slab remove the ``sawtooth'' in $\sqrt{\hat\Gamma_0}$, an artefact of the $t_c$-hull of the moving outer knots.

\emph{T8's hypothesis on $x$ (review B6).} T8 (\fn{camp.py}, \fn{n3far.py}) uses $\abs x\le\bar x=C_x\Yb+x_r$ with $C_x\le7.824$ and $x_r\le2.0\cdot10^{-16}$. The true $\abs x$ is $\le C_xY_{\rm cert}+x_r+x_{\rm NL}$ with $Y_{\rm cert}=6.21\cdot10^{-13}$ (\S\ref{ss:Y}), so the hypothesis holds with the slack $C_x(\Yb-Y_{\rm cert})\ge2.44\cdot10^{-12}$; $x_{\rm NL}\le1.95\cdot10^{-15}$ uses at most $8.0\cdot10^{-4}$ of it.

\subsection{Residuals and the (SE)-corrected base (T6; review B1)}\label{ss:res}
\emph{Verified by} \fn{p2/t6c/slabc.py} (T6's engine run on the corrected base); \emph{certificate} \fn{p2/t6c/residuals.json} and \fn{p2/t6c/jets/}, the one T6 data set of the final tail (\fn{FINAL_manifest.json}, key \fn{data}); method \fn{t6/T6_REPORT.md}; lemma \fn{LEMMA_F.md} \S6.1 and \fn{T8_REPORT.md} \S4.

\emph{Method \tC, no sampling.} On slab $i$ each unknown is the degree-8 polynomial in $\tau$ through the stored stage balls, with Radau nodes taken as exact rationals ($c_8=1$, so slabs join exactly). The outer knots $X_T(1+C\lt^3)^{1/3}$ are Taylor models in $\tau$. All objects are polynomials in $(s,\tau)$ with ball coefficients, and every truncation folds into the radii. $\rho_\kappa=\kt_\tau|_x-\Phi_x$ and $\rho_\theta=\int\rho_\kappa$ are exact. Then $r+iA=e^{i\theta}\int_0^xe^{-i\theta}\rho_\theta$, $(r,A)'=(\rho_\theta-\kt A,\kt r)$, and $S_1,S_2$ by exact integration. Per cell (length $\le0.25$ on $[0,40]$, split at the weight breakpoints; 2 per outer element) the exact polynomial $P_c(\tau)=\int_cf^2$ is formed. Weights are bounded on cell $\times$ quarter-slab, and $\sup_\tau\sum\bar w_cP_c$ uses a 32-piece Taylor-shift subdivision. Between nodes the far-field $\rho_\kappa$ is $\sim10^3$ times its nodal value, so the exact $\tau$-polynomial is necessary. Stored hulls and ranges are round-to-nearest; every reader widens them outward by 2 ulps (review B9).

\emph{The (SE) correction.} Lemma~F needs the symmetric-end hypothesis (SE), $\cB^S_j\le0$ (\S\ref{ss:far}); it holds with $\cB^S_j=0$ when the base is symmetric about the axis at its far end. The stored base meets this only at slab ends (\S\ref{s:base}). The correction is a \emph{comoving} curvature bump:
\[
\hat\kappa=\kt+\delta\kappa,\quad\delta\kappa=\frac{a_1}{\lt}b_1(u)+\frac{a_2}{\lt^2}b_2(u),\quad u=\frac x{\lt(\tau)}\in[0.6,0.8],
\]
\[
b_1=B,\quad b_2=B',\quad B(u)=\frac{\psi((u-0.6)/0.2)}{0.2\,I_\psi},
\]
with $\psi=t^6(1-t)^6$ ($C^5$), $\int B=1$, $\int b_2=0$ and $\int ub_2=-1$. On $x<0.6\lt$, i.e.\ on the whole tip region, nothing changes, so the tip residual norms and the T7 tip certificates are untouched.

\begin{lemma}[the corrected base and its residual \tC; \fn{LEMMA_F.md} \S6.1]\label{lem:SE}
On every slab $2..180$: $a_1=-\delta$ (the turning defect) is exact; $a_2$, the root of $\operatorname{Im}\hat W(\ell)=0$, is unique in a ball of radius $\le7.6\cdot10^{-26}$ by Banach's fixed-point theorem (contraction $\le2\cdot10^{-14}$), and $|a_2|\le1.6\cdot10^{-14}$. The corrected far end encloses the axis: $|\operatorname{Im}\hat W(\ell)|\le2.2\cdot10^{-20}$ and $|\hat\theta(\ell)+3\pi/2|\le1.9\cdot10^{-25}$, so (SE) holds. T6's engine, run on the corrected base itself (time degree 20, $e^{i\theta}$ from $\theta$ itself, window elements cut at $u=0.6$ and $0.8$, slab 11 also cut in time where a knot crosses $u=0.6$), gives the residual norms of Table~\ref{tab:res}. With the bump set to zero, the run reproduces T6.
\end{lemma}
\noindent\emph{The retired majorants.} The first correction run (\fn{se_corr.py}) added majorants to T6's uncorrected norms; they had no derivation (review B1), and its $\cR_0$ increment was in fact invalid: it left out that $S_1(\ell)$ moves, which shifts $\partial^{-1}r$ over the whole far field. The direct computation replaces it, and no final stage reads \fn{se_corr.json} or \fn{p2/t6/*} (\fn{FINAL_manifest.json}, \fn{note}). The direct level-0 hats are up to $\times3$ larger on some slabs, but their maxima are below the \fn{se_corr} values ($\cR^F_0$: $3.5\cdot10^{-9}$ against $8.7\cdot10^{-9}$; $\cR^F_3$: $4.9\cdot10^{-9}$ against $1.7\cdot10^{-7}$).

\begin{table}[htbp]\centering\small
\caption{Residual norms of the (SE)-corrected base, maxima over slabs 2--180 (\fn{p2/t6c/residuals.json}, regenerated summary). The P2 scripts read the per-slab values.}\label{tab:res}
\begin{tabular}{@{}lllll@{}}\toprule
norm & T5 target & max (slab) & role\\\midrule
$\cR_2=\norm{S_2}_{L^2_\rho(0,34.25)}$ & $2\cdot10^{-14}$ & $2.1\cdot10^{-17}$ (2) & (E0), (M)\\
$\cR_s=\norm r_{L^2_\rho(0,34.25)}$ & $5\cdot10^{-11}$ & $4.1\cdot10^{-14}$ (2) & (E2)\\
$\cR_1=\norm{\zeta S_1}_{W_\infty}$ & $5\cdot10^{-13}$ & $4.7\cdot10^{-15}$ (2) & (I0)\\
$\cR^I_s=\norm{\zeta r_x}_{W_\infty}$ & $10^{-10}$ & $2.8\cdot10^{-18}$ (2) & (I2)\\
$\cR^F_0$ (weights \fn{opt4n}) & $10^{-7}$ & $3.5\cdot10^{-9}$ (174) & Lemma F\\
$\cR^F_1$ & $10^{-7}$ & $4.5\cdot10^{-11}$ (153) & Lemma F\\
$\cR^F_2$ & $10^{-7}$ & $2.5\cdot10^{-10}$ (153) & Lemma F\\
$\cR^F_3$ & $10^{-7}$ & $4.9\cdot10^{-9}$ (153) & Lemma F\\
$\cR^F_4$ & $10^{-5}$ & $1.6\cdot10^{-7}$ (153) & not used (level 4 dropped)\\\bottomrule
\end{tabular}\end{table}
\noindent Far-end scalars of the corrected base: $\sup\abs{A(\lt)}\le1.92\cdot10^{-15}$, $\sup\abs{\eps_\ell}\le1.92\cdot10^{-15}$, and the transport defect $\abs D\le1.21\cdot10^{-12}$ (Lemma~B). The T6 corrections to T7 ($\delta\ip{\cL u}u_\rho=-\frac12\int(\rho A)'u^2$, with $\sup_{[0,34.25]}\abs{A'}\le1.2\cdot10^{-10}$) are $\ge10^7$ below T7's margins.

\subsection{The comparison system (Y)}\label{ss:Y}
\emph{Verified by} \fn{p2/ycomp.py} (\fn{T6RES=p2/t6c/residuals.json}, \fn{SLAB0=2}, argument $Y_0=6.21\cdot10^{-13}$); \emph{certificate} \fn{p2/out/ycomp_final.json}. All residual norms are multiplied by $1.001\ge\abs{1-\vartheta}$. Interpolations: $W_1\le(\frac12+\sqrt{c_1})A+\frac12B$, $W_3\le(\frac12+\sqrt{c_1})B+\frac12D_4$, $V_1\le(\frac12+\sqrt{c_{1W}})V_0+\frac12V_2$, $V_3\le(\frac12+\sqrt{c_{1W}})V_2+\frac12V_4$, with $c_1=0.5804$ and $c_{1W}\le0.09435$. Weak-pairing constants: $(\rho\chi u)''/\rho=\chi u''+2(\chi'+\chi\rho'/\rho)u'+(\chi''+2\chi'\rho'/\rho+\chi\rho''/\rho)u$, so $c_{W1}=2(\chi_1+\Phi_1)$ and $c_A=\chi_2+2\chi_1\Phi_1+\Phi_2$, with $\chi_1=0.73$, $\chi_2=0.84$, $\Phi_1=1.083$, $\Phi_2=1.3$. The review rechecked $\sup\abs{\chi'}=0.729$, $\sup\abs{\chi''}=0.835$, $\sup\abs{\rho'/\rho}=1.082$ and $\sup\abs{\rho''/\rho}\approx1.16$. With $x_r=\bar m\abs{c_\varphi}\cR_2$:
\begin{align*}
r_B&=(1+\tfrac12c_{W1})\cR_2+\eps\cR_s(2\Phi_1k_A+\Phi_2)+\eps g_{2E}x_r,\quad r_{D4}=\eps\cR_s(1+\Phi_1),\\ r_{I2}&=\mu_e\eps_I(\cR^I_s+g_2x_r),\\
r_{V2}&=\mu_e\cR_1(1+\Psi_1)+r_{I2}(2\Psi_1k_V+2c_{1W}),\quad r_{V4}=r_{I2}(1+\Psi_1),\\
r_A&=(c_A+c_{W1}k_A)\cR_2+\sqrt3\,x_r,\\ r_{V0}&=\mu_e\bigl(\cR_1(2\Psi_1k_V+\Psi_2)+g_0x_r\bigr),
\end{align*}
with $k_A=\frac12+\sqrt{c_1}$, $k_V=\frac12+\sqrt{c_{1W}}$, $\Psi_1\le0.38851$, $\Psi_2=2c_{1W}$, and $g_0,g_2,g_{2E}$ the per-slab modulation norms (\fn{farg}). Half of the reserve $\mathfrak D$ is kept for Lemma~J ($\kappa_{\mathfrak D}=\frac12$). The other half is shared out as $\theta_0/8$ (B), $\eps(\frac1{16}-\mathrm{D4CUB})$ ($D_4$), $\mu_e/16$ ($V_2$) and $\mu_e\eps_I/16$ ($V_4$). Young's inequality on each pair gives
\[
\tfrac12\Lam_e'\le-\varrho_m\Lam_e-\kappa_{\mathfrak D}\mathfrak D+a_cY+\Xi'^2,\qquad\varrho_m=\varrho_c-\delta_{T6}-\mathrm{DN},
\]
\[
a_c=r_A+\frac{r_{V0}}{\sqrt{\mu_e}},\qquad\Xi'^2=\sum_{k\in\{B,D_4,V_2,V_4\}}\frac{r_k^2}{4\,\mathrm{share}_k}.
\]
Lemma~4c puts every nonlinear term into $\varrho_c$, so $\mathrm{DN}=\mathrm{D4CUB}=0$ would suffice. \fn{ycomp.py} still subtracts $\mathrm{DN}=10^{-6}$ and reserves $\mathrm{D4CUB}=3\cdot10^{-9}$; this only lowers the rate and the share, and is conservative.
\begin{lemma}[constant majorant \tC]\label{lem:Y}
Let $\mathrm{lev}_m:=a_c/\varrho_m+\Xi'_m/\sqrt{\varrho_m}$ on slab $m$. If $Y(\tau_a)\le Y_c$ and $\mathrm{lev}_m\le Y_c$ for all slabs, then $Y\le Y_c$ on $[\tau_a,27.15]$.
\end{lemma}
\begin{proof}
With $Y_c=\mathrm{lev}+\delta$, $\delta\ge0$, write $\varrho Y_c^2-a_cY_c-\Xi'^2=Y_c(\varrho Y_c-a_c)-\Xi'^2\ge\Xi'(\sqrt\varrho Y_c-\Xi')\ge0$. So at any time with $Y=Y_c$ we have $(Y^2)'\le0$. Apply this with $Y_c+\delta'$, $\delta'\downarrow0$, as a strict first-touch argument.
\end{proof}
\noindent\emph{Result.} $\max_m\mathrm{lev}_m\le3.22\cdot10^{-14}$, at slab 2 (it was $1.33\cdot10^{-13}$ on the \fn{se_corr} residuals). With $Y(\tau_a)\le6.21\cdot10^{-13}$, the script sets $Y_c=\max(Y_0,\mathrm{lev},10^{-12})=10^{-12}=:\Yb$ and asserts $\Yb\ge Y_0$. So $Y\le Y_{\rm cert}:=6.21\cdot10^{-13}$ throughout, and $\Yb=10^{-12}$ is the majorant used by T8 and T9. Review B9: the output rounding \fn{up()} is now outward.

\subsection{Coercivity: the T7 certificates and Lemmas 4a--4c}\label{ss:coer}
\emph{Verified by} \fn{t7/run_all.py} (units \fn{cert7}, \fn{c2r3}, \fn{far7}, \fn{nl7}, \fn{nlterms7.py}, \fn{closure_cert.py}, \fn{merge.py}); \emph{certificate} \fn{t7/constants.json}, republished on Nov~12 with \fn{se7} on \fn{p2/t6c} (SHA-256 in \fn{FINAL_manifest.json}, key \fn{t7}); report \fn{t7/T7_REPORT.md} \S1--5, 9, 9a.

\paragraph{(C1)+(C2) \tC.} The constrained form is
\[
\cB_\tau(u):=\ip{\cL u}u_\rho-\sum_{i,j}\ip{\chi g_i}u_\rho(M^{-T})_{ij}\ip u{\ell_j}_\rho .
\]
For each of 26 intervals $I$ covering $\tau\in[0,27.15]$, every $\tau\in I$, $\theta_I=0.025$, and every even $u\in H^4_\rho$ with $\operatorname{supp}u\subset[0,34.25]$ and $\ip u{\varphi_j}_\rho=0$, we have $\cB_\tau(u)\le-c_I\norm u_\rho^2-\theta_I\norm{u''}_\rho^2$. One certificate runs at $\tau_c$ (Hermite block, Haynsworth inertia, tail and leak). It certifies $Q_c-\sum_e(\mu^0_e\norm w_e^2+\mu^1_e\norm{w'}_e^2)-c_R\norm w^2\ge0$, which covers all of $I$. The local coefficients $C_{22},C_{21},C_{11}$ do not depend on $\tau$. After integration by parts the change of the local form is $\int[-2\delta C_{20}w'^2+(\delta C_{20}''-\delta C_{10}'+\delta C_{00})w^2]$, and $\mu^{0,1}_e$ are its sups over each element and $I$, with 2 $t_c$-balls per slab. For the rank-3 part, $c_R$ is reserved as $1.6\,n_{\rm side}(F_ce^{0.3H}L+10^{-3})$, where $F_c$ is the exact operator norm of $d\Pi/d\tau$ on the constraint space. \fn{c2r3.py} certifies $V=\int\norm{d\Pi/d\tau}_Zd\tau$ per slab, and $\sum V\le0.66\,c_R$ on every side.
\begin{center}\small\begin{tabular}{@{}llllll@{}}\toprule
ids & slabs/interval & $k_c$ & $c_I$ & $c(\tau_c)$ & $c_R$\\\midrule
0--5 & 2 & 27--40 & 0.0629--0.0660 & 0.0713--0.0730 & 0.0052--0.0064\\
6--15 & 4 & 45--174 & 0.0599--0.0654 & 0.0700--0.0724 & 0.0050--0.0101\\
16--22 & 10 & 226--2151 & 0.0576--0.0614 & 0.0694--0.0700 & 0.0086--0.0118\\
23--25 & 20/20/19 & $3.8\cdot10^3$--$1.6\cdot10^4$ & 0.0528--0.0530 & 0.0700 & 0.0170--0.0172\\\bottomrule\end{tabular}\end{center}
So $\min c_I=0.05284$ (id 23, $k_c=3776$). The tip region lies in $x<0.6\lt$, where the corrected base equals the polished one, so these certificates are unchanged by the (SE) correction and by B1.

\paragraph{(C4) and the closure \tC.} The $E_I$ Gårding function on the whole far curve has raw infimum $0.1359$ ($k=26.3$, $x=10.25$; Hardy $\le1.83\cdot10^{-4}$). With the split at $\mu_f=0.15$ (below the floor only on $x\le10.52$, where $\chi=1$; $a_\mu\le2.5\cdot10^{-13}$ is added to $a_f$), $\mu_1=0.15$ on all 181 slabs and $\mu_2\ge0.643$. T8's maximal far rate is $\lambda_{\max}=0.2520$ (on \fn{constants_b1.json}: $0.25195$), so $2\mu_1/\lambda_{\max}-1=0.190$. The E2 level has $G_{22},G_{12},G_{02}\le15.63,19.37,12.74$; the strip has $s_B=3.47\cdot10^{-5}$, $s_A=2.18\cdot10^{-4}$; the ramp has $a_f\le0.0696$, $b_f\le0.0394$, $c^I_{0..3}\le(0.4251,0.4486,0.5390,0.4482)$, $\Psi_1\le0.38851$, $c_{1W}\le0.09435$. \fn{closure_cert.py} at $\mu_e=0.036$ (AM--GM relaxation and \fn{arb} $LDL^T$ on $6\times6$) certifies all 26 intervals in linear mode ($\min\varrho_c=0.0471$), and all 25 intervals of slabs 2--180 in nonlinear mode ($\min\varrho_c^{\rm NL}=0.04711$; Lemma~4c). The bound $\abs{e(x)}\le K_e\Yb$ holds for $x\le x_a=14+X_0$, with $K_e\le890.6$.

\begin{named}{Lemma 4a (tip nonlinearity without a pointwise bound on $e''$; \fn{T7_REPORT.md} \S4a)}
Let $\Omega=[0,34.25]$ and $w=\abs e+\abs{e'}$. Write the graph operator as $F=A_4e''''+A_3e'''+A_2$ and $N=F(e)-F(0)-DF(0)e$.
(i) \emph{Structure} (checked symbolically, \fn{nl7.derive}, \fn{poly_check}). $\mathsf a=1-\kt e$ and $\mathsf J=(\mathsf a^2+e'^2)^{1/2}$ contain no $e''$, so $F$ is a polynomial in $e''$: $A_4=A_4(e,e')$, $A_3=A_{30}+A_{31}e''$ with $A_{31}=10e'/\mathsf J^6$, and $A_2=A_{20}+A_{21}e''+A_{22}e''^2+A_{23}e''^3$ with $A_{22}=-3\kt$ at $e=e'=0$ and $A_{23}=3(\mathsf a^2-5e'^2)/\mathsf J^8$.
(ii) If $\abs e,\abs{e'}\le\eta_1=10^{-5}$ on $\Omega$, then for all $e'',e''',e''''$
\[
\abs N\le n_4w\abs{e''''}+n_3^*(w+\abs{e''})\abs{e'''}+n_2^*(w+\abs{e''})^2+m_{23}\abs{e''}^3,
\]
with $n_4=\sup\abs{\nabla A_4}$, $n_3^*=\max(\sup\abs{\nabla A_{30}},\sup\abs{A_{31}})$, $n_2^*=\max(\frac12\sup\abs{\nabla^2A_{20}},\frac12\sup\abs{\nabla A_{21}},\sup\abs{A_{22}})$, $m_{23}=\sup\abs{A_{23}}$. The gradients are taken in $(e,e')$, and the sups are \fn{arb} bounds over the cells of $\Omega$, the slab and the $(e,e')$-box. Certified values: $n_4\le6.17$, $n_3^*\le9.26$, $n_2^*\le16.9$, $m_{23}\le3.0003$, all maximal at $k\approx26$.
(iii) $S:=\sup_\Omega\abs{e''}$ satisfies $S^2\le S_W^2+2K_{\rm reg}BW_3$, with $S_W^2:=(V_1^2/18+2V_1V_2)/W(34.25)$, $K_{\rm reg}=1/\rho(16.25)\le3.021\cdot10^4$ and $W(34.25)\ge7.25\cdot10^{-7}$.
\end{named}
\begin{proof}
(ii) $N=[A_4-A_4(0)]e''''+[A_{30}-A_{30}(0)+A_{31}e'']e'''+[A_{20}-A_{20}(0)-\nabla A_{20}(0)\cdot(e,e')]+[A_{21}-A_{21}(0)]e''+A_{22}e''^2+A_{23}e''^3$. Apply the mean-value and Taylor theorems in $(e,e')$ to each bracket.
(iii) On $[16.25,34.25]$, $\zeta=1$ and $e''=v'$. The Sobolev bound $f^2\le L^{-1}\int f^2+2\int\abs{ff'}$ ($L=18$), with $W\ge W(34.25)$, gives $\abs{e''}\le S_W$. On $[0,16.25]$, $\chi=1$ and $\rho\ge1/K_{\rm reg}$, so $e''(x)^2\le e''(16.25)^2+2\int_x^{16.25}\abs{e''e'''}\le S_W^2+2K_{\rm reg}BW_3$.
\end{proof}
\noindent\emph{Use.} Lemma~4c uses (ii) and (iii) for the tip part of $\cN$. The earlier separate Young split of the cubic term (\fn{merge.tip_lemma}, slice $\mathrm{D4CUB}$) is superseded by Lemma~4c.

\subsubsection*{Lemma 4b: the tip box (review A3)}\label{ss:lem4b}
\begin{named}{Lemma 4b (pointwise tip box from $\Lam_e$; \fn{T7_REPORT.md} \S4b)}
Let $\Omega=[0,X_E]$, $X_E=34.249$ (the end of $\operatorname{supp}\chi_{32}$; $\chi_{32}=1$ up to $X_1=31.249$). At every time
\[
\abs{e(x)}\le c_e(x)\Lam_e^{1/2},\qquad\abs{e'(x)}\le c_{e1}(x)\Lam_e^{1/2}\qquad(x\in\Omega),
\]
\[
\sup_\Omega c_e\le1.6031\cdot10^4,\qquad\sup_\Omega c_{e1}\le1.6352\cdot10^4.
\]
On $[16.25,\lt]$, $\abs{e(x)}\le c_e(x)\Lam_e^{1/2}$ with $c_e(x)^2=K_e(16.25)^2+H(x)^2/\mu_e$ and $H(x)^2=\int_{16.25}^x(1+s^2)^2ds$.
\end{named}
\begin{proof}[Proof (\fn{T7_REPORT.md} \S4b; constants by \fn{nlterms7.lemma4b}, published in \fn{meta.tip_box})]
(i) For $x\le X_1$, with $u=e$: $\rho u^2(x)=\int_x^{X_E}(\phi'\rho u^2-2\rho uu')\le\Phi_1A^2+2AW_1$, with $\Phi_1=\sup_\Omega\phi'=1.0822$ and $W_1^2\le AB+c_1A^2\le s^2\Lam$, $s^2=1/(2\sqrt\eps)+c_1$; so $c_e(x)^2=(\Phi_1+2s)/\rho(x)$, e.g.\ $K_e(16.25)=890.95$.
(ii) For $x\le X_1$: $\rho u'^2(x)\le\Phi_1W_1^2+2W_1B\le2C_A\Lam$, with $C_A=(m_{31}\eps^{-3/2}+c_1/(4\eps))^{1/2}+(\Phi_1/2)(1/(2\sqrt\eps)+c_1)=3292.8$, $m_{31}=3\sqrt3/16$.
(iii) For $x_0\in[16.25,X_E]$, where $\zeta=1$ and $v=e'$: let $f=\psi_c(1+x^2)^{-1}v$, with $\psi_c=1$ on $[10.25,x_0]$, $\cos^2(\pi(x-x_0)/(2\ell_c))$ on $[x_0,x_0+\ell_c]$, $0$ beyond, $\ell_c=100$. Its support lies in $[10.25,\lt/2]$ ($\min\lt=411.8\ge268.5$), where $W=(1+x^2)^{-2}$ exactly. The sharp line inequality $\sup f^2\le K_{\rm line}(\norm f^2+\eps_I\norm{f''}^2)$, $K_{\rm line}=\eps_I^{-1/4}/(2\sqrt2)=4.7773$, with $\norm{f''}\le aV_2+bV_0$ ($a=1.2090$, $b=0.3998$, using $V_1\le k_{V1}V_0+V_2/2$), gives $\abs{e'(x_0)}\le(1+x_0^2)(K_{\rm line}c_\omega/\mu_e)^{1/2}\Lam^{1/2}$, $c_\omega=1.4619$. On $[16.25,X_1]$, $c_{e1}$ is the smaller of (ii) and (iii).
(iv) For $x\ge16.25$: $\abs{e(x)}\le\abs{e(16.25)}+H(x)V_0$, using $W\ge(1+x^2)^{-2}$, which also holds on the flattened part ($\Phi'\in[0,1]$ gives $X\le x$).
\end{proof}
\noindent\emph{Consequences.} $\eta_1:=c_{e1}\Yb\le10^{-5}$ (the tip box of T5 and T8) holds for every $\Yb\le6.11\cdot10^{-10}$; this is asserted. At $\Yb_{\rm box}=1.7\cdot10^{-12}$, $\eta_1=2.78\cdot10^{-8}$; at $\Yb=10^{-12}$, $\abs e,\abs{e'}\le1.61\cdot10^{-8},1.64\cdot10^{-8}$ on $[0,34.25]$. The old constant $2\cdot10^3\Yb/\sqrt{\mu_e}=1.054\cdot10^4\Yb$ was the value of (iii) at $x=30.25$ and was too small on $[30.25,34.25]$. The lemma's constants are passed to \fn{merge.tip_lemma}, to N3$'$ (through the $10^{-5}$ box), to \fn{clam.py}'s tail term, to \fn{n3boxes.py}, and to the position error of the Lemma~O scripts.

\subsubsection*{Lemma 4c: the nonlinear terms of (Y) (review A2)}\label{ss:lem4c}
\begin{named}{Lemma 4c (nonlinear closure; \fn{T7_REPORT.md} \S4c, \S9a)}
Assume the bootstrap box $\mathbb B_Y$: $\Lam_e\le\Yb_{\rm box}^2$ with $\Yb_{\rm box}=1.7\cdot10^{-12}$; Lemma~4b's tip box; T8's far box $\mathbb B'$ on $x\ge30.249$; and, per interval, $\abs{\hat x}\le\bar x=C_x\Yb_{\rm box}+x_r+x_{\rm NL}$ and $\abs\vartheta\le\bar x/n_V$. Let $N_0=\ip{\chi(\cN+\cQ)}u_\rho$, $N_2=\eps\ip{\chi(\cN+\cQ)}{D_\rho u}_\rho$, $N_{I0}=\int W\zeta^2e'\partial_x(\cN+\cQ)$, $N_{I2}=\eps_I\int(Wv'')''\zeta\partial_x(\cN+\cQ)$, and let $\delta x$ be the part of (M) beyond its linear and residual parts, entering through the modulation couplings of (E0), (E2), (I0), (I2). Then $\mathrm{NL}:=N_0+N_2+\mu_e(N_{I0}+N_{I2})+(\delta x\text{-couplings})\le\sum_kc_ka_kb_k$ over 55 explicit terms per interval (\fn{nl_terms.json}), each of dissipation degree $\le7/4<2$ (asserted), and \fn{closure_cert.py --nl} certifies
\[
\mathfrak Q_c(z)-\mathfrak D_s-\mathrm{NL}(z)\ge\varrho_c^{\rm NL}\Lam_e
\]
on all 25 intervals of slabs 2--180, with the same reserve $\mathfrak D_s$ as in linear mode and $\min\varrho_c^{\rm NL}=0.04711$ (linear mode $0.04713$). In (M), $x_{\rm NL}\le1.95\cdot10^{-15}$ and $\bar m\norm{\hat E}\le1.04\cdot10^{-4}$.
\end{named}
\noindent\emph{Estimates} (T8's exact split $\cN+\cQ=(1-\vartheta)\cN+\vartheta(e''''+\kt^2e''-c_0e+P)+ge'-(\vartheta c+x_1)e$, with $\cN=\tilde Ae''''+\tilde Be'''+\tilde C$, \fn{LEMMA_F.md} \S6.1). \emph{Tip (E0), (E2):} Lemma~4a's bound for $\cN$, with $A_{31}=10e'/\mathsf J^6\le10\eta_1/\alpha_{\min}^6$; $\abs u\le\eta_1$; $\norm{D_\rho u}\le D_4+2\Phi_1W_3+\Phi_2B$; $\rho$-norms on $\Omega$ by $N_k\le U_k+\sigma V_{k-1}$, $\sigma=5.28\cdot10^{-3}$; $S\le S_W+(2K_{\rm reg})^{1/2}(B^{3/4}D_4^{1/4}+c_1^{1/4}B)$. \emph{(I0):} one global weak identity on $[10.25,\lt]$, without boundary terms ($\zeta$ vanishes to order 4 at $10.25$, (SE) holds at $\lt$); its integrands are bounded region by region: on the far core by T8's per-range $\tilde A^+,\abs{\tilde A},\abs{\tilde A'},\abs{\tilde B}$, $\gamma$, $G$; on $[16.25,30.25]$ analytically with $\abs{e''}\le S_W$; on the ramp with $\abs{e''}\le S$; the $e$-terms by Lemma~4b(iv). \emph{(I2):} strong form, paired per range with $(Wv'')''/W$; on $\zeta=1$ the diagonal term $\int W\tilde A_{\rm eff}(v'''')^2\le\tilde A^+V_4^2$ keeps its sign; $\partial_x\cN$ on $[10.25,30.25]$ is a polynomial with 12 monomials whose coefficients are enclosed by second-order AD (\fn{nlsym/td2.py}); no $e'''$ box is used below $x=30.25$. \emph{(M):} the weak form $\int\psi(-(\mathsf J/\mathsf a)\kappa_{ss})=-\int((\psi/\mathsf a)'\mathsf J^{-1})'\kappa$ is exact (\fn{nlgen.py}) and quadratic in $e''$.

\noindent\emph{Results on the direct residuals} (\S9a; Nov~12). All 25 intervals are certified; $\min\varrho_c^{\rm NL}=0.04711$, losses $\le8\cdot10^{-5}$; the top terms are $\mu_e\tilde A^+V_2^2\le2.5\cdot10^{-4}$ with $\sup\tilde A^+\le6.95\cdot10^{-3}$; $x_{\rm NL}\le1.95\cdot10^{-15}$ against the slack $\ge2.44\cdot10^{-12}$. The reserve and its split are unchanged ($\kappa_{\mathfrak D}=\frac12$), and no dissipation slice is used, so the one-slab hand-off window does not interact with this lemma. T7's own (Y) level with $\varrho_c^{\rm NL}$ and $\mathrm{DN}=\mathrm{D4CUB}=0$ is $3.21\cdot10^{-14}$. Interval 0 (slabs 0--1) is outside P2; it keeps \fn{se_corr} widths, flagged, and no final stage reads it (\fn{cconst} and \fn{merge} take T7's worst cases over slabs $\ge2$ only, env \fn{T7SLAB0=2}).

\subsection{The far field: Lemmas F, B, J and N3\texorpdfstring{$'$}{'}}\label{ss:far}
\emph{Verified by} \fn{p2/t8se/}: \fn{lfsym.py}, \fn{cforms.py} (identity), \fn{n3sym.py}, \fn{n3qsym.py}, \fn{n3wcheck.py} (symbolic), \fn{cconst.py}, \fn{merge.py}, \fn{camp.py}, \fn{n3far.py}, \fn{n3boxes.py} (certified); \emph{certificates} \fn{constants/constants_b1.json}, \fn{n3_b1.json}, \fn{amp_A1_w180.json}; statement and proof \fn{t8/LEMMA_F.md} \S1--6. All on base v2 with the direct residuals of \fn{p2/t6c}, slabs 2--180. (\fn{T8_REPORT.md} \S1--5 and \fn{LEMMA_F.md} \S6.4 still speak of earlier data; review C.)

\paragraph{Setting.} On $y\ge y_a\ge14.001$ we have $\zeta_\infty=1$ and $v=e'$. Differentiating the error equation once gives
\[
v_\tau=-v''''-c\,y\,v'+(\beta v+b_2v'+c_0e)'+f,\qquad f=-x_1g_1'-x_2g_2'-x_3(g^s_3+g^r)'-r'+(\cN+\cQ)',
\]
with $g_3=\tilde V/n=g^s_3+g^r+r/n$, $g_3^s=(c\tilde p+\tilde b\tilde n)/n$ and $g^r=-\kt''/n$. The clock factor drops out of the residual forcing. Write $w_m:=\partial_x^mv$ and $w_{-1}:=e$. \emph{Parities.} About $x=\lt$ on the doubled curve, $\kt,b_2,c_0,e$ are even, $w_m$ has parity $(-1)^{m+1}$, and $b_1+\lt_\tau$ is odd up to the constant $D:=b_1(\lt)+\lt_\tau=A(\lt)+\eps_\ell$, $\abs D\le1.21\cdot10^{-12}$. \emph{Weights.} $T_j=28\,\eta_j(y)^2Y_\ell^{2j-2}$, with $\eta_j=S(\ln(y/y_a)/\ln(y_b/y_a))$, the $C^3$ smoothstep $S(u)=u^4(35-84u+70u^2-20u^3)$ and the ramp $[y_a,y_b]=[14,22]$ for all $j$ (design \fn{opt4n}, $\delta_P=(\frac14,\frac14,\frac14,\frac32)$). $Y_\ell=y_\ell\Phi(y/y_\ell)$ is a far-end flattening that is even about $\lt$: $\Phi(s)=s$ for $s\le\frac12$, and $\Phi'=1-\psi(2s-1)$ with $\psi(u)=u^5(\frac{21}4-\frac{49}8u+\frac{15}8u^2)$. The properties used are: (W1) $T_j\in W^{4,\infty}$; (W2) $T_j=0$ for $y\le y_a$ and $T_j\asymp(y-y_a)^8$ there; (W3) $T_j$ is even about $\lt$. \emph{Energies.} $F_j=\int_{x^*}^{\lt}T_jw_j^2$, $D_j=\int T_jw_{j+2}^2$, $N_{j+1}^2=\int T_jw_{j+1}^2$.

\paragraph{The identity $(I_j)$ \tC\ (machine-checked, \fn{lfsym.py}, Euler-operator test).} With $T=T_j$ and $\doteq$ meaning equality modulo a total $x$-derivative,
\[
2Tw_j\partial^j[v_\tau-f]\doteq-2Tw_{j+2}^2+P_jw_{j+1}^2+q_jw_j^2+\sum_{i=-1}^{j-1}p_{ji}w_jw_i+X_j,
\]
where $P_j=4T''+2\kt^2T$; $q_j=-T''''+(cy-\beta)T'+(1-2j)cT+(2j+1)\beta'T+(Tb_2)''-(j+1)(Tb_2')'+j(j+1)Tb_2''+2c_0T$;
\begin{multline*}
p_{ji}=2\tbinom{j+1}{j+1-i}T\beta^{(j+1-i)}[i\ge0,\,j+1-i\ge2]+2\tbinom{j+1}{j+2-i}Tb_2^{(j+2-i)}[1\le i\le j-1]\\
+2\tbinom{j+1}{j-i}Tc_0^{(j-i)}[i\le j-1];
\end{multline*}
and $X_j$ keeps all coefficient derivatives at $\kt^{(\le5)}$, via integrations by parts onto $Tw_jw_i$.

\begin{named}{Lemma F (far-field energies of the error; \fn{LEMMA_F.md} \S2)}
Let $j\in\{0,\dots,3\}$, let the weights satisfy (W1)--(W3), with $y_{a,i}<y_{b,j}$ for $i<j$, and assume (SE). Fix $a=0.01$, $\nu=10^{-3}$, $\delta_P\in(0,1.5]$, $\delta_s\in(0,\frac12]$. Let $n_j$ be the number of $(j+1,\cdot)$-pairs in $(I_j)$ plus one, and $\eps_j=\abs D/m_j$. Define
\begin{gather*}
\hat q_j=\partial_\tau T_j+q_j+\tfrac12P_j''+\frac{P_j^2}{4\delta_PT_j}+\Bigl(\frac{an_j}4+\frac{\eps_j}2\Bigr)(T_j'')^++\nu\abs{T_j''},\quad\Delta F_j=\frac{an_j}8+\eps_j\Bigl(\frac2{y_\ell}+\frac54\Bigr),\\
\lambda_j=-\sup_{C_j}\frac{\hat q_j}{T_j}-\Delta F_j,\qquad\theta_j=2-\delta_P-\frac{an_j}2-\delta_s-\eps_j,
\end{gather*}
$k_{jj}=\sup_{R_j}(\hat q_j+(\lambda_j+\Delta F_j)T_j)^+/W$, $c_{ji}=\sup_{C_j}\abs{p_{ji}}/\sqrt{T_jT_i}$ and $r_{ji}=\sup_{R_j}\abs{p_{ji}}/W$, with the $e$-terms and the $(j+1,0)$-pairs as in \fn{LEMMA_F.md} \S2. The $e$-terms use $\hat e(x)=K_e+\tilde I(x)/\sqrt{\mu_e}$, $\tilde I^2=\int_{x^*}^xW^{-1}$, $\hat\eps_0=\hat e(x_1)$ and $H^2=\int_{x_1}^xT_0^{-1}$. The modulation norms are $\gamma_j=(\sum_i\int T_j(\partial^{j+1}g_i^s)^2)^{1/2}$ and $\gamma^r_j=(\int T_j(\partial^{j-1}g^r)^2)^{1/2}$. The residual norms are $\cR_j=(\int T_j(\partial^{j-1}r)^2)^{1/2}$ and $\hat\cR_j^2=\frac8a\int\frac{T_j'^2}{T_j}\rho_j^2+\nu^{-1}\int\abs{T_j''}\rho_j^2$, with $\rho_j=\partial^{j-1}(r+x_3g^r)$ and $\partial^{-1}r:=-\int_x^{\lt}r$. If $\lambda_j>0$ and $\theta_j>0$, then along the linear part of the equation, for a.e.\ $\tau$,
\[
F_j'\le-\lambda_jF_j-\theta_jD_j+2\sqrt{F_j}\,g_j+h_j,\tag{F}
\]
\[
g_j=\sum_{i<j}\tfrac12c_{ji}\sqrt{F_i}+\abs x\gamma_j+\tfrac12a_j\Yb,\quad h_j=\sum_{i<j}d_{ji}F_i+\mathcal I_j+e_j\Yb^2+\delta_s^{-1}(\cR_j+\abs x\gamma^r_j)^2+\hat\cR_j^2,
\]
\[
\mathcal I_j=k_{jj}V_j^2+\sum_{i<j}(k_{ji}V_i^2+r_{ji}V_jV_i)+re_jV_j\Yb .
\]
\end{named}
\begin{proof}[Proof (\fn{LEMMA_F.md} \S3)]
\emph{3.1.} $T_j$ vanishes to order 8 at $y_a$, so nothing happens at $x^*$. At the moving end, after inserting $(I_j)$, what remains is $F_j'=\int[(\partial_\tau T_j)w_j^2+\text{density of }(I_j)]+\cB^D_j+\cB^S_j$. Here $\cB^D_j=D\,T_j(\lt)w_j(\lt)^2$ is the transport defect (Lemma~B), and $\cB^S_j$ is the sum of the fluxes at $\lt$ of the fourth-order, $b_2$-, $c_0$-, forcing and $X_j$-terms. $\cB^S_j$ contains no $b_1$, and $\cB^S_j=0$ under (SE) (closed-curve argument below).
\emph{3.3.} (S1) $\int Pw_{j+1}^2=-\int Pw_jw_{j+2}+\frac12\int P''w_j^2$, and $\abs{\int Pw_jw_{j+2}}\le\delta_PD_j+(4\delta_P)^{-1}\int(P^2/T)w_j^2$. (S2) For $0\le i<j$: $\int\abs{p}\abs{w_jw_i}\le c_{ji}\sqrt{F_jF_i}$ on $C_j$ and $\le r_{ji}V_jV_i$ on $R_j$. (S3) $N_{j+1}^2\le D_j+\frac14F_j+\frac12\int(T'')^+w_j^2$, then $\abs{\int pw_{j+1}w_i}\le\frac a2N_{j+1}^2+M_i^2/(2a)$. (S4) $\abs{\int pw_{j+2}e}\le\delta_6D_j+(4\delta_6)^{-1}\int(p^2/T)e^2$. (S5) $e=e(x^*)+\int_{x^*}^xv$ gives $\abs e\le\hat e\,\Yb$ on $R_0$ and $\abs e\le\hat\eps_0\Yb+H\sqrt{F_0}$ on $C_0$. (S6) For the residual and rough modulation, two integrations by parts give $-2\int(Tw_j)''\rho_j$, which is bounded by $\delta_sD_j+\delta_s^{-1}(\cR_j+\abs x\gamma_j^r)^2+\frac a2N_{j+1}^2+\hat\cR_j^2+\nu\int\abs{T''}w_j^2$. For the smooth modulation, Cauchy--Schwarz gives $2\abs x\sqrt{F_j}\gamma_j$.
\emph{3.4.} Collect the $w_j^2$-terms as $\int\hat q_jw_j^2+\Delta F_jF_j$: $\le-\lambda_j\int_{C_j}T_jw_j^2$ on the core, and $\le-\lambda_j\int_{R_j}T_jw_j^2+k_{jj}V_j^2$ on the ramp. The $D_j$-coefficient is $-\theta_j$.
\end{proof}

\begin{named}{Lemma B (transport defect)}
With $m_j=(\Phi(\frac12)/\Phi(1))^{2j-2}$ for $j\ge2$, $m_0=m_1=1$ ($\Phi(\frac12)=\frac12$, $\Phi(1)=0.8828125$) and $\eps_j=\abs D/m_j$: $\abs{\cB^D_j}\le\eps_j[(2/y_\ell+\frac54)F_j+D_j+\frac12\int(T_j'')^+w_j^2]$.
\end{named}
\begin{proof}
$P=\{y_\ell/2\le y\le y_\ell\}$ has length $d=y_\ell/2$, lies in $C_j$, and $T_j\ge m_jT_j(\lt)$ on $P$. The trace inequality $f(\lt)^2\le d^{-1}\int_Pf^2+2\norm f_{L^2(P)}\norm{f'}_{L^2(P)}$ with $f=w_j$ gives $T_j(\lt)w_j(\lt)^2\le m_j^{-1}((2/y_\ell+1)F_j+N_{j+1}^2)$. Conclude with (S3).
\end{proof}
\noindent\emph{(SE): $\cB^S_j\le0$.} If the error, the base coefficients and the residual are symmetric about the far-end normal line, write $2F_j$ as an integral over the doubled far field $[x^*,2\lt-x^*]$. On it the equation holds with the even/odd continuations. Under $x\mapsto2\lt-x$ every density is even, and the transport/$\partial_\tau T$ pair is mirror-symmetric, so no flux appears at $\lt$. The evenness of $T_j$ (W3) is what kills the fourth-order fluxes. On the corrected base the far end lies on the axis (Lemma~\ref{lem:SE}), so (SE) holds with $\cB^S_j=0$.

\begin{named}{Lemma J (time-weighted ramp data; \fn{LEMMA_F.md} \S4)}
Assume $\frac12\Lam_e'\le-\varrho\Lam_e-\kappa_{\mathfrak D}\mathfrak D+a_cY+\Xi'^2$ with $\kappa_{\mathfrak D}=\frac12$, and $\Yb'=-\varrho\Yb+\Phi$, $\Phi=a_c+\sqrt\varrho\Xi'$, $\sqrt{\Lam_e}\le\Yb\le1$, $\Yb\ge\Xi'/\sqrt\varrho$. Let $0<\lambda<2\min(\mu_1,\mu_2)$, $\omega(s)=e^{-\lambda(\tau-s)}$ and $J_X=\int_{\tau_0}^\tau\omega X^2$. Then:
(1) $\int\omega\mathfrak D\le\mathfrak B:=\kappa_{\mathfrak D}^{-1}[\frac12+(\bar\varrho-\lambda/2)^+/\lambda]$;
(2) $J_A\le1/\lambda$, $J_B\le4\mathfrak B/\theta_0$, $J_{D_4}\le8\mathfrak B/\eps$, $J_{W_1}\le\sqrt{J_AJ_B}+c_1J_A$, $J_{W_3}\le\sqrt{J_BJ_{D_4}}+c_1J_B$;
(3) $(\mu_1-\lambda/2)J_{V_0}+\frac12J_{V_2}\le K+C_x\gamma_0\sqrt{J_AJ_{V_0}}$, with $K=\frac12\omega(\tau_0)/\mu_e+a_fJ_{W_1}+b_fJ_B$;
(4) $J_{V_4}\le2\omega(\tau_0)/(\mu_e\eps_I)+M(\lambda)(c_3\sqrt{J_{D_4}}+c_2\sqrt{J_{W_3}}+c_1^I\sqrt{J_B}+c_0^I\sqrt{J_{W_1}}+C_x\gamma_2\sqrt{J_A})^2$, with $M(\lambda)=\min_\kappa4[(1+\Psi_1/\kappa)^2+(\Psi_1\kappa+2\Psi_1\sqrt{c_{1W}}+2c_{1W})^2/(4(\mu_2-\lambda/2))]$;
(5) $J_{V_1}\le\sqrt{J_{V_0}J_{V_2}}+c_{1W}J_{V_0}$ and $J_{V_3}\le\sqrt{J_{V_2}J_{V_4}}+c_{1W}J_{V_2}$;
(6) the residual terms of (I0), (I2) are absorbed with $\delta=10^{-3}$ and $\Yb_*=10^{-12}$.
\end{named}
\begin{proof}
(1) Multiply the closure by $\omega$ and integrate. Since $Y\le\Yb$ and $\Xi'^2\le\sqrt\varrho\Xi'\Yb$, we have $a_cY+\Xi'^2\le\Phi\Yb$, and $\int\omega\Phi\Yb=\frac12\Yb(\tau)^2-\frac12\omega(\tau_0)\Yb(\tau_0)^2+\int\omega(\varrho-\lambda/2)\Yb^2$. (2) follows from $\Lam_e\le\Yb^2\le1$, (1), the interpolations and Cauchy--Schwarz in time. (3), (4): multiply (I0), (I2) by $\omega$, use Young's inequality, and apply Minkowski in time.
\end{proof}
\noindent\emph{Certified inputs} (\fn{amp_A1_w180.json}, \fn{inputs.t7}): $\mu_e=0.036$, $\varrho=0.04711$, $\theta_0=0.025$, $\mu_1=0.15$, $\mu_2\ge0.6443$, $a_f\le0.0696$, $b_f\le0.0394$, $c^I\le(0.4251,0.4486,0.5390,0.4482)$, $c_{1W}\le0.09435$, $\Psi_1\le0.38851$, $C_{\gamma0}\le0.2069$, $C_x\le7.824$, $K_e\le890.6$, $\kappa_{\mathfrak D}=\frac12$. Result per unit $\Yb^2$: $J_0..J_4\le(114.1,64.9,25.7,6.81\cdot10^3,1.81\cdot10^6)$, with 74\% of $J_4$ from the $D_4$ term (the $1/\sqrt\eps$ loss). The condition $\lambda<2\mu_1$ holds for every rate up to T8's largest, $0.2520$, with $2\mu_1/\lambda_{\max}=1.19$; $\lambda$ is a free parameter of the lemma, so this is not a closure margin. $\varrho_c$ does not affect $\Yb$; it enters only through $(\bar\varrho-\lambda/2)^+$, which vanishes ($\lambda/2\ge0.12$).

\begin{named}{Corollary C (comparison and hand-off; \fn{LEMMA_F.md} \S5)}
\emph{C1.} If $F'\le-\lambda F+2g\sqrt F+h$ with $\lambda,g,h\ge0$, then $\sqrt F\le\sqrt H+G$, where $H=e^{-\int\lambda}F(\tau_0)+\int e^{-\int_s^\tau\lambda}h$ and $G=\int e^{-\frac12\int_s^\tau\lambda}g$.
\emph{C2.} Apply C1 level by level, with $\sqrt{F_i}\le\sqrt{H_i}+G_i$ for $i<j$ inside $g_j,h_j$ and $\int\omega h_j$ bounded through Lemma~J (\fn{camp.py}: per-slab sups at $\theta=0,\frac12,1$, running sup of $k_{jm}$; the $F_j$ start at the switch point).
\emph{C3} (windows). $z_j=(y\partial_y)^jv=\sum_mS(j,m)y^mw_m$, so $\sqrt{E_j^{[y,1.5y]}}[\mathrm{err}]\le\kappa_{\rm ch}\sum_mS(j,m)\sqrt{F_m}+c_{0,j}$; for $E_3$ the sum is $\sqrt{F_3}+3\sqrt{F_2}+\sqrt{F_1}$.
\emph{C4} (pointwise level 3). For $I=[y,1.5y]$: $\sup_I\abs{w_3}^2\le\int_Iw_3^2/\abs I+2\norm{w_3}_{L^2(I)}\norm{w_4}_{L^2(I)}$, with $\norm{w_3}^2\le F_3/(28y^4)$ and $\norm{w_4}^2\le D_2(\tau^*)/(28y^2)$.
\emph{C5} (good time). Let $W$ be the hand-off window (base slab 180, $\abs W=0.15$). Integrating (F) at $j=2$ over $W$ gives $\int_WD_2\le\theta_2^{-1}[F_2(\tau_W)+\int_W(2g_2\sqrt{F_2}+h_2)]$, so there is $\tau^*\in W$ with $D_2(\tau^*)\le\abs W^{-1}\int_WD_2$. All hand-off items are evaluated with sups over $W$, so they hold at $\tau^*$. Level 4 is not used anywhere.
\end{named}
\noindent C1 is proved by comparison with $Z=(\sqrt H+G)^2$, which satisfies $Z'\ge-\lambda Z+2g\sqrt Z+h$: on an interval where $F>Z>0$, $(F-Z)'\le(-\lambda+g/\sqrt Z)(F-Z)$, and Gr\"onwall from the last time with $F\le Z$ concludes.

\paragraph{Lemma N3$'$: nonlinear and modulation terms at levels 0--3.} \emph{Operator.} Let $\mathsf a=1-\kt e$ and $\mathsf J^2=\mathsf a^2+e'^2$. At the base modulation and with the clock fixed,
\[
F(e)=\frac{\mathsf J}{\mathsf a}\Bigl[-\kappa_{ss}+\frac{c(-\tilde qe'+(\tilde p+e)\mathsf a)+\tilde b(-\tilde t_1e'+\tilde n\mathsf a)}{\mathsf J}\Bigr]-\tilde V+\frac{e'(\varpi-e\tilde\theta_\tau)}{\mathsf a},
\]
$F=A_4e''''+A_3e'''+A_2$ with $A_4=-\mathsf J^{-4}$, and $\cN:=F(e)-F(0)-DF(0)e=\tilde Ae''''+\tilde Be'''+\tilde C$, where $\tilde A=1-\mathsf J^{-4}$, $\tilde B=A_3(z)-A_3(0)$, $\tilde C=A_2(z)-A_2(0)-\nabla A_2(0)\cdot z$ and $z=(e,e',e'')$. \emph{Cancellation} (\fn{n3sym.py}, symbolic): the quadratic tangential parts cancel those of the normal velocity, so $\cN=\cN_K+(\kt'''-r')ee'/\mathsf a$, with $\cN_K$ the nonlinear part of $-(\mathsf J/\mathsf a)\kappa_{ss}$; no coefficient grows with $x$. \emph{Modulation} (\fn{n3qsym.py}): $\cN+\cQ=(1-\vartheta)\cN+\vartheta e''''+\vartheta\kt^2e''-\vartheta c_0e+\vartheta P+g\,e'-(\vartheta c+x_1)e$, with $P=(\kt'''-r')ee'/\mathsf a$, $g=U/\mathsf a+5\vartheta\kt\kt'$ and $U=(x_1+\vartheta c)\tilde q+(x_2+\vartheta\tilde b)\tilde t_1$; the last two terms go through the strong form, and $g(\lt)=0$ on the corrected base. \emph{Weak form, level 3.} $S:=2\int T_3w_3\partial^4\cN=2\int(T_3w_3)''\partial^2\cN$ with $\partial^2\cN=\tilde Aw_5+a_4w_4+a_3w_3+a_\flat$; this is linear in the free jets, so there are no trilinear classes, and one more integration by parts gives
\[
S=2\!\int\!T\tilde Aw_5^2+\!\int\!\Psi_{45}w_4w_5+\!\int\!\Psi_{35}w_3w_5+\!\int\!\Psi_{34}w_3w_4+\!\int\!\Psi_{33}w_3^2+2\!\int\!(Tw_5+2T'w_4+T''w_3)a_\flat
\]
(coefficients checked by \fn{n3wcheck.py}). Levels 0--2 are analogous. \emph{Majorants.} Per $x$-range (the ramp, $[x_b,30.25]$, $[30.25,34.25]$ and dyadic ranges $[2^k\cdot34.25,2^{k+1}\cdot34.25]$ up to $\lt$), the base data are \fn{arb} hulls of $\kt^{(\le5)}$, $r',r'',r'''$, $\tilde p,\tilde q,\tilde n,\tilde t_1$. The box enters through $(Z_0,Z_1,Z_2,Z_3)=(1.7E(x_{hi}),b_0,b_1/x_{lo},b_2/x_{lo}^2)$ for $y\ge28$ and $(\eta_1,\eta_1,\eta_F/x_{lo},\eta_F/x_{lo}^2)$ below. Young's inequality with fixed weights produces $\Delta\theta_j,\Delta\sigma_j,\Delta k_{jj},\Delta d_{jm},\Delta k_{jm},\Delta c_{jm},hf_j$ (formulas in \fn{LEMMA_F.md} \S6.1).

\begin{named}{Lemma N3$'$}
Assume $\mathbb B'$, (SE), $\abs x\le\bar x$ and $\abs\vartheta\le\bar x/n$, with $\bar x=C_x\Yb_{\rm box}+x_r$. Then (F) holds at each $j\le3$ for the equation with $\cN+\cQ$ included, provided $\theta_j,\lambda_j,k_{jj},k_{jm},c_{jm}$ are replaced by $\theta_j-\Delta\theta_j$, $\lambda_j-\Delta\sigma_j$, $k_{jj}+\Delta k_{jj}$, $k_{jm}+\Delta k_{jm}$, $c_{jm}+\Delta c_{jm}$, and $h_j$ gains $\sum_{m<j}\Delta d_{jm}F_m+hf_j$. All of these are per-slab \fn{arb} bounds (\fn{n3far.py}, run at $\Yb_{\rm box}=1.7\cdot10^{-12}$, boxes $b_0,b_1,b_2=1.5\cdot10^{-3},3\cdot10^{-4},10^{-2}$, factor $1.7$; then \fn{camp.py --n3}, iterated once).
\end{named}

\begin{table}[htbp]\centering\small
\caption{Far-field constants and amplification on the direct residuals (\fn{constants_b1.json}; \fn{n3_b1.json}; \fn{amp_A1_w180.json}, window slab 180). $\sqrt{F_j}\le A_j\Yb+B_j$, where $B_j$ is the $\Yb$-independent part (initial data and residual).}\label{tab:far}
\begin{tabular}{@{}cccccccc@{}}\toprule
$j$ & $\min\lambda_j$ & $\theta_j$ & $\Delta\theta_j$ (N3$'$) & $\Delta\sigma_j$ & $A_j$ & $B_j$ & $\sqrt{F_j}$ at $\Yb=10^{-12}$\\\midrule
0 & 0.2400 & 1.2253 & $\le0.0161$ & $\le1.93\cdot10^{-3}$ & $6.58\cdot10^3$ & $7.61\cdot10^{-7}$ & $7.68\cdot10^{-7}$\\
1 & 0.2487 & 1.2253 & $\le0.0161$ & $\le1.21\cdot10^{-3}$ & $6.17\cdot10^4$ & $1.76\cdot10^{-9}$ & $6.34\cdot10^{-8}$\\
2 & 0.2484 & 1.2253 & $\le0.0161$ & $\le1.23\cdot10^{-3}$ & $1.13\cdot10^6$ & $5.11\cdot10^{-8}$ & $1.18\cdot10^{-6}$\\
3 & 0.2456 & 0.2201 & $\le0.0863$ & $\le2.84\cdot10^{-3}$ & $1.13\cdot10^8$ & $7.82\cdot10^{-6}$ & $1.21\cdot10^{-4}$\\\bottomrule
\end{tabular}\end{table}
\noindent\emph{Results.} $\theta_3-\Delta\theta_3\ge0.133$. The initial data contribute $\le3.8\%$ of each $\sqrt{F_j}$. At $\Yb=10^{-12}$ and the good time in slab 180: $D_2(\tau^*)\le9.40\cdot10^{-12}$ (parts: residual $1.44\cdot10^{-14}$, linear $6.46\cdot10^{-13}$, quadratic $8.74\cdot10^{-12}$; $\theta_2-\Delta\theta_2\ge1.209$), and the error part of the cone $K_4$ is $\le2.18\cdot10^{-3}$. The item limits (largest admissible $\Yb$, computed with the exact KCH, $c_0$ and $c_0^{\rm tip}$ of \S\ref{ss:win}) are: item~3 (windows) $E_3$ $1.575\cdot10^{-11}$ (binding), $E_2$ $1.81\cdot10^{-9}$, $E_1$ $3.82\cdot10^{-8}$, $E_0$ $3.31\cdot10^{-8}$; item~4 ($K_4$ at the good time) $2.89\cdot10^{-11}$; item~2 (shells) $E_2$ $3.63\cdot10^{-10}$, $E_1$ $9.85\cdot10^{-10}$, $E_0$ $2.39\cdot10^{-7}$. So $\Yb=10^{-12}$ has margin $\times15.7$. The final tail runs the T8 iteration and the window \fn{camp} twice (passes A, B), and \fn{t9.py} reproduces bit for bit between them.

\subsection{The nonlinear boxes}\label{ss:boxes}
\emph{Verified by} \fn{p2/t8se/n3boxes.py} (review B2: it reads the boxes from the N3$'$ output and asserts all four) with \fn{p2/boxes.py}; \emph{log} \fn{p2/out/n3boxes_final_B.log}. C4 on $I\subset[y/1.5,1.5y]$ with $\abs I\ge y/3$, and $\sqrt{F_j}\le A_j\Yb+B_j$ per slab, give the pointwise bounds; for example $\abs{e'}^2\le\frac{6.75}{28}yF_0+\frac3{28}y\sqrt{F_0F_1}$, up to $y_\ell$.
\begin{table}[htbp]\centering\small
\caption{Bootstrap boxes. N3$'$ is certified on these boxes at $\Yb_{\rm box}=1.7\cdot10^{-12}$; they close there, and the margins at the $\Yb=10^{-12}$ in use are in the last column.}\label{tab:boxes}
\begin{tabular}{@{}L{0.28\linewidth}L{0.22\linewidth}L{0.24\linewidth}L{0.16\linewidth}@{}}\toprule
box quantity & box & bound at $\Yb_{\rm box}=1.7\cdot10^{-12}$ & margin at $\Yb=10^{-12}$\\\midrule
far $\abs{e'}$ ($y\ge28$) & $b_0=1.5\cdot10^{-3}$ & $2.38\cdot10^{-4}$ & $\times6.42$\\
far $x\abs{e''}$ & $b_1=3\cdot10^{-4}$ & $1.08\cdot10^{-4}$ & $\times4.69$\\
far $x^2\abs{e'''}$ & $b_2=10^{-2}$ & $8.34\cdot10^{-3}$ & $\times1.99$\\
far $\abs e$: $\sqrt{F_0}$ ratio & factor $1.7$ & $1.00001$ & \thin{$\times1.69$}\\
tip $\abs e,\abs{e'}$ ($x\le34.25$) & $\eta_1=10^{-5}$ & $2.78\cdot10^{-8}$ (Lemma~4b) & ---\\
$B\le\Yb/\sqrt\eps$ & $\eta_B=10^{-4}$ & $3.2\cdot10^{-10}$ at $\Yb=10^{-12}$ & ---\\
$\abs\vartheta$ & $10^{-3}$ & $3.37\cdot10^{-8}$ at $\Yb=10^{-12}$ & ---\\\bottomrule\end{tabular}\end{table}
\noindent On the direct residuals the far $\abs{e'}$ box has margin $\times6.42$ (it was $\times1.29$ on the \fn{se_corr} residuals, review B2). The thinnest box is the $\abs e$ box, through the factor $1.7$ on $\sqrt{F_0}$.

\section{The hand-off (T9)}\label{s:T9}
\emph{Verified by} \fn{p2/t9h/} (\fn{h1.py}, \fn{h1dt.py}, \fn{clam.py}, \fn{win.py}, \fn{win_slab.py}), \fn{p2/t9.py} with \fn{p2/t9chart.py} and \fn{p2/curvjets.py}, \fn{p2/t9_far.py}, \fn{p2/t1clock.py}; \emph{certificates} in \fn{p2/out/} (Table~\ref{tab:hashes}). All functionals are evaluated on $\Gt(\tau)$ in \fn{arb} for $\tau$ anywhere in base slab 180; the error enters by the triangle inequality, through $C_\alpha$, $C_\Lam$, KCH, $c_0$ and $c_0^{\rm tip}$.

\subsection{The one-slab good-time window (review A1)}\label{ss:window}
The good-time window is all of base slab 180: $t_c\in[-1,1]$, $\Delta\tau=0.15$, $k\in[2.27\cdot10^4,2.38\cdot10^4]$. Every hand-off item below is certified uniformly for $\tau^*$ anywhere in the slab, and Corollary~C5 gives a good time $\tau^*$ in it with $D_2(\tau^*)\le9.40\cdot10^{-12}$. The earlier window $W'$ ($t_c\in[1-\frac1{512},1]$, length $1.46\cdot10^{-4}$) gave $D_2(\tau^*)\approx1.2\cdot10^{-8}$, because $D_2(\tau^*)$ is a mean over the window; Lemma~O H4 then failed by $\times5$ to $\times280$ (A1). With the whole slab, H4 holds with margin $\times4.35$ (\S\ref{ss:H4}). The price is that H1--H4 must hold on a window $1024$ times longer, which is what the Krawczyk points and derivative balls of \S\ref{ss:alpha} provide.

\subsection{\texorpdfstring{The $S_1$ chart and the certified set (review B8)}{The S1 chart and the certified set (review B8)}}\label{ss:S1}
The chart uses only step-1 objects (\fn{s1_tight/s1_tight.tex}). For $\sigma\le8$ it uses \fn{s1_tight/ext_tip/prof_tip.txt}, a drop-in replacement for \fn{prof_tight.txt}: same 51\,246 rows, every box nested in the \fn{prof_tight} box, $\theta$-width $\le3.2\cdot10^{-12}$ on $x\in[0,2]$. For $\sigma\in[8,13]$ it uses \fn{ext813/ext813_prof.txt} (25\,000 rows, width $\le5.7\cdot10^{-10}$, CAPD $C^0$ Lohner, order 25). Beyond $\sigma=13$ it uses tail envelopes. The asymptotic tangent angle $\Theta$ uses $\int_\sigma^\infty k=p/\sigma+4m/\sigma^2+\dots$.

\emph{Non-uniqueness.} The state of $S_1$ at $\sigma=8$ is a point $y^*$ of the certified set $Y=x_m+C\cdot r_0+r_r$ (\fn{in_br2.txt}) with $M(y^*)=0$, where $M$ is $m$ at the first return to $q=0$. \fn{ext_tip} evaluates
\[
\Phi_\tau(y^*)\in\Phi_\tau(x_m)-\frac{G_0}{A_0}(M(x_m)+DM\cdot e)+\sum_{j\ge1}\Bigl(G_j-\frac{G_0A_j}{A_0}\Bigr)\delta_j+D\Phi_\tau\cdot e,\qquad A=DM\cdot C,\ G=D\Phi_\tau\cdot C,
\]
by the mean value theorem on the convex $Y$, with $DM$ and $D\Phi_\tau$ enclosed over all of $Y$ ($A_0=-4850.43$). The formula holds for \emph{every} zero of $M$ in $Y$, so no uniqueness is used; step~1 does not prove uniqueness. The conditioned tip data are $\tau_{\rm tip}\in[10.249077651120778,10.249077651121361]$, $k(0)\in[1.5398206378718,1.5398206378722]$, $D=-p(\rm tip)\in[2.50080627661969,2.50080627662005]$ and $\varepsilon\in[1.5346198894026,1.5346198896816]\cdot10^{-4}$, for all zeros. Every object T9 uses is an interval enclosure valid for every zero: \fn{prof_tip}, \fn{ext813}, the 3425 IBP hull boxes of step~3 (\fn{outA5.txt}), and the fixed functions $\hat\psi_3,\hat h_2$ of \fn{c4fun}. T9 multiplies hulls, so its bounds hold for each zero, and Theorem~4 is stated for any $S_1$ in the certified set. Distinct zeros differ by $\le3\cdot10^{-12}$, far below the step-3 floor $4.9\cdot10^{-10}$. Result: $\Theta=1.8617364190148\pm5.2\cdot10^{-11}$.

\emph{Containment check} \tC\ (\fn{/home/claude/s1_tight/consistency/check_outA5.py} $\to$ \fn{check_outA5.json}, ball arithmetic, 128 bits): on each overlap the tight enclosures, restricted by second-order Taylor forms, lie inside the step-3 hulls of \fn{outA5.txt}, component by component: all $1025$ rows of \fn{prof_tip} and all $1025$ rows of \fn{prof_tight} against stage C, and all $501$ overlapping rows of \fn{ext813} against stage B, with positive relative margins and $\sigma_{\rm tip}$ inside. So the step-2/3 objects and the tight $S_1$ describe the same certified set.

\subsection{\texorpdfstring{(H1): $\alpha$ and $\sqrt\Lambda$ by the step-3 IFT}{(H1): alpha and sqrt(Lambda) by the step-3 IFT}}\label{ss:alpha}
The step-3 coordinates at $\tau$ are $(x_c,\lambda,\alpha)$, determined by the three constraints $\ip{\chi w}{f_j}_\rho=0$, with $w=u-\alpha\hat\psi_3-\alpha^2\hat h_2$ and $f_3$ the adjoint function, on the 3425 step-3 IBP hull boxes. The tip offset $C=-\pi/2-\theta({\rm tip})+dT\,k_{\rm tip}$ enters the projection as one scalar (\fn{COHERENT=1}): $\Lam$ is evaluated at the midpoint of $C$, and the rotation mode's norm (by \fn{h1.py}, $\le17.65$) times $\operatorname{rad}C=1.72\cdot10^{-12}$ bounds the rest.

\begin{lemma}[H1 on slab 180 \tC; \fn{t9h/h1.py}, \fn{h1dt.py}, \fn{clam.py}]\label{lem:H1}
(i) At each of the 16 exact points $t_k=-1+(2k+1)/16$, $k=0..15$, a Krawczyk test proves that the three constraints have exactly one zero $p(t_k)=(x_c,\ln\lambda,\alpha)$ in the certified box; there $\sqrt{\Lam(t_k)}\le2.01\cdot10^{-9}$, with $\Lam$ evaluated using the exact $u''''$ cell hull.
(ii) On each of 32 $t_c$-balls of radius $1/32$, $\sup\norm{\partial_tw}_\Lam\le3.33\cdot10^{-7}$, where $dw/dt=u_t+\sum_i\dot p_im_i$, with exact first variations $m_i$ and $\dot p$ from the IFT by defect correction; the parameter box is verified a priori.
(iii) Hence, for the base, $\sqrt{\Lam(t)}\le\sqrt{\Lam(t_k)}+\abs{t-t_k}\sup\norm{w_t}_\Lam\le2.282\cdot10^{-8}$ and $\abs{\tilde\alpha}\le1.1739\cdot10^{-4}$ on the whole slab (\fn{h1slab/summary.json}).
(iv) For the true curve (\fn{clam.py}): $\sqrt{\Lam[w]}\le\sqrt{\Lam[\tilde w]}+C_\Lam\Yb$ and $\abs{\alpha-\tilde\alpha}\le C_\alpha\Yb$, with $C_\Lam\le24.2$ and $C_\alpha\le0.721$, using $\Lam^{S3}[e]\le(0.08/\mu_e)\Lam_e$, Cauchy--Schwarz on $\ip{\chi e}{f_j}_\rho$ with $\norm{f_j}_\rho\le1.838,1.191,0.0196$, and Lemma~4b on the cutoff tail. At $\Yb=10^{-12}$: $\sqrt\Lam\le2.29\cdot10^{-8}$ (it holds for $\Yb\le1.12\cdot10^{-9}$) and $\abs\alpha\le1.1739\cdot10^{-4}+7.3\cdot10^{-13}$.
\end{lemma}
\noindent Review B3: the Krawczyk test replaces the earlier Newton step with a typed remainder ball; the fall-back bound in \fn{fU4_safe} is gone (any failure is an error); the factor $1.01$ in \fn{clam.py} is now justified by $\abs{u'},\abs u\le2\cdot10^{-4}$ and $k_1\le1.55$. The rotation-mode norm still uses generous hand bounds for the frame-mode derivatives on $[10.25,60]$ (\fn{h1.py}, block ``E\_I part''); they enter only through $\operatorname{rad}C$, a contribution $\le3.1\cdot10^{-11}$ to $\sqrt{\Lam(t_k)}$, against the threshold $5\cdot10^{-8}$. Review B9: cells are exact dyadic intervals tiling $[0,60]$.

\subsection{(H2), (H3): shells and windows; KCH and the chart (review B4, B5)}\label{ss:win}
The step-3 definitions are reproduced in \fn{win.py} and were checked by the review: $w=u-\alpha\hat\psi_3-\alpha^2\hat h_2$; $\chi_{32}$; $\zeta$ on $[10.25,16.25]$ and $[57,60]$; $\eps,\eps_I,\mu=0.08$; $\omega=28/y$; $U=1.3056\,\alpha\,y^{1.0246}$; a sliding cover of $y\in[50,253]$. The base values on slab 180 use 32 $t_c$-balls of radius $1/32$ with the parameter paths $p(t_k)+(t-t_k)\dot p$ (\fn{win_slab.py} $\to$ \fn{h1slab/win_summary.json}): shells $(8.01\cdot10^{-11},1.34\cdot10^{-10},2.94\cdot10^{-10})$ and windows $(3.48\cdot10^{-8},1.27\cdot10^{-7},3.76\cdot10^{-7},9.71\cdot10^{-7})$. \fn{camp.py} reads these JSON outputs (review B4).

\emph{The chart constants, computed (review B5; \fn{p2/t9chart.py}, arb Taylor mode, no sympy at run time).} The windows' error parts come from C3 with the chart constant KCH and the chain-rule terms $c_0$, $c_0^{\rm tip}$, all computed with no typed allowance. The split is exact:
\[
D^jf_t(y)-D^jf_b(y)=[Q^t_j-Q^b_j](X_t(y))+[Q^b_j(X_t(y))-Q^b_j(X_b(y))],
\]
where $Q_j$ is $D^j$ of the slope at one base point $x$; the true curve has $F_t=(a\sin\varphi+e'\cos\varphi)/(a\cos\varphi-e'\sin\varphi)$, $a=1-\kt e$, $Y_t=Y_b-e\sin\varphi$, $D_t=(Y_t/Y_t')d/dx$. The first bracket is an exact linear part plus a remainder charged to $\eta_i=y_0^{i-1}e^{(i)}$ with the N3$'$ boxes and $\abs e\le c_e\Yb+\sqrt{F_0}(y^3/84)^{1/2}$ (Lemma~4b); the second is $\le\sup\abs{dQ^b_j/dx}\,\abs e\tan\varphi^*$. \fn{t9chart_check.py} checks offline that the linear part equals $D^j(\sec^2\varphi\,e'+\sec^2\varphi\tan\varphi\,\kt e)$ (sympy, exact) and that the remainder bound holds on 200 random jets (worst $0.99$ of it). Results (\fn{t9_cone_kch.json}): $\mathrm{KCH}\le1.18778$ (linear part $1.18355$ times the computed nonlinear factor $1.00358$); $c_0\le(2.69\cdot10^{-3},2.20\cdot10^{-2},4.71\cdot10^{-2},9.87\cdot10^{-2})$; $c_0^{\rm tip}\le(0.111,0.114,0.118,0.125)$ per unit $\Yb$; $\sup\abs\varphi\le0.0779$.

\emph{The ray coordinate} is certified by a tip anchor: $y_{\rm ray}(x)=p(0)\cos\Theta+x-\int_0^x(1-\cos\varphi)$, with the integral bounded on the arrangement of the cells' $x$-hulls, and $q(0)=0$, $\theta(0)=-\pi/2$ by symmetry (asserted against the tip hulls). The hand-off frame enters through $y_h=(y_{\rm ray}-x_c\cos\Theta)/\lambda$, in $D_h$ and in the weight $28\lambda M/(Y-c)^2dx$, for the windows and the shells. $(x_c,\lambda)$ is the union of the 32 IFT balls of slab 180, widened by the true-curve shift $2\max(C_{\rm tr},C_{\rm sc})\Yb_{\rm box}\le3.3\cdot10^{-11}$ ($C_{\rm tr}\le9.66$, $C_{\rm sc}\le5.18$, from \fn{clam.py}'s transfer): $x_c\in[1.007,1.047]\cdot10^{-4}$, $\lambda\in[0.99999490,0.99999511]$, $c=x_c\cos\Theta\in[-3.013,-2.878]\cdot10^{-5}$.

\emph{The shells, by their own bound.} T8's $F_m$ cannot carry the shells: their ramp weights vanish at $y_{T8}=14$. \fn{t9.py} (block \fn{shell}) bounds the shells by $\Lam_e$ alone, per unit $\Yb$: with $v=\zeta_\infty e'$, $V_m=\norm{v^{(m)}}_W$, $\Lam_e\ge\mu_e(V_0^2+\eps_IV_2^2)$ and $V_1\le k_{V1}V_0+V_2/2$, it gets $V_0,V_1,V_2\le5.28,486,966$ per unit $\Yb$ (with the $\zeta$ corrections on $x\in[15.94,16.25]$, $\zeta\ge0.99977$), then the same exact split as the windows, $\abs e$ from Lemma~4b and T7's far bound, and Agmon on $\psi_c(1+x^2)^{-1}v^{(m)}$ for the remainder boxes. Result: $\sqrt{E_j^{\rm sh}[\mathrm{err}]}\le1.833\cdot10^3,\ 1.014\cdot10^7,\ 1.232\cdot10^9$ times $\Yb$ ($j=0,1,2$), so the shell limits are $2.39\cdot10^{-7}$, $9.85\cdot10^{-10}$, $3.63\cdot10^{-10}$, and (H2) holds with margin $\times363$ in $\Yb$.

\emph{The windows.} With these constants, \fn{camp.py --window 180-180} gives the item-3 limit $\Yb\le1.575\cdot10^{-11}$ ($E_3$, binding) (\S\ref{ss:far}), so (H3) holds with margin $\times15.7$ in $\Yb$. \fn{t9.py} uses only the amplification's per-slab $a_F,b_F$ and the values at $\Yb$; the final tail runs \fn{t9}, \fn{camp}, \fn{t9}, and the second \fn{t9} output is identical (manifest check \fn{t9_fixed_point}).

\subsection{(H3) \texorpdfstring{$K_4$}{K4} and (H4)}\label{ss:cone}
\emph{Cone cylinders} (\fn{t9.py}; slabs 177--180 hulls; exact normal-graph curvature jets by \fn{curvjets.py}). For $r_c\in[200,380]$, $L=r_c/4$, $X=3$ (half-width $0.75r_c$), the graph is taken over the centre tangent: $h''=k\sec^3$, $h'''=\sec^4(k_s+3k^2t)$, $h''''=\sec^5(k_{ss}+10kk_st+12k^3t^2+3k^3\sec^2)$, with $K_j=L^{j-1}\sup\abs{h^{(j)}}$ and $N_3=L^{3/2}\norm{h'''}_{L^2(\abs x\le XL)}$. Any frame with $\abs{x_c}\le0.5$, $\abs{\lambda-1}\le10^{-3}$ is covered. The error parts come from C4 at the inner end $y=0.25r_c-X_0$, with $D_2(\tau^*)$ from \fn{camp}. \emph{(H3) $K_4$:} on the cylinder $r_c=200$ ($y\in[50,350]$), $K_4\le5.60\cdot10^{-3}$ (base $3.1633\cdot10^{-4}$) against $0.06$: $\times10.7$. \emph{(H4):} the base satisfies $\abs{\theta-\Theta}\le0.078$ over the ray to $y=350$, and the error adds $\le5.5\cdot10^{-4}$, so the curve is a graph over the $S_1$ ray. The cone and far Lemma~O data are in \S\ref{s:LO}.

\subsection{\texorpdfstring{The hand-off time $t_1$ (Theorem~4(b); review B7)}{The hand-off time t1 (Theorem 4(b); review B7)}}\label{ss:t1}
\begin{lemma}[\tC; \fn{p2/t1clock.py} $\to$ \fn{t1clock.json}]\label{lem:t1}
For $\tau^*$ anywhere in slab 180, $t_1=t(\tau^*)\in[4.597857477538968,\,4.597858614388128]\subset[4.5978574775,\,4.5978586144]$. On the slab, $\ln\tilde\lambda\in[-9.64236,-9.60485]$, and the base curve's maximal curvature is in $[22796.37,\,23759.61]\subset[2.279\cdot10^4,\,2.376\cdot10^4]$.
\end{lemma}
\begin{proof}
$t_1-t_{96}=\int_{\tau_a}^{\tau^*}\lambda^4(1-\vartheta)\,d\tau$ with the true $\lambda$. The script encloses $\int\tilde\lambda^4d\tau$ per slab by $0.15\,[e^{4\,lo},e^{4\,hi}]$ from T6's certified range of $\ln\tilde\lambda$ (outward); it takes the fraction $(1-t_{c,a})/2$ of slab 2 after the switch and a fraction in $[0,1]$ of slab 180. The clock factor $1+[-\theta_s,\theta_s]$ and the scale factor $(\lambda/\tilde\lambda)^4\in e^{\pm4\rho_s}$ are per slab (Remark~\ref{rem:t1}), with a start term $\pm s\lambda^4\le1.1\cdot10^{-17}$; on slab 180, $\theta_{\max}=3.3641\cdot10^{-8}\le3.37\cdot10^{-8}$ and $\abs{\ln(\lambda/\tilde\lambda)}\le3.38\cdot10^{-8}$. The base $k_{\max}$ is the tip-curvature hull over slab 180 divided by $\tilde\lambda$ (unchanged).
\end{proof}
\begin{remark}[the true scale factor, in the script (cleanup O6a)]\label{rem:t1}
\fn{t1clock.py} now uses $dt=\lambda^4(1-\vartheta)d\tau$ with the true $\lambda$ (the first draft integrated the base's $\tilde\lambda^4$ and bounded the omission by hand, $\le2.32\cdot10^{-7}$ in $\ln(\lambda/\tilde\lambda)$). Since $d\ln(\lambda/\tilde\lambda)/d\tau=a-\tilde a=x_1+\vartheta c$, per slab $\abs{\ln(\lambda/\tilde\lambda)}\le\rho_s=\rho_a+\sum\Delta\tau\,(\bar x_s+\theta_s\bar c_s)$. Here $\rho_a=\mu_1+s\bar c_2\le4.14\cdot10^{-13}$ comes from the switch's IFT (\fn{switch_s2.json}), $\bar x_s=C_x\Yb+x_{r,s}+x_{\rm NL}$ with $x_{r,s}$ computed per slab as in \fn{n3far} (maximum $2.57\cdot10^{-16}$; the typed $2.0\cdot10^{-16}$ was below it) and $x_{\rm NL}\le1.95\cdot10^{-15}$, and $\theta_s=\bar x_s/n_{V,s}$. The result is $t_1\in[4.597857477538968,4.597858614388128]$ (before the cleanup $[4.597857477538703,4.597858614388391]$: the per-slab $\theta_s$ tightens more than $\lambda/\tilde\lambda$ widens), inside the quoted $[4.5978574775,4.5978586144]$; $\abs{\ln(\lambda/\tilde\lambda)}\le3.38\cdot10^{-8}$ on slab 180 (the hand bound was $2.32\cdot10^{-7}$). The base $k_{\max}$ hull is unchanged. This item is closed.\n\end{remark}
\noindent\emph{The true curve's $k_{\max}$} is not certified: the curvature of the true curve at the tip is $\kt+e''+\dots$, and P2 controls $e''$ on the tip only in $L^2$ and through $S\le S_W+\dots$ with dissipation (Lemma~4a(iii)), not pointwise in time. Only the base value is certified \OPEN{\S\ref{s:open}}. No hypothesis of step~3 uses $k_{\max}$; the scale $\lambda_1$ comes from $\ln\tilde\lambda$ (\S\ref{s:main}).

\subsection{(H6): regularity}\label{ss:H6reg}
(H6) asks that $\gamma(t_1)$ be a smooth closed immersed curve, and fixes the forcing norm $\norm{\chi\Phi_F}$ as the one determined by $\alpha(\tau_1)$ (a definition in step~3 \S3.3; nothing to check). By Lemma~\ref{lem:gal} the solution lies in $C^0H^5$ on every P1 slab, and Theorem~P2 continues it to $t_1$; it is smooth for $t>0$ by parabolic regularity, and closed and immersed as the CDF evolution of a closed immersed curve \tA. It is not embedded at $t_1$ (\S\ref{ss:embed}).

\section{Lemma O data (H5)}\label{s:LO}
\subsection{What step 3 asks}\label{ss:LOfam}
(H5) of step~3 (v3.1 l.97, as amended by v3.2, \fn{step3_amend_v32/amend_v32.tex} \S3) asks that Lemma~O's hypotheses H1--H5, H7, H9 hold \emph{arc by arc} for the cylinder families covering physical $r\ge150\lambda_1$, and adds six data thresholds and far thresholds in $\lambda_1$ units. The families are: \emph{cone}, centres $r_c\in[200,380]$, unit $L=r_c/4$, $X=3$, $M=0.12$, $K_2=K_3=0.18$ (unchanged by v3.2); \emph{far}, centres $\ge376$, unit $L=113$, $X=2$, now $M=0.15$, $K_2=K_3=0.3$ (it was $M=1$, $K=1.5$ in v3.1). Lemma~O's hypotheses (\fn{writeup/lemmaO/lemmaO_v3.tex}) are: H1 graph, H2 slope, H3 jets $\le M$, H4 $\int_{-X}^X(h_0^{(5)})^2\le M^2$, H5 $\norm{h_0'''}_{L^2}$, H6 the a-priori bounds close at every centre, H7 time ratio, H8 distance to other arms, H9 regularity. v3.2 proves H6 in step~3, conditional on Lemma~R (\S\ref{s:H6}); H7 and H9 are internal to step~3; H8 is not required. Table~\ref{tab:LO} collects the step-4 data.

\begin{table}[htbp]\centering\small
\caption{Lemma~O data at $\tau^*$ (any $\tau^*$ in base slab 180) against step~3 v3.2. Own cylinder units unless marked. Sources (\fn{step4/p2/out/}): \fn{t9_cone_kch.json}, \fn{t9_v32_checks.json} (cone); \fn{t9_far.json}, \fn{t9_far_v32.json} (far); \fn{h4loc_A1_w180.json}, \fn{t8se/constants/h4loc_v32.json} (H4); \fn{t9_lemmaO_v32.json}, \fn{t9_lemmaO_far.json} (H1).}\label{tab:LO}
\begin{tabular}{@{}L{0.13\linewidth}L{0.40\linewidth}L{0.40\linewidth}@{}}\toprule
Lemma~O & cone: $r_c\in[200,380]$, $L=r_c/4$, $X=3$, $M=0.12$ & far: centres $\ge376$, $L=113$, $X=2$, $M=0.15$\\\midrule
H1 graph & graph over $[-X,X]$ in all 18 cylinders; $\sup\abs{h_0}\le0.1872$, $H\le0.1895$ ($C_0\sqrt{t_1}=0.0023$) & every centre from base arclength $375.125$: turning $\le0.04095$ rad within $\pm227.125$, $\sup\abs{h_0}\le0.04072$, $H=0.0417$ ($C_0\sqrt{t_1}\le9.212\cdot10^{-4}$); closed, \fn{t9_lemmaO_far.py} \tC\ (\S\ref{ss:LOH1})\\
H2 slope & $\le0.0573$; data threshold $s_0=0.08$: $0.05724$, \thin{$\times1.40$} & $\le0.04153$; $s_0=0.08$: $\times1.92$; against $M$: $0.0416$, $\times3.61$\\
H3 jets & $K_2,K_3,K_4\le1.89\cdot10^{-2},4.78\cdot10^{-4},1.23\cdot10^{-2}$ against $0.12$: $\times6.3$, $\times251$, $\times9.7$; data threshold $m_2=0.03$: $0.01883$, \thin{$\times1.59$} & $\lambda_1$ units: $\abs{h_{xx}},\abs{h_{xxx}},\abs{h_{xxxx}}\le1.797\cdot10^{-4},2.595\cdot10^{-8},6.964\cdot10^{-9}$ against $3.53\cdot10^{-4},1.17\cdot10^{-5},1.03\cdot10^{-7}$: $\times1.96$, $\times450.9$, $\times14.79$; $m_2=0.04$: $0.02031$, $\times1.96$; $h_3,h_4$ against $M$: $3.313\cdot10^{-4},0.01005$, $\times452.8$, $\times14.92$\\
H4 $\int(h_0^{(5)})^2$ & $\le0.2393\,M^2$ (worst, $r_c\in[370,380]$): $\times4.18$ & $\le0.1761\,M^2$ (centre 376): $\times5.68$\\
H5, $\norm{h_0'''}_{L^2}$ & $N_3\le1.170\cdot10^{-3}$; data threshold $n_3=0.004$: $\times3.41$ & $\le6.626\cdot10^{-4}$; $n_3=0.004$: $\times6.03$; $\lambda_1$ units $5.516\cdot10^{-7}$ against $3.32\cdot10^{-6}$: $\times6.01$\\
H6 & \multicolumn{2}{L{0.80\linewidth}}{proved in step~3 v3.2 from the data thresholds and the strip statement (closure in \S\ref{ss:v32}), conditional on Hypothesis~(F) (Lemma~R is proved, \S\ref{ss:v32status}); not met under v3.1 (\S\ref{ss:H6v31})}\\
H7, H9 & \multicolumn{2}{L{0.80\linewidth}}{internal to step~3: $t^*\ge3.683\cdot10^{-7}$ against $1.6001\cdot10^{-7}$ (cone, $\times2.30$) and $9.511\cdot10^{-8}$ against $6.134\cdot10^{-9}$ (far, $\times15.5$); regularity}\\
H8 & \multicolumn{2}{L{0.80\linewidth}}{not required (arc by arc). Cone: distance $\ge0.136L$ against $0.1$ (informational). Far: the separation test fails at the self-crossing (\S\ref{ss:embed})}\\\bottomrule
\end{tabular}\end{table}

\subsection{Slope, jets and \texorpdfstring{$L^2$}{L2} norms (H2, H3, H5, and the data thresholds)}\label{ss:LOjets}
\emph{Cone} \tC\ (\fn{t9.py} $\to$ \fn{t9_cone_kch.json}; thresholds checked by \fn{t9_v32_checks.py} $\to$ \fn{t9_v32_checks.json}, all pass): on the cylinders of \S\ref{ss:cone}, base plus error at $\Yb$: $K_2\le1.89\cdot10^{-2}$, $K_3\le4.78\cdot10^{-4}$, $K_4\le1.23\cdot10^{-2}$, $N_3\le1.170\cdot10^{-3}$, slope $\le0.05724$. The three cone data thresholds of v3.2 hold: slope $0.05724\le0.08$ ($\times1.40$), $\sup\abs{h_0''}=0.01883\le0.03$ ($\times1.59$), $\norm{h_0'''}_{L^2}=1.170\cdot10^{-3}\le0.004$ ($\times3.41$). The slope and $h''$ thresholds are design choices that trade against the closure margins (\fn{h6/fable/DESIGN.md} \S6).

\emph{Far} \tC\ (\fn{t9_far.py --v32} $\to$ \fn{t9_far_v32.json}; the same data as \fn{t9_far.json}, against the v3.2 thresholds). Exact normal-graph jets (\fn{curvjets.py}, review B5) from the slab-180 hulls ($\abs\kt\le1.77\cdot10^{-4}$, $\abs{\kt'}\le1.71\cdot10^{-8}$, $\abs{\kt''}\le2.33\cdot10^{-10}$) and error parts from C4 at $y=147-X_0$; graph over the centre tangent with $\abs\psi\le K_0(XL+2)\cdot1.01+e_1$. In $\lambda_1$ units: $\abs{h_{xx}}\le1.797\cdot10^{-4}$, $\abs{h_{xxx}}\le2.595\cdot10^{-8}$, $\abs{h_{xxxx}}\le6.964\cdot10^{-9}$, slope $\le0.04153$, $\norm{h_{xxx}}_{L^2(\abs x\le2L)}\le5.516\cdot10^{-7}$, against $3.53\cdot10^{-4}$, $1.17\cdot10^{-5}$, $1.03\cdot10^{-7}$, $0.08$, $3.32\cdot10^{-6}$. Each threshold is asserted to be $\le$ the exact value from $M_{2F}/113$, $M/113^2$, $M/113^3$, $N_{3F}/113^{1.5}$. In cylinder units the two far data thresholds give $0.02031\le0.04$ ($\times1.96$) and $6.626\cdot10^{-4}\le0.004$ ($\times6.03$) (\fn{t9_h6_v32.json}).

\paragraph{The v3.1 values, for the record.}\label{ss:LOv31} Against the v3.1 thresholds (far $M=1$: $\abs{h_{xx}}\le8.8\cdot10^{-3}$, $\abs{h_{xxx}}\le7.0\cdot10^{-5}$, $\abs{h_{xxxx}}\le6.2\cdot10^{-7}$, slope $\le0.21$, $\norm{h_{xxx}}_{L^2}\le113^{-3/2}$; cone $N_3\le0.14$), the same data give $\times48$, $\times2698$, $\times89$, $\times5.0$, $\times1509$ (far) and $\times119$ (cone), and far H4 $\le0.00397\,M^2$ at $M=1$. These are superseded by v3.2.

\subsection{Lemma O H4 per cylinder (review A1)}\label{ss:H4}
\begin{lemma}[\tC; \fn{p2/t8se/h4loc.py} $\to$ \fn{h4loc_A1_w180.json}; at $M=0.15$: \fn{constants/h4loc_v32.json}]\label{lem:H4}
For every cylinder of both families, at the good time $\tau^*$ of slab 180,
\[
\int_{-X}^X(h_0^{(5)})^2dX\le\bigl(\sqrt{\mathrm{err}}+\sqrt{\mathrm{base}}\bigr)^2,\qquad\mathrm{err}=L^7\cdot c_{\rm ch}^2\cdot\frac{D_2(\tau^*)}{28y_0^2},\qquad\mathrm{base}=L^7\abs S\sup_S\abs{h^{(5)}_{\rm base}}^2,
\]
where $S=[x_{lo},x_{hi}]$ is the cylinder's support, $y_0=x_{lo}-X_0-1\ge22$ (so $T_2=28Y_\ell^2\ge28y_0^2$ on $S$, and $S$ lies in the core of every level), and $h^{(5)}_{\rm base}$ is the exact fifth derivative of the base graph over the centre tangent (formula in $k^{(\le3)}$ and the angle, sympy-checked). Here $c_{\rm ch}$ is the computed chart factor, $c_{\rm ch}\le1.0103$ (cone, $r_c\in[370,380]$) and $1.0067$ (far). With $D_2(\tau^*)\le9.40\cdot10^{-12}$: worst cone ($r_c\in[370,380]$) $\le0.2393\,M^2$ at $M=0.12$; far (centre 376) $\le0.1761\,M^2$ at $M=0.15$ (at $M=1$: $0.00397\,M^2$); base parts $\le6.8\cdot10^{-6}$. Margins $\times4.18$ (cone) and $\times5.68$ (far).
\end{lemma}
\begin{proof}
In cylinder units ($X=x/L$, $h_0=h/L$), $\int(h_0^{(5)})^2dX=L^7\int_S(h^{(5)})^2dx$. The error part of $h^{(5)}$ is the fifth derivative of the error; $\int_Sw_4^2\le D_S/(28y_0^2)$ with $D_S:=\int_ST_2w_4^2\le D_2$, and $D_2(\tau^*)$ is bounded by Corollary~C5 on the one-slab window. A localized energy has the same leading term, so the window length, not the localization, decides; the global $D_2$ is used. The chart factor $c_{\rm ch}$ passes from the normal graph over the base to the graph over the centre tangent (below). Minkowski combines the two parts.
\end{proof}
\noindent\emph{The chart factor (cleanup O6b; closed).} The typed factor $1.001^2$ of the first draft is replaced by a computed one (\fn{h4loc.py}, with \fn{t9chart.py}/KCH). The true curve's $h^{(5)}$ over the cylinder's own centre tangent is expanded in \fn{arb} Taylor mode at the same base point; its $e^{(5)}$ coefficient is $A_5\le1.0103$, and the lower-order terms $e,\dots,e''''$ are charged explicitly ($\le5.4\cdot10^{-12}$ pointwise in $h^{(5)}$). The weight is $d\xi=M_t\,dx$, and the angle includes the true-versus-base centre tangent. The hand-off frame ($\lambda_f\le1$, $\abs{x_c}$) and the arc's excess over $XL$ fit the supports' slack $1$ (asserted, $\le0.43$). The chart factor is $c_{\rm ch}=A_5\sqrt{M_t}\le1.0103$ (cone $r_c\in[370,380]$) and $1.0067$ (far). Effect: worst cone $0.2297\to0.2393\,M^2$ ($\times4.35\to\times4.18$); far at $M=0.15$ $0.1712\to0.1761\,M^2$ ($\times5.84\to\times5.68$); far at $M=1$ $0.00386\to0.00397\,M^2$.

\subsection{Lemma O H1: the graph property, arc by arc}\label{ss:LOH1}
Lemma~O H1 asks: after translating to the centre and rotating the tangent to the horizontal, $\gamma\cap([-X,X]\times[-H,H])$ is the graph of $h_0$ over $[-X,X]$ with $\sup\abs{h_0}\le H-C_0\sqrt{t_1}$. Step~3 reads it arc by arc (v3.1 l.97, l.105, l.111; amended (H5)): only the arc through the centre must be a graph; other arcs passing through the box are irrelevant, CDF being local.

\emph{Cone family} \tC\ (\fn{t9_lemmaO.py --v32} $\to$ \fn{t9_lemmaO_v32.json}; T6-engine positions on 48 $s$-balls per element over the whole slab, correlated in $t_c$; position error from Lemma~4b and C4). In each of the 18 cylinders the graph arc is monotone in $\xi$ and leaves $[-XL,XL]$ at both ends; $\sup\abs{h_0}\le0.1872$ and $H=\sup\abs{h_0}+C_0\sqrt{t_1}\le0.1895$ with $C_0\sqrt{t_1}=0.0023$ (v3.1, $C_0\sqrt{t_1}=0.0021$: $H\le0.1893$).

\emph{Far family} \tC\ (closed, cleanup O1, Nov~12; \fn{p2/t9_lemmaO_far.py} $\to$ \fn{p2/out/t9_lemmaO_far.json}, stage 28 of the final tail, \fn{ok = true}). Data of v3.2: $X=2$, $L=113$, $(M,K)=(0.15,0.3)$. The script is a pass/fail check arc by arc: for every centre on the arm (and its mirror) from the first far centre to the far end, on base arclength $\pm G=\pm227.125$ around the centre the true tangent turns by at most $K_0J_{\max}G\le0.04095$ rad, with $K_0\le h_2=1.797\cdot10^{-4}$ (\fn{t9_far_v32.json}) and $J=\abs{\gamma'}\in[0.99667,1.00333]$ from $k_0$, $\abs e$ and $\abs{e'}$. So $\xi$ is increasing and $\abs{\xi(\pm G)}\ge226.18>XL=226$: the arc through the centre is a graph over $[-X,X]$ in its own tangent frame. Its height is $\sup\abs{h_0}\le\frac{X^2}2LK_0\sec^3\le0.04072$; with $C_0\sqrt{t_1}\le9.212\cdot10^{-4}$ (\fn{lemmaO_v32_consts.json}, far\_v32), $H=0.0417$ and $\sup\abs{h_0}\le H-C_0\sqrt{t_1}$, with slack $6.7\cdot10^{-5}$. The arc stays where \fn{t9_far.py}'s bounds hold ($x\ge148$) for every centre at base arclength $\ge375.125$. This covers ``centres $\ge376$'' in hand-off units read as arclength ($\ge375.998$ in base units) or as radius (points with $x<375.125$ have $r\le373.013$); the frame scale $\lambda_f\le1$ makes the base-unit computation conservative. The old script's pairwise separation test (H8's content, which failed at the self-crossing) is dropped (backup \fn{p2/t9_lemmaO_far_pre_cleanup.py}); H8 is not required. The far rows of \fn{t9_lemmaO.py} remain informational, since each covers a range of centres, so the frame moves within a row.

\subsection{H7, H8, H9 and the loss of embeddedness}\label{ss:embed}
\emph{H7 and H9} are internal to step~3: the time ratio holds by step~3's certificates (cone $1.6001\cdot10^{-7}$ against $t^*\ge3.683\cdot10^{-7}$; far at $(0.15,0.3)$ $6.134\cdot10^{-9}$ against $9.511\cdot10^{-8}$; v3.1 had $6.13\cdot10^{-9}$ against $1.23\cdot10^{-8}$), and regularity holds by the parabolic flow and Lemma~O's conclusion. Step~4 adds nothing to them.

\emph{The curve is not embedded at the hand-off, and need not be.} $S_1$ is not embedded: its loop has a transversal self-crossing, and a $C^0$-small normal perturbation of two transversally crossing strands still crosses, so every curve in the step-3 trapping set self-intersects (\fn{STEP4_DESIGN.md} \S2). An embedded initial curve must therefore be carried through the crossing by the flow itself, which is what steps P1 and P2 do. $\Gamma_0$ is embedded (Proposition~\ref{prop:G0}). In the float runs the curve first self-intersects at clip time $t\approx0.0028$, i.e.\ $t\approx0.003$, about $0.0018$ after $\Gamma_0$ (\fn{T3_REPORT.md} \S1) \tF, inside the first P1 slabs. Neither P1 (Fourier variables) nor P2 (a normal graph over the base) is sensitive to it. At the hand-off the self-intersection is certified (\fn{p2/t9_axis_crossing.py} $\to$ \fn{p2/out/t9_axis_crossing.json}, stage 29 of the final tail): on each of 16 $t_c$-balls covering slab 180, $\operatorname{Im}$ of the true arm changes sign on element 93, $s\in[29/32,30/32]$, at arclength $x\in[3.269\cdot10^4,3.398\cdot10^4]$ from the tip ($\operatorname{Re}\in[2.235\cdot10^4,2.316\cdot10^4]$), with $\abs e\le0.525$ and a certified $\abs{\operatorname{Im}}\ge\abs e+7.9$ at both ends; so the arm crosses the axis and, with its mirror image, the curve self-intersects there (informational). Step~3 applies Lemma~O arc by arc and needs only H1--H7 and H9 (v3.1 l.97, l.105, l.111): \emph{H8 is not a requirement}, and the crossing affects no hypothesis.

\section{Lemma O H6 and the step-3 amendment v3.2}\label{s:H6}
\subsection{H6 under v3.1 (for the record)}\label{ss:H6v31}
Lemma~O H6 (\fn{lemmaO_v3.tex}): $K_2=K_3=1.5M$ are assumed on $[0,t_1]$ in every cylinder of the family and must be recovered at every centre simultaneously, which needs Lemma~O at all centres in $c+[-X,X]$. With the v3.1 families this fails (\fn{p2/t9_h6.py} $\to$ \fn{t9_h6.json}, kept to document the failure):
\begin{itemize}
\item cone family by overlapping centres: at a cylinder's inner end $L_c/L_p=4$, so the recovered bound is $4\cdot(0.0189+0.0395)=0.234$ against $K_2=0.18$ ($\times0.77$); it closes only for $L_c/L_p\le3.08$;
\item cone cylinders beyond the last cone centre ($r>380$) take their bound from the far family: $(95/113)\cdot0.374=0.315$ against $0.18$ ($\times0.57$);
\item both far checks pass: far by overlap $0.374$ against $1.5$ ($\times4.0$); far inner strip from the cone family $0.176$ against $1.5$ ($\times8.5$).
\end{itemize}
So under v3.1 alone, (H5) is not met. The design review of the failure (\fn{p2/h6/fable/DESIGN.md} \S1--2) locates the cause in the far family: with $(M,K)=(1,1.5)$ its deviation $C_2\sqrt t=0.35$ is too large for the cone cylinders that reach $r>475$ (in its closure, cone $K_2$ $\times0.54$, $K_3$ $\times0.44$). It also shows that the cone family need not change: on the strip $r\in[50,150)$ the (B) windows present at $\tau_1$ are fixed in physical space and keep their hand-off energies (step~3's own D5 input), which removes the ratio-4 problem. (An earlier plan, \fn{p2/h6/H6_PLAN.md}, raised the cone $K$ to $0.25$; it is superseded.)

\subsection{The amendment v3.2}\label{ss:v32}
\emph{Status: approved by the authors and built:} \fn{writeup/step3_proof_v32/step3_proof_v32.tex} (step~3 v3.2; v3.1 is kept intact in \fn{writeup/step3_proof/}), \fn{writeup/lemmaO_v32/lemmaO_v32.tex} (Lemma~O v3.2, with Theorem~H6 and \fn{lemmaO_v32_h6.py}), and \fn{writeup/step3_amend_v32/amend_v32.tex} (summary of the amendment). \emph{Lemma~R is proved; the step-3 gap in Lemma~S(B) is the subject of Hypothesis~(F)} (\S\ref{ss:v32status}). Design: \fn{step4/p2/h6/fable/DESIGN.md} \S4, \S7. In units of $\lambda_1$, with $r$ the distance along the $S_1$ ray and $t\in[t_1,T)$:
\begin{itemize}
\item \emph{Families.} Cone unchanged: centres $r_m\in[200,380]$, $L_m=r_m/4$, $X=3$, cylinder $[r_m/4,7r_m/4]$, conclusion region $[3r_m/4,5r_m/4]$, $M=0.12$, $K_2=K_3=0.18$. Far: centres $\ge376$, unit 113, $X=2$, conclusion region $\pm113$ (covering $r\ge263$), \emph{$M=0.15$, $K_2=K_3=0.3$}. Strip: $r\in[50,150)$, the (B) windows present at $\tau_1$.
\item \emph{Covering of a cone cylinder $c$.} $r_p\in[r_c/4,150)$: windows; $r_p\in[150,285]$: the cone member $r_m=4r_p/3$; $r_p\in(285,475]$: the member $r_m=380$; $r_p\in(475,7r_c/4]$: the far member centred at $p$. A far cylinder is covered by cone members for $r_p\in[150,263)$ and by far members beyond.
\item \emph{Scaling and frames.} The $C_j$ do not depend on $t$; a member of unit $L_m$ has the deviations at $50^{-4}(50/L_m)^2$ times the horizon; $\abs{h_{xx}}$ scales by $L_c/L_m$ and $\norm{h_{xxx}}_{L^2}$ by $(L_c/L_m)^{3/2}$; frame factors $(1+s^2)^{3/2}$, $(1+s^2)^2$ with the slope bound $s$.
\item \emph{Order.} Lemma~O in every member gives $\abs{g_x}\le C_1\sqrt t$, $\abs{g_{xx}}\le C_2\sqrt t$, $\int_{-1}^1g_{xxx}^2\le C_3t$; then the slope, $K_2$ (worst at $r_p=150$, $r_c=380$) and $K_3$ (sum over the pieces) are recovered strictly below $\frac14$, $K_2$, $K_3$; the families are compact in the centre, so the continuity argument closes. In step~3 this is D5a (Lemma~O conclusions from B4 and (H5)), then D4 (the windows, using D5a at their outer edge), then D5b (Theorem~H6 of \fn{lemmaO_v32.tex}, which improves B4).
\item \emph{The strip, frozen profile (the ``Risk-1'' repair).} On the strip the (B) window deviation is now measured against the profile frozen at entry, $U_e(r)=1.3056\,\alpha(\tau_e)(r/\lambda(\tau_e))^p$, not against $U(\alpha(\tau))$, whose forcing at fixed physical $r$ is not controlled (\fn{step3_proof_v32.tex} \S3.7; F-U in \fn{DESIGN.md} \S6). The rerun (\fn{lemmaS_v32/lemmaS_Bbox.py}, \fn{FROZEN=1}) leaves the boxes, $\hat G$ and $\gamma'$ unchanged.
\item \emph{(H5) amended.} Lemma~O H1--H5, H7, H9 arc by arc with the $M$ above, plus six data thresholds (own cylinder units, whole cylinder): slope $s_0=0.08$, $\sup\abs{h_0''}\le m_2$ with $m_2=0.03$ (cone), $0.04$ (far), and $\norm{h_0'''}_{L^2}\le n_3=0.004$ (both); far thresholds in $\lambda_1$ units $\abs{h_{xx}}\le3.53\cdot10^{-4}$, $\abs{h_{xxx}}\le1.17\cdot10^{-5}$, $\abs{h_{xxxx}}\le1.03\cdot10^{-7}$, $\norm{h_{xxx}}_{L^2(\abs x\le2L)}\le3.32\cdot10^{-6}$, slope $\le0.08$, $\int(h_0^{(5)})^2\le M^2=0.0225$. The old $L^2$ thresholds ($0.14$ cone, $1.0$ far) are dropped. The thresholds are needed: with the data bounded only by $M=0.12$ the closure fails ($0.12+1.9\cdot0.0395>0.18$).
\end{itemize}

\begin{table}[htbp]\centering\small
\caption{The H6 closure of v3.2 as computed in step~4 (\fn{p2/t9_h6_v32.py} $\to$ \fn{t9_h6_v32.json}; \fn{arb}, outward rounding, no sampling, cone centres in cells of width 1). It belongs to the step-3 side of v3.2 and uses the step-4 data only through the thresholds; its rows depend only on design values. Its strip terms (the worst cone slope, the strip part $0.04969$ of cone $K_2$, the strip piece of cone $K_3$) come from the strip window energies at levels 0--2, whose level-2 part used Lemma~R (now proved) and which belong to the strip input of Hypothesis~(F): \emph{the closure rows are conditional on (F)} (\S\ref{ss:v32status}). The $t^*$ rows are not.}\label{tab:v32}
\begin{tabular}{@{}L{0.52\linewidth}L{0.14\linewidth}L{0.10\linewidth}L{0.12\linewidth}@{}}\toprule
item & recovered bound & against & margin\\\midrule
cone slope (strip is worst; covered points $0.08753$) & $0.13532$ & $0.25$ & $\times1.85$\\
far slope & $0.08753$ & $0.25$ & $\times2.86$\\
cone $K_2$ ($r_c=380$, $r_p=150$; strip $0.04969$) & $0.10664$ & $0.18$ & $\times1.69$\\
cone $K_3$ & $0.09348$ & $0.18$ & $\times1.93$\\
far $K_2$ & $0.13125$ & $0.3$ & $\times2.29$\\
far $K_3$ & $0.12549$ & $0.3$ & $\times2.39$\\
Lemma~O $t^*$, cone (unchanged, grid 14) & $3.683\cdot10^{-7}$ & $1.6001\cdot10^{-7}$ & $\times2.30$\\
Lemma~O $t^*$, far at $(0.15,0.3)$ & $9.511\cdot10^{-8}$ & $6.134\cdot10^{-9}$ & $\times15.5$\\
time-equivalent margin of the closure (cone $K_2$) & & & $\times3.86$ in time\\\bottomrule
\end{tabular}\end{table}
\noindent The thinnest closure item is cone $K_2$, $\times1.69$; the binding time margin stays the cone $t^*$, $\times2.30$. Every v3.1 output is unchanged; the v3.2 outputs sit beside them. The step-4 checks of v3.2 are the stages ``v3.2: \dots'' of \fn{FINAL_manifest.json}, run on the final data. (\fn{p2/out/v32_manifest.json} records the first v3.2 run, on earlier inputs; its hashes no longer match the files, which the final tail regenerated.)

\subsection{The step-3 dependency: Lemma R proved, Hypothesis (F)}\label{ss:v32status}
The step-3 edits of v3.2 (\fn{DESIGN.md} \S7) are in place in the built texts: Lemma~O's H5, H6 and far case in \fn{lemmaO_v32.tex}; in \fn{step3_proof_v32.tex}, (H5), (W4) (the far a-priori bound $1.5$ becomes $0.3$), the frozen profile in \S3.7, \S3.8, B4, and D5 split into D5a and D5b.

\emph{Lemma R is proved.} The (B) energies of the strip windows need a cutoff at the outer end of the window region; its commutator adds to the level-2 window energy a term that needs $z_3$, i.e.\ $h_{xxxx}$ in $L^2$, time-averaged, on the ramp (\fn{lemmaS_v32/ramp_out.py}). Lemma~R, a level-4 companion of Lemma~O (the time average over $[t_1,T]$ of $\int_{-1}^1h_{xxxx}^2$ on the conclusion regions of the cone members $r_m\in[253,380]$, used with the ramp moved to $[190,380]$), is now proved, with arb-certified constants (\fn{writeup/lemmaO_v32/lemmaR.pdf}). It closes the outer ramp.

\emph{A gap in step~3 itself.} An independent referee confirmed a gap in step~3, present already in v3.1 (\fn{writeup/step3_gap_v32/gap_v32.pdf}, \fn{REFEREE_check.md}). Lemma~S(B)'s far-field window-energy inequalities drop the window-edge (total-derivative) terms. Step~3's trapping argument uses the resulting far-field bound $\abs f\le\frac12$ (W4). This is a missing estimate with large numerical room (derived $\abs f\le0.093$ against $\frac12$), not a failed inequality.

\emph{Consequence.} The main theorem (\S\ref{s:main}) is stated \textbf{conditional on Hypothesis~(F)} (\fn{writeup/main/main.pdf}). (This is not the item (F) of the bootstrap hypotheses of Theorem~P2, \S\ref{s:P2}.) A fix, a uniformly local Lemma~S(B), is being evaluated (go/no-go tests); until it is in, (F) stays a hypothesis.

\emph{The dependency, precisely.} Strip statement of Lemma~S(B) for $t\in[t_1,T)$ (part of (F)) $\Rightarrow$ Theorem~H6 of \fn{lemmaO_v32.tex} (D5b, which improves B4) $\Rightarrow$ Theorem~W of step~3 v3.2. \emph{Step~4's own results do not depend on (F)}: Theorem~4's data are certified at $\tau_1$ and do not involve the strip energies for $t>t_1$, and neither do the six thresholds or the $t^*$ rows. Step~4's v3.2 H6 closure margins (Table~\ref{tab:v32}: the slope, $K_2$ and $K_3$ rows) use the strip input, which is part of (F), so those rows are conditional on (F). The frozen-profile repair is done: its rerun leaves the boxes unchanged (\S\ref{ss:v32}).

\emph{Other step-3 caveats} (\fn{step3_proof_v32.tex} \S6, final-rigor-pass items, not step-4 items): the Lemma~S assembly arithmetic runs in IEEE double; the class weights of the $L^4$ terms use sympy floats; the sliver where $\zeta<1$; literature references; the frozen-profile runs of \fn{lemmaS_F4B.py} and of the $\alpha_1=1.25\cdot10^{-4}$ point were not redone.

\section{The main theorem}\label{s:main}
We combine the four steps. Step~1 constructs $S_1$ (\fn{s1_tight/s1_tight.tex} and the step-1 supplement, with the geometric checks (a)--(d): double point, embedding, mirror, tail). Step~2 proves linear stability of $S_1$ modulo symmetries: every eigenvalue with $\operatorname{Re}\lambda>-\frac18$ is $1$ or $\frac14$ in the even class, and $\frac14$ or $0$ in the odd class (\fn{step2_proof.tex}). Step~3 proves trapping with a floor (Theorem~W, \fn{step3_proof_v32/step3_proof_v32.tex}: v3.1 as amended by v3.2), conditional on (H1)--(H6) at $t_1$ and on Hypothesis~(F) (\S\ref{ss:v32status}); Lemma~R, used in D5b, is proved. Step~4 is Theorem~4.

\begin{named}{Main Theorem (conditional on Hypothesis~(F))}
Let $S_1$ be a shrinker in the certified set of step~1, and let $\gamma$ be the CDF evolution of the curve $\Gamma_0$ of Theorem~4. Assume Hypothesis~(F) of step~3 (\S\ref{ss:v32status}; \fn{writeup/main/main.pdf}). Then:
\begin{enumerate}[label=(\roman*)]
\item $\Gamma_0$ is embedded, and $\gamma$ has a finite maximal existence time $T$ with $0<T-t_1\le\lambda_1^4/(1-4\bar a)$, where $\bar a:=\sup_{\tau\ge\tau_1}\abs a\le2.19103\cdot10^{-7}$, $t_1\in[4.5978574775,4.5978586144]$ and $\lambda_1\in[6.49\cdot10^{-5},6.75\cdot10^{-5}]$; so $T-t_1\le2.1\cdot10^{-17}$.
\item The singular set at time $T$ is a single point $p_*$, and the curvature is unbounded as $t\uparrow T$.
\item The rate is self-similar: $\gamma(t)=p(t)+\lambda(t)\Gamma(\tau)$ with $\lambda(t)=(T-t)^{1/4}(1+O(\bar a))$ and $p(t)\to p_*$.
\item The rescaled profile stays near $S_1$ for all $\tau\ge\tau_1$: $\abs\alpha\le1.5\cdot10^{-4}$ and $\sqrt\Lam\le5.31\cdot10^{-8}$; after the tail start the barrier decays at rate $\varrho=2.69\cdot10^{-3}$ in $\Lam$ down to the floor $\sqrt{\Lam_F}=4.94\cdot10^{-10}$, and eventually $\abs\alpha\le3.43\cdot10^{-7}$. Convergence to $S_1$ is not claimed.
\item Away from $p_*$ the curve converges smoothly as $t\uparrow T$.
\item Embeddedness is lost during step~4: $\Gamma_0$ is embedded and $\gamma(t_1)$ is not.
\end{enumerate}
\end{named}
\begin{proof}
By Theorem~4 the curve $\Gamma_1$ satisfies (H1)--(H6) of step~3 v3.2 at $t_1$ (Lemma~O H1 for the far family is closed, \S\ref{ss:LOH1}). Theorem~W of step~3 v3.2, whose Lemma~O H6 (D5b) uses Lemma~R (proved) and the strip statement of Lemma~S(B), then gives, under Hypothesis~(F), (W1)--(W5), which are (ii)--(v) and the bound in (i). Embeddedness of $\Gamma_0$ is Proposition~\ref{prop:G0}, and the self-intersection at $t_1$ is \S\ref{ss:embed}. $T$ is the maximal time by the step-3 continuation criterion (P1) \tA. $\lambda_1$ is the base scale on slab 180 ($\ln\tilde\lambda\in[-9.64236,-9.60485]$, Lemma~\ref{lem:t1}) times the step-3 frame factor $\lambda\in[0.99999490,0.99999511]$ and the modulation factor $\lambda/\tilde\lambda\in e^{\pm3.38\cdot10^{-8}}$ (Remark~\ref{rem:t1}). The tip curvature is $k(0)/\lambda+\dots$ with $k(0)\in[1.53982063787,1.53982063788]>0$ (\S\ref{ss:S1}); this settles the step-3 caveat ``$k(0)\neq0$ used without proof''.
\end{proof}
\noindent\emph{Dependencies.} The theorem inherits the caveats of step~3 (\fn{step3_proof_v32.tex} \S6): Hypothesis~(F) (Lemma~R is proved), the literature facts (P1) (local well-posedness for immersed data, the continuation criterion and interior regularity), and the final-rigor-pass items listed there (\S\ref{ss:v32status}). So $T$ lies within $2.1\cdot10^{-17}$ above the $t_1$ interval (clip time $T+0.001$). In (iv), $\varrho$ is step~3's tail rate, not P2's $\varrho_c$. V3.2 changes only the far-cylinder constants of (W4); (W1), (W2), (W3) and (W5) are unchanged.

\section{Open items}\label{s:open}
\begin{enumerate}[label=\textbf{O\arabic*.},leftmargin=*]
\item \textbf{Lemma~O H1 for the far family: closed} (cleanup O1, Nov~12; \S\ref{ss:LOH1}). \fn{p2/t9_lemmaO_far.py} is a pass/fail check arc by arc (stage 28 of the final tail, \fn{ok = true}): tangent turning $\le0.04095$ rad within $\pm227.125$ of each centre, $\abs{\xi(\pm G)}\ge226.18>226$, $\sup\abs{h_0}\le0.04072$, $C_0\sqrt{t_1}\le9.212\cdot10^{-4}$, $H=0.0417$. It holds for every centre from arclength $375.125$ on. H8 is not required.
\item \textbf{Step-3 dependency: Lemma~R proved; Hypothesis~(F) open} (\S\ref{ss:v32status}). Lemma~R is proved, with arb-certified constants (\fn{writeup/lemmaO_v32/lemmaR.pdf}), and closes the outer ramp. An independent referee confirmed a gap in step~3 itself, present in v3.1 (\fn{writeup/step3_gap_v32/gap_v32.pdf}, \fn{REFEREE_check.md}): Lemma~S(B)'s far-field window-energy inequalities drop the window-edge (total-derivative) terms, and the trapping argument uses the resulting bound $\abs f\le\frac12$ (W4). It is a missing estimate with large numerical room (derived $\abs f\le0.093$ against $\frac12$), not a failed inequality. The main theorem is therefore conditional on Hypothesis~(F) (\fn{writeup/main/main.pdf}). Step~4's own results do not depend on (F); its v3.2 H6 closure margins use the strip input, which is part of (F). A fix, a uniformly local Lemma~S(B), is being evaluated (go/no-go tests). Under v3.1 alone, H6 is not met (\S\ref{ss:H6v31}), and the same gap is present there.
\item \textbf{The true curve's $k_{\max}$ at $t_1$} (Theorem~4(b); \fn{STEP4_RESULT.md} \S3). Only the base value $[2.279\cdot10^4,2.376\cdot10^4]$ is certified; the true curve's needs a pointwise tip bound on $e''$. No hypothesis of step~3 uses it.
\item \textbf{Kit B} (P2 and the hand-off on the user's machine): \textbf{still pending}. Until then the P2 and T9 chain is reproduced only by the two passes of the final tail on the sandbox (\S\ref{app:repro}).
\item \textbf{Lemma~O H7 and H9} are internal to step~3 and hold there (\S\ref{ss:embed}); \fn{STEP4_RESULT.md} \S3 lists them for completeness. No step-4 action.
\item \textbf{Script items found while writing: closed, except (c)} (cleanup of Nov~12; backups \fn{*_pre_cleanup.py}; pre-cleanup manifest \fn{FINAL_manifest_pre_cleanup.json}). (a) \emph{Closed:} \fn{t1clock.py} now includes $\lambda/\tilde\lambda$ and $x_{\rm NL}$ slab by slab; the quoted $t_1$ interval stands (Remark~\ref{rem:t1}). (b) \emph{Closed:} \fn{h4loc.py} uses a computed chart factor ($\le1.0103$ cone, $1.0067$ far); H4 is now $\le0.2393\,M^2$ ($\times4.18$) and $\le0.1761\,M^2$ ($\times5.68$) (\S\ref{ss:H4}). (d) \emph{Closed:} the new script \fn{p2/t9_axis_crossing.py} gives the crossing on 16 of 16 time balls, element 93, $\abs e\le0.525$ (replacing $1.17$ and $1.24$), and it and the far-H1 stage are stages 28 and 29 of \fn{FINAL_manifest.json}, which now has 31 stages (\S\ref{app:repro}). (c) \emph{Remaining, not covered by the cleanup:} the rotation-mode norm in \fn{h1.py} uses generous hand bounds, a contribution $\le3.1\cdot10^{-11}$ to $\sqrt\Lam$ (\S\ref{ss:alpha}); it changes no conclusion.
\end{enumerate}
\noindent Not proof items: the AI-use statement and the literature check (\S\ref{s:rel}, \S\ref{s:ai}) are for the authors.

\section{Related work}\label{s:rel}
CDF (surface diffusion of plane curves) is fourth order. It has no comparison principle, and it need not preserve embeddedness or convexity; see Giga and Ito~\cite{GI}. Near circles the flow is global and converges, e.g.\ \cite{Wh13}. Self-similar shrinkers, among them the figure eight, are constructed and classified in \cite{EGMWW}. Our $S_1$ is the non-embedded one-loop shrinker, with a computer-assisted enclosure from step~1. Wheeler~\cite{Wh26} treats pinch-off by surface diffusion for surfaces: a conical pinch-off, for a family of data. The present work is different in kind: a single explicit planar curve, a non-embedded singularity model, and a computer-assisted proof. The validated integration uses \fn{arb} \cite{Arb}, CAPD \cite{CAPD}, and the self-consistent bounds and first-exit method of Zgliczy\'nski \cite{Zg}. \AUTH{full literature review; verify all bibliographic details.}

\section{AI use statement}\label{s:ai}
\AUTH{write the AI use statement. Facts to cover: the computations, certificates and design notes of steps 1--4 were developed with AI agents (Claude, Anthropic) under the authors' direction; this document was assembled by Claude from the project files; an independent AI review (\fn{STEP4_REVIEW.md}) was run on step 4; the authors' own verification, the reproduction kits and the independent Windows runs of step 3 and of kit A; responsibility for the content.}

\appendix
\section{What is computed and how}\label{app:how}

\subsection{Ball arithmetic and outward rounding}\label{app:arb}
\begin{itemize}
\item \emph{Arithmetic.} Every certified quantity of step 4 is computed with python-flint's \fn{arb} (midpoint-radius ball arithmetic; 96 to 192 bits; $\Gamma_0$ at 192 bits, the T9 chart at 128 bits). Step~1's \fn{ext813} uses CAPD's $C^0$ Lohner method (order 25). Sympy is used offline only: for symbolic identities checked once (\fn{lfsym.py}, \fn{cforms.py}, \fn{n3sym.py}, \fn{n3qsym.py}, \fn{n3wcheck.py}, \fn{nl7.derive}, \fn{t9chart_check.py}, \fn{h4loc.py}'s $h^{(5)}$ formula), never for a bound at run time.
\item \emph{Floats.} Floats are used to \emph{find} objects: the P1 collocation and Lohner bases, the polish of the P2 base, the step sizes. Every certified statement is about the exact data that come out (dyadic centres, stored balls), and the residuals of those data are enclosed exactly. No certified bound comes from sampling.
\item \emph{Rounding.} When a ball is stored as doubles, its ends are moved outward with \fn{nextafter}; stored round-to-nearest hulls (T6) are widened by 2 ulps on read. JSON outputs are read back exactly. The review's rounding spots are fixed (B9): \fn{switch.py} weight arguments, \fn{ycomp.py}'s \fn{up()}, \fn{h1.py}'s cell tiling (exact dyadic cells), T6's stored scalars and ranges.
\item \emph{No typed inputs on the wrong side.} Inputs passed between stages are read from the JSON outputs of the earlier stage (B4), and the typed initial data of the switch ($Y_0=6.21\cdot10^{-13}$, $F^0$) are checked against the switch output in the manifest (\fn{switch_vs_typed}). Typed allowances remain only where stated: \fn{h1.py}'s rotation-mode hand bounds (\S\ref{ss:alpha}) and \fn{ycomp.py}'s conservative $\mathrm{DN}$ and $\mathrm{D4CUB}$ (\S\ref{ss:Y}).
\item \emph{Rounding of quoted numbers.} Upper bounds are rounded up and lower bounds and limits down. Where \fn{STEP4_RESULT.md} rounds to nearest, this document quotes the certified file value rounded outward: the P1 $k_{\max}$ enclosure (\fn{t3/p1/status.json}), $c_0^{\rm tip}$, the frame $\lambda$ and the shell constants (\fn{t9_cone_kch.json}).
\end{itemize}

\subsection{Proof of the Galerkin-limit lemma}\label{app:gal}
\begin{proof}[Proof of Lemma~\ref{lem:gal}]
\emph{1. Galerkin systems.} For $N\ge M$, $u^N$ solves $u^{N\prime}_n=-\lambda_nu^N_n+N_n(u^N)$ for $n\le N$, with $u^N_m\equiv0$ for $m>N$ and $u^N(0)=\Pi_Nu(0)$. This is a polynomial ODE, so a local solution exists and continues while bounded. The difference $z^N=u^N-\bar u-P\rho$ satisfies $z^{N\prime}_n=-\lambda_nz^N_n+F^N_n$ with $F^N_n=(Jz^N)_n+Q_n-r_n-(D\rho)_n$.
\emph{2. First exit.} $\abs{z_n^N(0)}\le a_n<Y_n$. Let $s^*=\sup\{s\le\Delta:z^N([0,s])\subset W\}>0$. On $[0,s^*]$ we have $\abs{F_n^N}\le F_n$, and Duhamel gives $\abs{z^N_n(s^*)}\le g_n(s^*)$. Since $g_n'=e^{-\lambda_nt}(F_n-\lambda_na_n)$ has constant sign, $g_n(s^*)\le\max(g_n(0),g_n(\Delta))<Y_n$. Finitely many coordinates move, so continuity gives $s^*=\Delta$.
\emph{3. Compactness.} Uniformly in $N$, $\abs{u^N_n}\le B_n:=\sup_I\abs{\bar u_n}+\sup_I\abs{(P\rho)_n}+Y_n$, and $B_n=Tn^{-6}$ for large $n$. The convolution sums in $N_n$ involve $\sum m^kB_m$ with $k\le3$, which converge since $6>4$. So $\abs{\partial_su^N_n}\le\lambda_nB_n+\tilde N_n$. Arzel\`a--Ascoli and a diagonal argument give a subsequence converging for every $n$, uniformly on $I$.
\emph{4. Limit.} In the mild form $u_n^N(s)=e^{-\lambda_ns}u_n(0)+\int_0^se^{-\lambda_n(s-\sigma)}N_n(u^N)d\sigma$, dominated convergence passes to the limit, and the closed bounds persist.
\emph{5. Regularity.} For $n$ large, $\abs{u_n}\le Tn^{-6}$ and $\abs{N_n}\le Cn^{-4}$. So $u\in C^0H^5\cap C^1L^2$, and $u$ solves the PDE in $L^2$.
\emph{6. Uniqueness.} For two solutions, $\delta=u-u^\infty$ satisfies $\frac12\frac d{ds}\norm\delta^2=-\norm{\delta_{\alpha\alpha}}^2+\norm{\delta_\alpha}^2+\ip{N(u)-N(u^\infty)}\delta$. Each term of the difference is bounded by $C_A(\norm{\delta_\alpha}^2+\norm{\delta_\alpha}\norm\delta+\norm{\delta_{\alpha\alpha}}\norm\delta)$, with $C_A$ depending on $\norm{v}_\infty,\norm w_\infty$ of both solutions ($H^5\subset C^3$). Interpolation and Gr\"onwall give $\delta\equiv0$.
\end{proof}
\noindent Steps 1--5 also give existence from $H^5$ data, so chaining sub-intervals and slabs \emph{constructs} the solution on all of P1.

\subsection{The manifest}\label{app:repro}
\emph{One command} (from \fn{step4/p2}): \fn{python3 -I final_tail.py}. It writes \fn{p2/out/FINAL_manifest.json}: every stage in order, with its command, working directory, environment, and the SHA-256 of its inputs and outputs, plus the SHA-256 of the scripts and of 117 upstream files.
\begin{itemize}
\item \emph{One version of every input.} T6 $=$ \fn{p2/t6c} (residuals and jets of the corrected base); T7 $=$ \fn{t7/constants.json} (republished Nov~12); T8 $=$ \fn{constants_b1.json} and \fn{n3_b1.json}, rebuilt in the tail by \fn{cconst} and \fn{merge} (T7 slabs $\ge2$ only); base windows $=$ \fn{h1slab/win_summary.json}. No final stage reads \fn{p2/t6/*}, \fn{step4/t6/residuals.json}, \fn{se_corr.json}, \fn{constants_fin.json} or \fn{n3.json} (stated in the manifest and enforced by the \fn{T6DIR} and \fn{T8CONSTS} environment of each stage). The manifest checks that \fn{t6c/residuals.json} is the version \fn{constants_b1} records.
\item \emph{The 31 stages}: \fn{cconst} (179 slabs), \fn{merge}, \fn{switch}, \fn{ycomp}, \fn{clam}; the T8 iteration (\fn{camp}, \fn{n3far}, \fn{camp --n3}, \fn{n3far}) and the window \fn{camp}, \fn{n3boxes}, \fn{t9} in two passes A and B; then \fn{t9_far}, \fn{h4loc}, Lemma~O H1/H8 (\fn{t9_lemmaO.py}), H6 under v3.1 (\fn{t9_h6.py}), the five v3.2 stages (\fn{t9_far.py --v32}, \fn{h4loc} with \fn{FAR_M=0.15}, \fn{t9_v32_checks.py}, \fn{t9_h6_v32.py}, \fn{t9_lemmaO.py --v32}), the far-family Lemma~O H1 (\fn{t9_lemmaO_far.py}, stage 28), the axis crossing (\fn{t9_axis_crossing.py}, stage 29) and \fn{t1clock}. Pass B's \fn{t9} reproduces pass A's byte for byte (\fn{t9_fixed_point}).
\item \emph{Upstream} (hashes in the manifest; not rerun by the tail): the B8 containment check, the 16 Krawczyk points, the 32 derivative balls, the 32 window balls and their summaries of the H1 stage (\fn{t9h}). They read only the base polynomials, which equal the corrected base on $x<0.6\lt$.
\item \emph{Second pass.} With \fn{cconst}'s slab files moved aside, so that every stage recomputes, all 199 non-log outputs reproduce byte for byte, except the CPU and wall-time fields of the 179 \fn{cconst} slab files, \fn{constants_b1.json} and \fn{switch_s2.json}; these three are identical once those fields are removed (checked). The pass-1 manifest is \fn{p2/out/FINAL_manifest_pass1.json}.
\item \emph{Cleanup rerun (O6a--c).} \fn{final_tail.py --from 'h4loc on the final amp'} was run twice; stages 0--19 were reused (442 paths verified). Pass 2 reproduces all 11 non-log outputs of stages 20--30 byte for byte; the only differences are the stages' seconds and \fn{basegeo}'s build-time lines in three logs (\fn{checks.cleanup_second_pass}).
\end{itemize}

\begin{table}[htbp]\centering\scriptsize
\caption{SHA-256 (first 16 hex digits) of the certified outputs used here, as recorded in \fn{FINAL_manifest.json} and checked against the files on Nov~12. Paths relative to \fn{step4/}; those starting with \fn{s1_tight/} are under \fn{/home/claude/}.}\label{tab:hashes}
\setlength{\tabcolsep}{1.5pt}
\begin{tabular}{@{}ll|ll@{}}\toprule
\fn{p2/out/FINAL_manifest.json} & \fn{cd37f42214d97561} & \fn{p2/out/h1slab/summary.json} & \fn{5813699c98e56dcf}\\
\fn{p2/final_tail.py} & \fn{534e4ce4e2bd7fc0} & \fn{p2/out/h1slab/win_summary.json} & \fn{f3a48403bd4aaed5}\\
\fn{t2/gamma0_final.json} & \fn{e3a0763ef5e56209} & \fn{p2/out/t9_h1_clam_s180.json} & \fn{0a759b5b786d38bf}\\
\fn{t3/p1/ckpt/slab_00096.json} & \fn{9a53a8b6be713c3a} & \fn{p2/out/t9_cone_kch.json} & \fn{6312bb98696838dc}\\
\fn{t1/polished_v2/manifest.json} & \fn{bffa87863a836020} & \fn{p2/out/t9_far.json} & \fn{20ae0224e32a4365}\\
\fn{p2/t6c/residuals.json} & \fn{64e7fe04ec80770a} & \fn{p2/out/h4loc_A1_w180.json} & \fn{2031df5f1a749b12}\\
\fn{t7/constants.json} & \fn{ddb6097f3f3d0bee} & \fn{p2/out/t9_lemmaO.json} & \fn{553ba87fca5a32c3}\\
\fn{p2/out/switch_s2.json} & \fn{e4d6f89f6a778d96} & \fn{p2/out/t9_h6.json} (v3.1) & \fn{3774ce9d6fe5846c}\\
\fn{p2/out/ycomp_final.json} & \fn{e33e8e414dfaaa3f} & \fn{p2/out/t9_far_v32.json} & \fn{e0c9c533517ddb3a}\\
\fn{p2/t8se/constants/constants_b1.json} & \fn{b99fff442520c83f} & \fn{p2/t8se/constants/h4loc_v32.json} & \fn{77c52bec8666359a}\\
\fn{p2/t8se/constants/n3_b1.json} & \fn{c9deb0fc0e512856} & \fn{p2/out/t9_v32_checks.json} & \fn{7193fbe5543d4c1b}\\
\fn{p2/t8se/constants/amp_A1_w180.json} & \fn{b7f3191940906596} & \fn{p2/out/t9_h6_v32.json} & \fn{0841b98780ef812e}\\
\fn{s1_tight/ext_tip/prof_tip.txt} & \fn{1e07195d48e879a9} & \fn{p2/out/t9_lemmaO_v32.json} & \fn{07ea3d750ed3c379}\\
\fn{s1_tight/ext813/ext813_prof.txt} & \fn{49e01a0933e65aaf} & \fn{p2/out/t1clock.json} & \fn{2ede3d87d9199a56}\\
\fn{s1_tight/consistency/check_outA5.json} & \fn{43a4ae8b56ed562f} & \fn{p2/out/t9_lemmaO_far.json} & \fn{8915a8a49a3f1b3a}\\
 & & \fn{p2/out/t9_axis_crossing.json} & \fn{ce63ff1087e9347e}\\\bottomrule\end{tabular}\end{table}

\subsection{Reproduction kits}\label{app:kits}
\paragraph{Step-3 kit.} \fn{s1spec/kit_step3.zip} (v2, v3.1 outward-rounded), SHA-256 \fn{1d1818251248b523be684653bf2958cbce0ff3cbe90c26e3cd8252a5bc740cf3}. Parts: verifier, IBP modulation, Lemma~S, Lemma~O. The independent run on Windows~11 (\fn{kit_step3/user_runs/windows11_v2_nov10_check.txt}) reports \textbf{91 PASS, 0 FAIL, 0 WARN}, with the verifier bit-identical to the reference.

\paragraph{Kit A ($\Gamma_0$ and P1).} \fn{s1spec/kit_step4a.zip}, current SHA-256 \fn{0a496d4969a713d4798760e5231440556854cabafcb9dda43b9b3f2c0ce12a38}. It reproduces the $\Gamma_0$ certificate (\fn{t2/gamma0.py final}, about 3~min) and the P1 chain through all 96 slabs (\fn{t3/run_p1.py}), on Python~3.13 with pinned python-flint, numpy and scipy, single-threaded BLAS and \fn{PYTHONHASHSEED=0}; it is resumable and touches nothing outside its folder. The authors' run on Windows~11 (Nov~11; Python 3.13.7; \fn{s1spec/kit_step4a/user_runs/}):
\begin{itemize}
\item \emph{$\Gamma_0$: bit-identical} to the reference ($b_n$ doubles, the 192-bit $b_2$ centre and half-width, the arc nodes), with the same Krawczyk contraction $3.01\cdot10^{-12}$ and embeddedness data.
\item \emph{P1: 96 of 96 slabs validated} (no retries), in 6.18~h of CPU. The final slab-96 enclosures of $t$, $L$, $X_0$ and $k_{\max}$ overlap the reference; for example $k_{\max}\in[29.352874545326006,29.352874545327026]$ against the reference $[29.352874545325992,29.352874545327012]$.
\item \emph{Intermediate slabs.} The kit's acceptance rule also asks for plain overlap of the slab-end enclosures of $t$, $L$, $X_0$ at every slab. On Windows that test fails at most intermediate slabs, and the checker's summary line reads ``TOTAL FAIL (12 PASS, 3 FAIL, 1 WARN)''. The reason is that platform floats change the step sizes: the chains choose float step sizes that differ in the last digits, so their slab-end times differ (by up to $8.7\cdot10^{-16}$), and enclosures taken at different times need not overlap; from slab 82 on the discrete choices ($M$, block, $h$) differ too. Transported to a common time, $\abs{q_u-q_r}\le\operatorname{rad}_u+\operatorname{rad}_r+S_q(\abs{t_u-t_r}+\operatorname{rad}t_u+\operatorname{rad}t_r)$ with $S_q$ twice the reference's slope (\fn{user_runs/transport_check.py}, a sanity check outside the proof), all 97 slab-end enclosures of $L$ and $X_0$ agree (0 failures). The plain-overlap rule was the wrong criterion for intermediate slabs; the transported comparison should replace it in the kit's checker.
\end{itemize}

\paragraph{Kit B (P2 and the hand-off).} \textbf{Pending.} It will cover the base store, T6 on the corrected base, T7, the final tail and the H1 stage.

\subsection{Commands}\label{app:cmd}
{\scriptsize\begin{verbatim}
cd /home/claude/s1spec/step4/p2
python3 -I final_tail.py     # the whole final tail on the B1 data -> out/FINAL_manifest.json
# upstream H1 stage (hashes in the manifest's upstream list)
cd t9h; export COHERENT=1 S1PROF=/home/claude/s1_tight/ext_tip/prof_tip.txt \
    S1TT=10.249077651120778,10.249077651121361
./run_h1slab.sh 0 7 & ./run_h1slab.sh 8 15   # 16 x h1.py p<t_k> (Krawczyk + Lambda), 32 x h1dt.py <tc> 1/32
python3 -I win_slab.py > ../out/h1slab/win_cmds.txt; ./run_winslab.sh 0 15 & ./run_winslab.sh 16 31
# B8 containment
python3 -I /home/claude/s1_tight/consistency/check_outA5.py
# offline check of the chart algebra (sympy)
cd ..; python3 -I t9chart_check.py
\end{verbatim}}
\noindent Upstream of the tail: \fn{t2}: \fn{gamma0.py final}; \fn{t3}: \fn{run_p1.py --target-k 25 --Mmax 3600}, then \fn{--target-k 28}; \fn{t1}: \fn{seed_p1.py 95 80}, \fn{run3.py}, \fn{polish_v2.py}, \fn{verify_v2.py}, \fn{se_v2.py}; \fn{p2/t6c}: \fn{slabc.py}; \fn{t7}: \fn{run_all.py} (\fn{T7_REPORT.md}, with \fn{--nlclosure}); \fn{s1_tight}: \fn{run_tight.sh}, \fn{ext813/run_ext813.sh}, \fn{ext_tip/run_tip.sh}.

\subsection{Notes on the sources}\label{app:notes}
\fn{STEP4_RESULT.md} is authoritative. Where its text and its tables differ, the tables agree with the certified outputs and are used here: its status paragraph and \S5 still name Lemma~O H6 and the step-3 amendment as open, while this document follows the later step-3 status (Lemma~R proved, main theorem conditional on Hypothesis~(F)); its \S3 paragraph on B5 quotes $E_3$ limits $1.517$--$1.518\cdot10^{-11}$, computed before B1, while \S1, \S2 and the final \fn{camp} give $1.575\cdot10^{-11}$. Other sources that predate the final tail: \fn{LEMMA_F.md} \S6.1 (a one-slab requirement of $1.79\cdot10^{-11}$, before B5), \fn{LEMMA_F.md} \S6.4 and \fn{T8_REPORT.md} \S1--5 (earlier data), the bracketed step-4 values in the threshold table of \fn{amend_v32.tex} (e.g.\ far slope $0.0428$ against the final $0.04153$), and \fn{v32_manifest.json}. The closure table of \fn{lemmaO_v32.tex} (step~3's own computation, \fn{lemmaO_v32_h6.py}, with the frozen-profile strip) gives cone $K_3\le0.0938$ ($\times1.92$, strip piece $0.0194$) against step~4's $0.09348$ ($\times1.93$), and quotes H4 as $0.304\,M^2$ (cone) and $0.229\,M^2$ (far), the design estimates, against step~4's certified $0.2393\,M^2$ and $0.1761\,M^2$; the other closure rows agree. Step~3's numbers govern the step-3 side; step~4 quotes its own outputs.

\begin{thebibliography}{9}\small
\bibitem{Arb} F.~Johansson, Arb: efficient arbitrary-precision midpoint-radius interval arithmetic, \emph{IEEE Trans.\ Comput.}\ 66 (2017).
\bibitem{CAPD} T.~Kapela, M.~Mrozek, D.~Wilczak, P.~Zgliczy\'nski, CAPD::DynSys: a flexible C++ toolbox for rigorous numerical analysis of dynamical systems, \emph{Commun.\ Nonlinear Sci.\ Numer.\ Simul.}\ 101 (2021).
\bibitem{EGMWW} M.~Edwards, A.~Gerhardt-Bourke, J.~McCoy, G.~Wheeler, V.-M.~Wheeler, The shrinking figure eight and other solitons for the curve diffusion flow, \emph{J.\ Elasticity} 119 (2015); arXiv:1505.02978.
\bibitem{GI} Y.~Giga, K.~Ito, Loss of convexity of simple closed curves moved by surface diffusion, in: \emph{Topics in Nonlinear Analysis}, Progr.\ Nonlinear Differential Equations Appl.\ 35, Birkh\"auser (1999).
\bibitem{Wh13} G.~Wheeler, On the curve diffusion flow of closed plane curves, \emph{Ann.\ Mat.\ Pura Appl.}\ 192 (2013); arXiv:1201.3735.
\bibitem{Wh26} G.~Wheeler, Pinchoff by surface diffusion, arXiv:2608.21882 (2026).
\bibitem{Zg} P.~Zgliczy\'nski, Rigorous numerics for dissipative partial differential equations II. Periodic orbit for the Kuramoto--Sivashinsky PDE --- a computer-assisted proof, \emph{Found.\ Comput.\ Math.}\ 4 (2004).
\end{thebibliography}
\end{document}
